% =====================================================================
%  Algorithmic Game Theory — Appunti (UniCal, a.a. 2026/2027)
%  Compilazione:  powershell -File compila.ps1   ->  main.pdf
%
%  Struttura: una cartella per lezione in lezioni/ (NN-slug/lezione.tex),
%  inclusa qui in ordine. Ogni lezione si compila anche da sola:
%  vedi lezione-standalone.tex e compila.ps1.
%  Le slide del docente stanno in ../lessons/, i notebook in ../labs/.
% =====================================================================
\documentclass[11pt, a4paper, oneside, openany]{book}
% =====================================================================
%  preambolo.tex — pacchetti, stile e ambienti comuni a tutto il documento
% =====================================================================

% ---------- Lingua e codifica ----------
\usepackage[T1]{fontenc}
\usepackage[utf8]{inputenc}
\usepackage[italian]{babel}

% ---------- Font ----------
% Times-like (newtx): tratto pieno e scuro, larghezza compatta.
% Anche la matematica è coerente col testo.
\usepackage{newtxtext}
\usepackage{newtxmath}
\usepackage[scaled=0.88]{beramono}   % monospace leggibile per il codice
\usepackage{microtype}

% ---------- Pagina ----------
\usepackage[a4paper, top=2.5cm, bottom=2.5cm, left=2.5cm, right=2.5cm, headheight=14pt]{geometry}
\setlength{\parindent}{0pt}
\setlength{\parskip}{0.55em}
\linespread{1.08}

% ---------- Colori ----------
\usepackage[dvipsnames]{xcolor}
\definecolor{unical}{HTML}{B00000}     % rosso delle slide
\definecolor{blu}{HTML}{0B3A91}        % blu delle slide di sezione
\definecolor{verde}{HTML}{1E7B34}
\definecolor{arancio}{HTML}{C0620A}
\definecolor{grigio}{HTML}{4A4A4A}
\definecolor{sfondo}{HTML}{F5F6F8}

% ---------- Matematica ----------
\usepackage{amsmath}
\usepackage{mathtools}

% ---------- Tabelle, liste, figure ----------
\usepackage{booktabs}
\usepackage{array}
\usepackage{multirow}
\usepackage{tabularx}
\usepackage{enumitem}
\setlist{itemsep=0.2em, topsep=0.3em}
\setlist[itemize,1]{label=\textcolor{unical}{\small$\blacktriangleright$}}
\usepackage{graphicx}
\usepackage{float}
\usepackage{caption}
\usepackage{needspace}   % \needspace{n\baselineskip}: va a capo pagina se manca spazio
\captionsetup{font=small, labelfont={bf,color=blu}}

\usepackage{tikz}
\usetikzlibrary{arrows.meta, positioning, trees, calc, shapes.geometric}

% ---------- Titoli ----------
\usepackage{titlesec}
% Capitolo = una lezione
\titleformat{\chapter}[display]
  {\normalfont\sffamily\bfseries\color{blu}}
  {\large\MakeUppercase{\chaptertitlename}\ \thechapter}
  {0.3em}
  {\Huge}
  [\vspace{0.3em}{\color{unical}\titlerule[1.2pt]}]
\titlespacing*{\chapter}{0pt}{0pt}{1.5em}
\titleformat{\section}
  {\normalfont\sffamily\Large\bfseries\color{blu}}{\thesection}{0.8em}{}
\titleformat{\subsection}
  {\normalfont\sffamily\large\bfseries\color{black}}{\thesubsection}{0.8em}{}
\titleformat{\subsubsection}
  {\normalfont\sffamily\normalsize\bfseries\color{grigio}}{\thesubsubsection}{0.8em}{}
\addto\captionsitalian{\renewcommand{\chaptername}{Lezione}}

% ---------- Intestazioni ----------
\usepackage{fancyhdr}
\pagestyle{fancy}
\fancyhf{}
\fancyhead[L]{\small\sffamily\color{grigio}Algorithmic Game Theory}
\fancyhead[R]{\small\sffamily\color{grigio}\leftmark}
\fancyfoot[C]{\small\sffamily\thepage}
\renewcommand{\headrulewidth}{0.4pt}
\renewcommand{\chaptermark}[1]{\markboth{Lezione \thechapter\ -- #1}{}}
\fancypagestyle{plain}{\fancyhf{}\fancyfoot[C]{\small\sffamily\thepage}\renewcommand{\headrulewidth}{0pt}}

% ---------- Box (definizioni, esempi, ...) ----------
\usepackage[most]{tcolorbox}
\tcbset{
  appunti/.style={
    enhanced, breakable, sharp corners,
    boxrule=0pt, leftrule=3pt,
    left=8pt, right=8pt, top=5pt, bottom=5pt,
    before skip=0.8em, after skip=0.8em,
    fonttitle=\sffamily\bfseries, coltitle=black,
    attach title to upper, after title={.\ },
  }
}
\newtcolorbox[auto counter, number within=chapter]{definizione}[1][]{
  appunti, colback=blu!5, colframe=blu,
  title={\textcolor{blu}{Definizione~\thetcbcounter}\ifx\relax#1\relax\else\ (#1)\fi}}
\newtcolorbox[use counter from=definizione]{teorema}[1][]{
  appunti, colback=unical!5, colframe=unical,
  title={\textcolor{unical}{Teorema~\thetcbcounter}\ifx\relax#1\relax\else\ (#1)\fi}}
\newtcolorbox{esempio}[1][]{
  appunti, colback=verde!5, colframe=verde,
  title={\textcolor{verde}{Esempio}\ifx\relax#1\relax\else\ (#1)\fi}}
\newtcolorbox{osservazione}[1][]{
  appunti, colback=sfondo, colframe=grigio,
  title={\textcolor{grigio}{Osservazione}\ifx\relax#1\relax\else\ (#1)\fi}}
% Punti delicati, errori tipici, imprecisioni nelle slide
\newtcolorbox{attenzione}[1][]{
  appunti, colback=arancio!7, colframe=arancio,
  title={\textcolor{arancio}{Attenzione}\ifx\relax#1\relax\else\ (#1)\fi}}
% Contenuti aggiunti negli appunti, non presenti nelle slide
\newtcolorbox{extra}[1][]{
  appunti, colback=white, colframe=grigio!60, leftrule=1.5pt,
  borderline west={0.4pt}{4pt}{grigio!60},
  title={\textcolor{grigio}{Approfondimento}\ifx\relax#1\relax\else\ (#1)\fi}}
% Riepilogo di fine lezione
\newtcolorbox{riepilogo}{
  enhanced, breakable, sharp corners, colback=blu!4, colframe=blu, boxrule=0.8pt,
  fonttitle=\sffamily\bfseries, title={Riepilogo della lezione}}

% ---------- Teoria vs Laboratorio ----------
% \parteteoria e \partelab aprono le due metà di una lezione con un'etichetta colorata.
\newcommand{\etichetta}[2]{%
  \tikz[baseline=(t.base)]\node[fill=#1, text=white, font=\sffamily\bfseries\small,
    inner xsep=6pt, inner ysep=3pt](t){#2};}
\newcommand{\parteteoria}{%
  \par\medskip\noindent\etichetta{blu}{TEORIA}\par\smallskip}
\newcommand{\partelab}[1]{%
  \clearpage
  \section{Laboratorio: #1}
  \noindent\etichetta{unical}{LAB}\par\smallskip}
\newcommand{\fonte}[1]{\hfill{\small\sffamily\color{grigio}[#1]}}

% ---------- Codice ----------
\usepackage{listings}
\lstdefinestyle{python}{
  language=Python,
  basicstyle=\ttfamily\small,
  keywordstyle=\color{blu}\bfseries,
  commentstyle=\color{verde}\itshape,
  stringstyle=\color{unical},
  numberstyle=\tiny\color{grigio}, numbers=left, numbersep=8pt,
  backgroundcolor=\color{sfondo},
  frame=leftline, framerule=2pt, rulecolor=\color{blu},
  xleftmargin=1.8em, framexleftmargin=1.8em,
  showstringspaces=false, breaklines=true, tabsize=4,
  columns=fullflexible, keepspaces=true,
  morekeywords={True,False,None},
  literate={à}{{\`a}}1 {è}{{\`e}}1 {é}{{\'e}}1 {ì}{{\`i}}1 {ò}{{\`o}}1 {ù}{{\`u}}1
           {−}{{-}}1 {∞}{{$\infty$}}1 {×}{{$\times$}}1,
}
\lstdefinestyle{pseudo}{
  style=python, numbers=none, mathescape=true,
  morekeywords={def,if,else,for,each,return,break,in,and,or,not},
}
\lstset{style=python}

% ---------- Tris (tic-tac-toe) ----------
% \tris{X,O,e, e,X,e, O,e,e}: 9 celle riga per riga, e = vuota.
\newcommand{\trisX}{\textbf{X}}
\newcommand{\trisO}{\textcolor{unical}{\textbf{O}}}
\newcommand{\trise}{}
\newcommand{\tris}[2][0.42]{%
  \begin{tikzpicture}[scale=#1, baseline=(current bounding box.center)]
    \draw[thick] (1,0)--(1,3) (2,0)--(2,3) (0,1)--(3,1) (0,2)--(3,2);
    \foreach \c [count=\k from 0] in {#2} {
      \pgfmathtruncatemacro\riga{\k/3}
      \pgfmathtruncatemacro\colonna{mod(\k,3)}
      \node[font=\sffamily] at (\colonna+0.5, 2.5-\riga) {\csname tris\c\endcsname};
    }
  \end{tikzpicture}}

% ---------- Utilità ----------
\newcommand{\en}[1]{\textit{#1}}            % termine inglese originale
\newcommand{\kw}[1]{\textbf{#1}}            % parola chiave
\newcommand{\R}{\mathbb{R}}
\newcommand{\N}{\mathbb{N}}
\newcommand{\Ev}{\mathbb{E}}
% Relazioni di preferenza
\newcommand{\pref}{\succsim}                % x pref y: x e' almeno buono quanto y
\newcommand{\npref}{\not\succsim}
\newcommand{\spref}{\succ}                  % preferenza stretta
\newcommand{\indif}{\sim}                   % indifferenza
\newcommand{\lott}{\mathcal{L}}             % insieme delle lotterie semplici
% Teoria dei giochi
\newcommand{\BR}{\mathrm{BR}}               % best response
\newcommand{\uno}{\mathbf{1}}               % vettore di tutti 1
\DeclareMathOperator{\supp}{supp}           % supporto di una strategia mista
\DeclareMathOperator{\core}{Core}
% Tabella dei payoff di un gioco a due giocatori: celle "a, b"
\newcommand{\pay}[2]{$#1,\,#2$}
\DeclareMathOperator*{\argmax}{arg\,max}
\DeclareMathOperator*{\argmin}{arg\,min}

% ---------- Link (per ultimo) ----------
\usepackage[hidelinks]{hyperref}
\hypersetup{pdftitle={Algorithmic Game Theory -- Appunti}, pdflang={it}}
\usepackage{bookmark}


\begin{document}
\frontmatter

% ---------- Frontespizio ----------
\hypersetup{pageanchor=false}
\begin{titlepage}
  \centering
  \vspace*{3cm}
  {\sffamily\bfseries\color{unical}\large UNIVERSITÀ DELLA CALABRIA\par}
  \vspace{0.3cm}
  {\sffamily\color{grigio}Anno accademico 2026/2027\par}
  \vspace{2.5cm}
  {\color{blu}\rule{\linewidth}{1.2pt}\par}
  \vspace{0.6cm}
  {\sffamily\bfseries\Huge\color{blu} Algorithmic Game Theory\par}
  \vspace{0.4cm}
  {\sffamily\Large Appunti del corso: teoria e laboratori\par}
  \vspace{0.6cm}
  {\color{blu}\rule{\linewidth}{1.2pt}\par}
  \vspace{1.5cm}
  {\large Corso tenuto da Sebastiano A.\ Piccolo\par}
  \vfill
  {\small\color{grigio} Ultima compilazione: \today\par}
\end{titlepage}
\hypersetup{pageanchor=true}

\tableofcontents

\mainmatter
% ---------- Lezioni ----------
% =====================================================================
%  Lezione 1 — Introduzione all'AGT, giochi combinatori, Minimax
%  Materiale: Lesson_1_Introduction_to_AGT.pdf, Lab_Minimax_blank.ipynb
% =====================================================================
\chapter{Introduzione all'Algorithmic Game Theory}

\parteteoria

% ---------------------------------------------------------------------
\section{Informazioni sul corso}

\begin{itemize}
  \item \kw{Lezioni}: mercoledì 10:30--13:30 (teoria), giovedì 11:30--13:30 (laboratorio).
  \item \kw{Teams}: \emph{Algorithmic Game Theory 2026/27} (il codice per entrare è nelle slide).
  \item \kw{Esame}: un \kw{progetto} in cui si deve \emph{formalizzare} un problema e
        \emph{implementarne} una soluzione usando gli strumenti dell'AGT.
\end{itemize}

\begin{osservazione}[conseguenza pratica]
Dato che l'esame è un progetto, contano due cose: saper \emph{modellare} un problema
(chi sono i giocatori, quali azioni hanno, quali utilità) e saper \emph{scrivere codice}
che calcola il concetto di soluzione giusto. I laboratori vanno quindi studiati quanto la teoria.
\end{osservazione}

\subsection*{Programma (provvisorio)}
\begin{enumerate}[label=\arabic*.]
  \item Giochi combinatori (\en{Combinatorial Games})
  \item Teoria dell'utilità (\en{Utility theory})
  \item Teoria dei giochi non cooperativi (\en{Non-cooperative game theory})
  \item Teoria dei giochi coalizionali (\en{Coalitional game theory})
  \item Scelta sociale computazionale (\en{Computational Social Choice})
  \item Mercati di abbinamento (\en{Matching Markets})
  \item Progettazione di meccanismi (\en{Mechanism Design})
  \item Divisione equa (\en{Fair division})
\end{enumerate}

\subsection*{Testi di riferimento}
Non esiste un unico libro che copra tutto il corso; queste sono le risorse gratuite indicate:
\begin{itemize}
  \item Y.~Shoham, K.~Leyton-Brown, \emph{Multiagent Systems};
  \item N.~Nisan, T.~Roughgarden, É.~Tardos, V.~Vazirani, \emph{Algorithmic Game Theory};
  \item A.~Karlin, Y.~Peres, \emph{Game Theory, Alive};
  \item F.~Brandt et al., \emph{Handbook of Computational Social Choice}.
\end{itemize}

% ---------------------------------------------------------------------
\section{Che cos'è la teoria dei giochi}

\begin{definizione}[Teoria dei giochi]
La \kw{teoria dei giochi} studia le \emph{decisioni prese da più soggetti}
(\en{multiperson decision making}), con lo scopo di prevedere come più agenti
\kw{razionali} interagiranno e si comporteranno in situazioni di \kw{conflitto}
e/o di \kw{cooperazione}.
\end{definizione}

I quattro ingredienti della definizione:
\begin{enumerate}
  \item \kw{Conflitto}: le azioni di un giocatore influenzano gli altri.
  \item \kw{Cooperazione}: i giocatori possono unire le forze per raggiungere un obiettivo comune.
  \item \kw{Comportamento razionale}: ogni giocatore vuole \emph{massimizzare} la propria
        \kw{utilità attesa} (\en{expected utility}).
  \item \kw{Predizione}: vogliamo sapere cosa succederà. Lo si fa tramite i
        \kw{concetti di soluzione} (\en{solution concepts}), come l'equilibrio di Nash o il valore di Shapley.
\end{enumerate}

\section{Che cos'è l'Algorithmic Game Theory}

\begin{definizione}[Algorithmic Game Theory]
L'\kw{Algorithmic Game Theory} (AGT) è l'area all'intersezione tra \emph{teoria dei giochi}
e \emph{progettazione di algoritmi}, il cui obiettivo è progettare algoritmi in
\emph{ambienti strategici}. \fonte{Nisan et al.\ 2007}
\end{definizione}

In concreto, l'AGT si occupa di progettare algoritmi per:
\begin{itemize}
  \item \kw{calcolare} in modo efficiente i concetti di soluzione (es.\ trovare un equilibrio);
  \item \kw{aggregare} voti e preferenze;
  \item \kw{allocare} risorse scarse.
\end{itemize}

% ---------------------------------------------------------------------
\section{Panoramica delle aree dell'AGT}

\subsection{Giochi non cooperativi}

I giochi non cooperativi modellano situazioni di \kw{conflitto}: i giocatori competono
tra loro e non possono necessariamente cooperare.
Esempi: scacchi, sasso-carta-forbici, dilemma del prigioniero.

L'obiettivo è modellare una situazione e capire cosa succede. Il concetto di soluzione
tipico è l'\kw{equilibrio}.

\subsubsection*{Il dilemma del prigioniero}

Due giocatori scelgono simultaneamente se \emph{cooperare} (C) o \emph{tradire} (D, \en{defect}).
I numeri sono le utilità (anni di carcere con segno negativo): in ogni cella il primo valore è di P1,
il secondo di P2.

\begin{table}[H]
\centering
\renewcommand{\arraystretch}{1.4}
\begin{tabular}{cc|c|c|}
  & \multicolumn{1}{c}{} & \multicolumn{2}{c}{\textbf{P2}} \\
  & \multicolumn{1}{c}{} & \multicolumn{1}{c}{Cooperate} & \multicolumn{1}{c}{Defect} \\ \cline{3-4}
  \multirow{2}{*}{\textbf{P1}} & Cooperate & $-1,\,-1$ & $-5,\,0$ \\ \cline{3-4}
                               & Defect    & $0,\,-5$  & $-3,\,-3$ \\ \cline{3-4}
\end{tabular}
\caption{Dilemma del prigioniero in \emph{forma normale} (gioco bi-matrice).}
\end{table}

Questa rappresentazione a tabella si chiama \kw{forma normale} (o gioco \emph{bi-matrice}):
una matrice di utilità per ciascuno dei due giocatori.

\begin{definizione}[Ottimo di Pareto]
Un esito è \kw{Pareto-ottimo} se non esiste un altro esito in cui \emph{qualcuno} sta meglio
e \emph{nessuno} sta peggio. In altre parole: non si può migliorare la situazione di un
giocatore senza peggiorare quella di qualcun altro.
\end{definizione}

\begin{definizione}[Equilibrio di Nash, idea]
Un profilo di strategie è un \kw{equilibrio di Nash} se nessun giocatore può migliorare
la propria utilità cambiando \emph{da solo} la propria scelta.
\end{definizione}

\begin{esempio}[analisi del dilemma]
\textbf{Nash.} Per P1, tradire conviene \emph{qualunque} cosa faccia P2:
se P2 coopera, $0 > -1$; se P2 tradisce, $-3 > -5$. Lo stesso vale per P2.
Quindi (D,\,D) è l'unico equilibrio di Nash, con utilità $(-3,-3)$.
Da (C,\,C), invece, ciascuno ha interesse a deviare (passa da $-1$ a $0$).

\textbf{Pareto.} (C,\,C) con $(-1,-1)$ è Pareto-ottimo. L'equilibrio (D,\,D) \emph{non} lo è:
(C,\,C) fa stare meglio entrambi.
\end{esempio}

Il dilemma mostra la tensione centrale della teoria dei giochi:
\begin{center}
\textbf{Ottimalità di Pareto $\Rightarrow$ efficienza}\qquad
\textbf{Equilibrio di Nash $\Rightarrow$ stabilità}
\end{center}
L'esito efficiente non è stabile e l'esito stabile non è efficiente.

\begin{attenzione}[definizione nelle slide]
Le slide definiscono l'ottimo di Pareto come \emph{``nessun giocatore può deviare senza peggiorare
qualcun altro''}. È una formulazione informale: l'ottimo di Pareto \emph{non} riguarda le deviazioni
unilaterali (quelle sono il criterio di Nash), ma il confronto tra \emph{esiti}. Nel dilemma sono
Pareto-ottimi anche (C,\,D) e (D,\,C): passando da $(-5,0)$ a qualunque altro esito P2 starebbe
peggio. L'unico esito \emph{non} Pareto-ottimo è (D,\,D).
\end{attenzione}

\subsubsection*{Applicazioni}
\begin{itemize}
  \item algoritmi per giocare: scacchi, poker, \ldots;
  \item \en{Generative Adversarial Networks} (GAN): generatore e discriminatore sono due giocatori;
  \item sicurezza: pattugliamento, sorveglianza, \ldots;
  \item decisioni sequenziali robuste (\en{robust sequential decision making}).
\end{itemize}

\subsection{Giochi coalizionali}

Nei giochi coalizionali (o cooperativi) i giocatori \kw{uniscono le forze} per un obiettivo
comune, ad esempio mettono insieme i soldi per comprare un bene.
Le domande sono due:
\begin{itemize}
  \item quale coalizione è più probabile che si formi?
  \item come si divide il bene (il \en{payoff}) tra i membri?
\end{itemize}
Concetto di soluzione principale: il \kw{valore di Shapley}.

\begin{esempio}[tre imprese]
Tre imprese $A, B, C$ possiedono ciascuna una parte della tecnologia per sviluppare un prodotto.
Due imprese insieme guadagnano 100; una da sola guadagna 0; tutte e tre insieme guadagnano
comunque 100. La \kw{funzione caratteristica} $v$ assegna a ogni coalizione il suo valore:
\begin{align*}
  v(\{A\}) &= v(\{B\}) = v(\{C\}) = 0,\\
  v(\{A,B\}) &= v(\{A,C\}) = v(\{B,C\}) = 100,\\
  v(\{A,B,C\}) &= 100.
\end{align*}
\begin{itemize}
  \item Se si forma la coalizione di tutti e tre e si divide in parti uguali: $33{,}33$ a testa.
  \item Ma $B$ e $C$ da soli ottengono 100 e, dividendo, $50$ a testa: stanno meglio
        abbandonando $A$. La divisione in parti uguali della grande coalizione non è stabile.
\end{itemize}
\end{esempio}

\begin{esempio}[variante: un giocatore più importante]
Ora $v(\{A,B\}) = v(\{A,C\}) = 100$ ma $v(\{B,C\}) = 0$ (il resto invariato).
Senza $A$ non si produce nulla: $A$ è più ``importante'' e la divisione in parti uguali
non è più equa. Come dividere in modo equo? Con il \kw{valore di Shapley}.
\end{esempio}

\begin{extra}[calcolo del valore di Shapley nella variante]
Il valore di Shapley (sarà definito formalmente più avanti) dà a ogni giocatore il suo
\emph{contributo marginale medio} su tutti gli ordini in cui i giocatori possono entrare
nella coalizione. Con 3 giocatori gli ordini sono $3! = 6$:

\begin{center}
\renewcommand{\arraystretch}{1.2}
\begin{tabular}{lccc}
\toprule
Ordine & contributo di $A$ & contributo di $B$ & contributo di $C$ \\
\midrule
$A,B,C$ & 0   & 100 & 0   \\
$A,C,B$ & 0   & 0   & 100 \\
$B,A,C$ & 100 & 0   & 0   \\
$B,C,A$ & 100 & 0   & 0   \\
$C,A,B$ & 100 & 0   & 0   \\
$C,B,A$ & 100 & 0   & 0   \\
\midrule
Media   & $66{,}67$ & $16{,}67$ & $16{,}67$ \\
\bottomrule
\end{tabular}
\end{center}
Esempio di lettura: nell'ordine $B,C,A$, $A$ entra quando ci sono già $B$ e $C$, e porta il valore
da $v(\{B,C\}) = 0$ a $v(\{A,B,C\}) = 100$: contributo 100.
Il totale è $100$, cioè tutto il valore della grande coalizione. Nel caso simmetrico
di prima si ottiene invece $33{,}33$ a testa.
\end{extra}

\subsubsection*{Applicazioni}
\begin{itemize}
  \item divisione equa (\en{fair division});
  \item classifiche di importanza, potere di voto, \ldots;
  \item spiegabilità dell'IA: \kw{SHAP} (valori di Shapley applicati alle feature di un modello);
  \item \en{network science}: massimizzazione dell'influenza, centralità, \ldots
\end{itemize}

\subsection{Scelta sociale computazionale}

La \kw{scelta sociale computazionale} (\en{Computational Social Choice}) si occupa del calcolo
dei risultati delle regole di voto e dei problemi legati a rappresentare e raccogliere
le preferenze dei giocatori. Domande tipiche:
\begin{itemize}
  \item come si decide il vincitore di un'elezione?
  \item a quali condizioni la democrazia è possibile?
\end{itemize}

\begin{esempio}[paradosso di Condorcet]
Tre votanti ordinano tre candidati:
\begin{center}
\renewcommand{\arraystretch}{1.2}
\begin{tabular}{lccc}
\toprule
Votante & 1\textsuperscript{a} preferenza & 2\textsuperscript{a} preferenza & 3\textsuperscript{a} preferenza \\
\midrule
Votante 1 & A & B & C \\
Votante 2 & B & C & A \\
Votante 3 & C & A & B \\
\bottomrule
\end{tabular}
\end{center}
Chi dovrebbe vincere? Confrontiamo i candidati a coppie, a maggioranza:
\begin{itemize}
  \item $A$ batte $B$ (votanti 1 e 3);
  \item $B$ batte $C$ (votanti 1 e 2);
  \item $C$ batte $A$ (votanti 2 e 3).
\end{itemize}
La preferenza collettiva è \kw{ciclica}: $A \succ B \succ C \succ A$.
Anche se ogni votante ha preferenze coerenti (transitive), la maggioranza non le ha,
e nessun candidato batte tutti gli altri.
\end{esempio}

\subsubsection*{Applicazioni}
Elezioni, decisioni collettive, e nel machine learning i \en{random forest} e gli
\en{ensemble} di modelli (ogni modello ``vota'' una predizione).

\subsection{Mechanism design}

Il \kw{mechanism design} è una sorta di teoria dei giochi \emph{al contrario}: invece di
prendere un gioco come dato e analizzarne gli esiti, si vuole \emph{progettare} un gioco
che incentivi i giocatori a comportarsi nel modo desiderato.

\begin{esempio}[aste]
In un'asta si vuole che le persone dichiarino il \emph{vero} prezzo a cui comprerebbero
un oggetto o venderebbero un servizio.
\end{esempio}

Casi di successo: eBay, Google Ads, appalti pubblici e aste.

\subsection{Matching markets}

Nei \kw{mercati di abbinamento} si vogliono abbinare agenti ad altri agenti, o agenti a
oggetti, beni, servizi, garantendo qualche proprietà sociale desiderabile.
Esempi: studenti e alloggi universitari, pazienti e donatori di rene, lavoratori e compiti.
Nel corso si studierà un algoritmo premiato con il Nobel per lo \kw{stable matching}.

% ---------------------------------------------------------------------
\section{Giochi combinatori}

\begin{definizione}[Gioco combinatorio]
Un \kw{gioco combinatorio} ha le seguenti proprietà:
\begin{enumerate}
  \item \kw{due giocatori}, di solito chiamati Giocatore~1 e Giocatore~2, oppure MAX e MIN;
  \item \kw{informazione perfetta}: entrambi conoscono l'intero stato del gioco;
  \item \kw{nessun elemento casuale}: niente dadi, mazzi mescolati o eventi aleatori;
  \item \kw{finito}: il gioco termina sempre (niente partite infinite);
  \item \kw{turni alternati}: i giocatori muovono uno alla volta;
  \item \kw{somma zero}: il guadagno di un giocatore è esattamente la perdita dell'altro.
\end{enumerate}
\end{definizione}
Esempi: scacchi, dama, go, tris (\en{tic-tac-toe}).

\subsection{Dalla teoria dei giochi al giocare}

\kw{Problema centrale}: dato lo stato attuale del gioco, qual è la mossa \emph{migliore}?

La mossa migliore è quella che porta al miglior esito possibile \emph{assumendo che anche
l'avversario giochi in modo perfetto}. L'idea chiave è modellare l'avversario come
intelligente e capace di dare sempre la contromossa migliore.

\begin{osservazione}[perché assumere un avversario perfetto]
Se un giocatore vince contro la \kw{risposta migliore} (\en{best response}) dell'avversario,
vince anche contro qualunque risposta non ottimale. Ragionare sul caso peggiore dà quindi
una \emph{garanzia}: il risultato non può essere peggiore di quello calcolato.
\end{osservazione}

\subsection{L'albero di gioco}

\begin{definizione}[Albero di gioco]
Ogni gioco di questo tipo si può rappresentare come un \kw{albero di gioco}:
\begin{itemize}
  \item i \kw{nodi} sono stati del gioco (una configurazione della scacchiera);
  \item gli \kw{archi} sono mosse, che portano da uno stato a un altro;
  \item la \kw{radice} è lo stato attuale;
  \item le \kw{foglie} sono stati terminali (vittoria di qualcuno o pareggio).
\end{itemize}
\end{definizione}

\begin{figure}[H]
\centering
\begin{tikzpicture}[>=Stealth, x=1cm, y=1cm]
  \node (r)  at (0,0)    {\tris{e,e,e, X,O,e, e,e,e}};
  \node (a1) at (3.2,2)  {\tris{O,e,e, X,O,e, e,e,e}};
  \node (a2) at (3.2,0)  {\tris{e,e,e, X,O,e, e,e,O}};
  \node (a3) at (3.2,-2) {\tris{e,O,e, X,O,e, e,e,e}};
  \node      at (3.2,-1) {$\vdots$};
  \node (b1) at (6.4,0)  {\tris{X,e,e, X,O,e, e,e,O}};
  \node (c1) at (9.6,0)  {\tris{X,e,e, X,O,e, O,e,O}};
  \node (d1) at (12.8,1.5)  {\tris{X,e,O, X,O,e, O,X,O}};
  \node (d2) at (12.8,-1.5) {\tris{X,e,X, X,O,e, O,O,O}};
  \draw[->, blu] (r.east) -- (a1.west);
  \draw[->, blu] (r.east) -- (a2.west);
  \draw[->, blu] (r.east) -- (a3.west);
  \draw[->, blu] (a2.east) -- (b1.west);
  \draw[->, blu] (b1.east) -- (c1.west);
  \draw[->, blu, dashed] (c1.east) -- node[above, sloped]{\small\ldots} (d1.west);
  \draw[->, blu, dashed] (c1.east) -- node[below, sloped]{\small\ldots} (d2.west);
  \draw[unical, thick] ($(d1.center)+(-0.52,-0.52)$) -- ($(d1.center)+(0.52,0.52)$);
  \draw[unical, thick] ($(d2.center)+(-0.6,-0.42)$) -- ($(d2.center)+(0.6,-0.42)$);
\end{tikzpicture}
\caption{Una porzione dell'albero di gioco del tris. A destra due foglie in cui vince O.}
\end{figure}

\subsection{Un'IA per il tris}

Servono tre elementi:
\begin{enumerate}
  \item un modo per \kw{rappresentare} la scacchiera e lo stato del gioco;
  \item una \kw{funzione di valutazione} per sapere chi vince;
  \item una \kw{strategia di ricerca} nell'albero di gioco.
\end{enumerate}

\paragraph{Funzione di punteggio.}
Il gioco è piccolo e a somma zero, quindi si assegna un punteggio agli stati finali:
\[
  \text{valore}(s) =
  \begin{cases}
    +1 & \text{se vince MAX,}\\
    -1 & \text{se vince MIN,}\\
    \phantom{+}0 & \text{in caso di pareggio.}
  \end{cases}
\]
La grandezza dei numeri non conta: conta solo l'\emph{ordine} ($+1 > 0 > -1$).

\paragraph{Rappresentazione.}
La scacchiera è una matrice $3\times 3$ di numeri: X $\to 1$, O $\to 2$, cella vuota $\to 0$.
Servono una funzione che controlla se c'è un tris su una riga, colonna o diagonale, e una che
controlla se la scacchiera è piena.

\subsection{L'algoritmo Minimax}

\kw{Minimax} è un algoritmo ricorsivo che esplora l'albero di gioco per determinare
la mossa ottimale.
\begin{itemize}
  \item Il giocatore \kw{MAX} vuole \emph{massimizzare} il punteggio: sceglie la mossa con il
        valore più alto restituito dalla chiamata ricorsiva.
  \item Il giocatore \kw{MIN} vuole \emph{minimizzare} il punteggio (che è quello di MAX):
        sceglie la mossa con il valore più basso.
\end{itemize}

\begin{lstlisting}[style=pseudo]
def minimax(board, player):
    if board is leaf node:
        return value(board)
    if player = MAX:
        best = $-\infty$
        for each move available:
            new_board = apply move to board
            val = minimax(new_board, MIN)
            best = max(val, best)
        return best
    else:
        best = $+\infty$
        for each move available:
            new_board = apply move to board
            val = minimax(new_board, MAX)
            best = min(val, best)
        return best
\end{lstlisting}

In formula, il valore minimax di uno stato $s$ è:
\[
  \text{MM}(s) =
  \begin{cases}
    \text{valore}(s) & \text{se $s$ è terminale,}\\
    \max_{s' \in \text{figli}(s)} \text{MM}(s') & \text{se in $s$ muove MAX,}\\
    \min_{s' \in \text{figli}(s)} \text{MM}(s') & \text{se in $s$ muove MIN.}
  \end{cases}
\]

\begin{extra}[esempio numerico]
Un albero a due livelli (MAX alla radice, MIN al livello sotto). Si parte dalle foglie e si
risale: ogni nodo MIN prende il minimo dei figli, la radice MAX il massimo.

\begin{center}
\begin{tikzpicture}[
  level distance=1.5cm,
  level 1/.style={sibling distance=4.3cm},
  level 2/.style={sibling distance=1.25cm},
  every node/.style={font=\sffamily\small},
  max/.style={draw, thick, fill=blu!12, isosceles triangle, isosceles triangle apex angle=60,
              shape border rotate=90, minimum size=0.95cm, inner sep=0pt},
  min/.style={draw, thick, fill=unical!12, isosceles triangle, isosceles triangle apex angle=60,
              shape border rotate=-90, minimum size=0.95cm, inner sep=0pt},
  foglia/.style={draw, minimum size=0.6cm}]
\node[max]{3}
  child{node[min]{3}
    child{node[foglia]{3}} child{node[foglia]{12}} child{node[foglia]{8}}}
  child{node[min]{2}
    child{node[foglia]{2}} child{node[foglia]{4}} child{node[foglia]{6}}}
  child{node[min]{2}
    child{node[foglia]{14}} child{node[foglia]{5}} child{node[foglia]{2}}};
\node[font=\sffamily\bfseries\small, blu]   at (-6.3,0)    {MAX};
\node[font=\sffamily\bfseries\small, unical] at (-6.3,-1.5) {MIN};
\end{tikzpicture}
\end{center}
I tre nodi MIN valgono $\min(3,12,8) = 3$, $\min(2,4,6) = 2$, $\min(14,5,2) = 2$;
la radice vale $\max(3,2,2) = 3$. MAX sceglie quindi la prima mossa.
\end{extra}

\subsection{I limiti di Minimax}

Minimax nella versione base esplora \emph{tutto} l'albero. Già per il tris l'albero completo
ha $549\,946$ nodi (partendo dalla scacchiera vuota). Molti rami però sono inutili:
se un ramo è già sicuramente peggiore di un'alternativa trovata prima, non serve esplorarlo.

\subsection{Potatura Alpha-Beta}

Si tengono traccia di due valori, i migliori trovati finora lungo il percorso:
\begin{itemize}
  \item $\alpha$: il valore migliore (più \emph{alto}) che \kw{MAX} può garantirsi
        al livello corrente o sopra;
  \item $\beta$: il valore migliore (più \emph{basso}) che \kw{MIN} può garantirsi
        al livello corrente o sopra.
\end{itemize}
All'inizio $\alpha = -\infty$ e $\beta = +\infty$. Appena $\alpha \geq \beta$, il nodo corrente
non potrà mai essere raggiunto con gioco ottimale e i figli restanti si possono \kw{potare}
(\texttt{break}).

\begin{lstlisting}[style=pseudo]
def minimax(board, player, alpha, beta):
    if board is leaf node:
        return value(board)
    if player = MAX:
        best = $-\infty$
        for each move available:
            new_board = apply move to board
            val = minimax(new_board, MIN, alpha, beta)
            best = max(val, best)
            alpha = max(best, alpha)
            if alpha >= beta:
                break
        return best
    else:
        best = $+\infty$
        for each move available:
            new_board = apply move to board
            val = minimax(new_board, MAX, alpha, beta)
            best = min(val, best)
            beta = min(best, beta)
            if alpha >= beta:
                break
        return best
\end{lstlisting}

\begin{extra}[alpha-beta sull'esempio numerico]
Riprendiamo l'albero di prima, esplorando i figli da sinistra a destra.
\begin{enumerate}
  \item Primo nodo MIN ($\alpha=-\infty$, $\beta=+\infty$): vede 3, 12, 8 e restituisce 3.
        La radice aggiorna $\alpha = 3$.
  \item Secondo nodo MIN ($\alpha=3$): la prima foglia vale 2, quindi $\beta = 2$.
        Ora $\alpha = 3 \geq \beta = 2$: \kw{potatura}. Le foglie 4 e 6 non vengono visitate.
        \emph{Perché è corretto}: questo nodo MIN varrà al massimo 2, ma MAX ha già
        un'alternativa che vale 3, quindi non sceglierà mai questo ramo.
  \item Terzo nodo MIN ($\alpha=3$): vede 14 ($\beta=14$), 5 ($\beta=5$), 2 ($\beta=2$) e restituisce 2.
\end{enumerate}
Il risultato è 3, identico a Minimax, ma con due foglie in meno. Il risultato di alpha-beta
è \emph{sempre} identico a quello di Minimax: cambia solo quanto lavoro si fa.

\begin{center}
\begin{tikzpicture}[
  level distance=1.5cm,
  level 1/.style={sibling distance=4.3cm},
  level 2/.style={sibling distance=1.25cm},
  every node/.style={font=\sffamily\small},
  max/.style={draw, thick, fill=blu!12, isosceles triangle, isosceles triangle apex angle=60,
              shape border rotate=90, minimum size=0.95cm, inner sep=0pt},
  min/.style={draw, thick, fill=unical!12, isosceles triangle, isosceles triangle apex angle=60,
              shape border rotate=-90, minimum size=0.95cm, inner sep=0pt},
  foglia/.style={draw, minimum size=0.6cm},
  potata/.style={draw, dashed, gray, minimum size=0.6cm, text=gray}]
\node[max]{3}
  child{node[min]{3}
    child{node[foglia]{3}} child{node[foglia]{12}} child{node[foglia]{8}}}
  child{node[min]{$\le 2$}
    child{node[foglia]{2}}
    child[gray, dashed]{node[potata]{4}}
    child[gray, dashed]{node[potata]{6}}}
  child{node[min]{2}
    child{node[foglia]{14}} child{node[foglia]{5}} child{node[foglia]{2}}};
\end{tikzpicture}
\end{center}
\end{extra}

\begin{osservazione}[l'ordine delle mosse conta]
Quanto si pota dipende dall'ordine in cui si esplorano i figli. Se le mosse migliori vengono
esplorate per prime, $\alpha$ e $\beta$ si stringono subito e si pota molto di più. Nel caso
ideale alpha-beta visita circa $b^{m/2}$ nodi invece di $b^{m}$ ($b$ = numero di mosse per
stato, $m$ = profondità): a parità di tempo si può cercare circa due volte più in profondità.
\end{osservazione}

\subsection{Oltre alpha-beta}

Alpha-beta può ridurre l'albero di oltre il 99\%, ma alcuni giochi restano troppo grandi
(scacchi, go, \ldots). Tecniche usate in pratica:
\begin{itemize}
  \item \kw{ricerca a profondità fissa} / \en{iterative deepening}: ci si ferma a una certa
        profondità, aumentandola finché c'è tempo;
  \item \kw{valutazione euristica} degli stati: una stima del valore di stati non terminali;
  \item \en{quiescence search}: si continua a cercare oltre il limite solo nelle posizioni
        ``instabili'' (es.\ in mezzo a uno scambio di pezzi);
  \item \en{hashing} e \en{transposition tables}: si memorizza il valore di stati già visti,
        raggiungibili con sequenze di mosse diverse;
  \item \en{move ordering}: si provano prima le mosse promettenti, per potare di più;
  \item \kw{basi di conoscenza}: aperture e tabelle dei finali negli scacchi.
\end{itemize}

% =====================================================================
\partelab{Minimax per il tris}
% =====================================================================
\fonte{Lab\_Minimax\_blank.ipynb}

\textbf{Obiettivo}: implementare Minimax (e poi alpha-beta) per il tris, farlo giocare
contro sé stesso e contro un umano. Il codice qui sotto è una soluzione completa ed è stata
verificata: le asserzioni del notebook passano e le due IA pareggiano sempre.

\subsection{Rappresentazione e condizioni di vittoria}

La scacchiera è una lista di liste $3\times3$: \texttt{0} = vuota, \texttt{1} = X, \texttt{2} = O.
\texttt{check\_winner} restituisce \texttt{1} o \texttt{2} se quel giocatore ha un tris,
altrimenti \texttt{0}. \texttt{is\_full} dice se non ci sono celle vuote.

\begin{lstlisting}
def check_winner(board):
    lines = []
    lines.extend(board)                                    # righe
    lines.extend([[board[i][j] for i in range(3)]
                  for j in range(3)])                      # colonne
    lines.append([board[i][i] for i in range(3)])          # diagonale
    lines.append([board[i][2 - i] for i in range(3)])      # antidiagonale
    for line in lines:
        if line[0] != 0 and line[0] == line[1] == line[2]:
            return line[0]
    return 0

def is_full(board):
    return all(cell != 0 for row in board for cell in row)
\end{lstlisting}

\begin{attenzione}[la cella vuota]
Il controllo \texttt{line[0] != 0} è essenziale: senza, una riga vuota \texttt{[0,0,0]}
verrebbe vista come ``tre uguali'' e la funzione restituirebbe un vincitore inesistente.
\end{attenzione}

\subsection{Minimax}

X è MAX, O è MIN. La funzione \texttt{score} traduce il vincitore nel punteggio $+1/-1/0$.
Invece di copiare la scacchiera a ogni mossa, si \emph{modifica sul posto} e poi si
\kw{annulla la mossa} (\en{rollback}) prima di provare la successiva.

\begin{lstlisting}
from math import inf

def score(board):
    w = check_winner(board)
    if w == 1:
        return 1       # vince X (MAX)
    if w == 2:
        return -1      # vince O (MIN)
    return 0           # pareggio

def minimax(board, is_maximizing):
    if check_winner(board) or is_full(board):   # nodo foglia
        return score(board)
    if is_maximizing:
        best = -inf
        for i in range(3):
            for j in range(3):
                if board[i][j] == 0:
                    board[i][j] = 1                          # prova la mossa
                    best = max(best, minimax(board, False))
                    board[i][j] = 0                          # annulla la mossa
        return best
    else:
        best = inf
        for i in range(3):
            for j in range(3):
                if board[i][j] == 0:
                    board[i][j] = 2
                    best = min(best, minimax(board, True))
                    board[i][j] = 0
        return best
\end{lstlisting}

\begin{attenzione}[pseudocodice del notebook]
Nel notebook, nel ramo MAX, il \texttt{return best} è indentato \emph{dentro} il ciclo
\texttt{for}: così la funzione restituirebbe il valore dopo aver provato solo la prima mossa.
Il \texttt{return} va dopo il ciclo, come nelle slide e nel codice qui sopra.
\end{attenzione}

\subsection{Scegliere la mossa migliore}

\texttt{minimax} restituisce un \emph{valore}, non una mossa. \texttt{find\_best\_move} prova
ogni mossa legale, valuta la posizione risultante con \texttt{minimax} e tiene quella con il
valore più alto (se muove il giocatore 1) o più basso (se muove il giocatore 2).

\begin{lstlisting}
def find_best_move(board, player):
    move = (-1, -1)                      # nessuna mossa se la scacchiera è piena
    best = -inf if player == 1 else inf
    for i in range(3):
        for j in range(3):
            if board[i][j] == 0:
                board[i][j] = player
                # dopo la mossa di 'player' tocca all'avversario:
                # se ha mosso X (1) ora minimizza O, e viceversa
                val = minimax(board, player == 2)
                board[i][j] = 0
                if (player == 1 and val > best) or (player == 2 and val < best):
                    best, move = val, (i, j)
    return move
\end{lstlisting}

\begin{attenzione}[errore tipico]
Dopo aver piazzato la mossa di \texttt{player} bisogna chiamare \texttt{minimax} con il turno
dell'\emph{avversario}. Chiamarla con \texttt{is\_maximizing = (player == 1)} è un errore comune:
l'IA valuterebbe le posizioni come se toccasse di nuovo a lei.
\end{attenzione}

\subsection{Minimax contro sé stesso}

Il notebook fornisce il ciclo di gioco: i giocatori si alternano chiamando
\texttt{find\_best\_move} finché qualcuno vince o la scacchiera è piena. Con due giocatori
ottimali la partita finisce \kw{sempre in pareggio}. Con l'implementazione sopra la partita è:
\[
(0,0),\ (1,1),\ (0,1),\ (0,2),\ (2,0),\ (1,0),\ (1,2),\ (2,1),\ (2,2)
\quad\Rightarrow\quad \text{pareggio}
\]

\begin{center}
\tris[0.55]{X,X,O, O,O,X, X,O,X}
\end{center}

\begin{extra}[perché il pareggio]
In ogni gioco combinatorio finito si può calcolare il valore di ogni stato risalendo dalle foglie
(\kw{induzione all'indietro}, \en{backward induction}), che è esattamente ciò che fa Minimax.
Ne segue che uno dei due giocatori ha una strategia che garantisce la vittoria, oppure entrambi
possono garantirsi almeno il pareggio (teorema di Zermelo). Per il tris vale la seconda:
il valore minimax della scacchiera vuota è $0$.
\end{extra}

\subsection{Minimax con potatura alpha-beta}

Si aggiungono i parametri \texttt{alpha} e \texttt{beta} e il controllo di potatura dopo ogni
figlio. Il resto del codice (incluso \texttt{find\_best\_move}) non cambia, salvo chiamare
\texttt{minimax\_ab} al posto di \texttt{minimax}.

\begin{lstlisting}
def minimax_ab(board, is_maximizing, alpha=-inf, beta=inf):
    if check_winner(board) or is_full(board):
        return score(board)
    if is_maximizing:
        best = -inf
        for i in range(3):
            for j in range(3):
                if board[i][j] == 0:
                    board[i][j] = 1
                    best = max(best, minimax_ab(board, False, alpha, beta))
                    board[i][j] = 0
                    alpha = max(alpha, best)
                    if alpha >= beta:
                        return best              # potatura
        return best
    else:
        best = inf
        for i in range(3):
            for j in range(3):
                if board[i][j] == 0:
                    board[i][j] = 2
                    best = min(best, minimax_ab(board, True, alpha, beta))
                    board[i][j] = 0
                    beta = min(beta, best)
                    if alpha >= beta:
                        return best              # potatura
        return best
\end{lstlisting}

\begin{attenzione}[\texttt{break} con due cicli annidati]
Qui le mosse si scorrono con \emph{due} cicli (\texttt{i} e \texttt{j}). Un \texttt{break}
uscirebbe solo dal ciclo interno e l'esterno continuerebbe a esplorare le righe successive,
annullando parte della potatura. Per questo si usa \texttt{return best} (va bene perché la mossa
è già stata annullata). In alternativa si può prima costruire la lista delle celle vuote e
usare un solo ciclo.
\end{attenzione}

\paragraph{Risultati misurati.}
Stessa partita (pareggio), stesse mosse, ma molto meno lavoro:

\begin{center}
\renewcommand{\arraystretch}{1.2}
\begin{tabular}{lrr}
\toprule
 & Minimax & Alpha-beta \\
\midrule
Nodi visitati dalla scacchiera vuota & $549\,946$ & $18\,297$ \\
Nodi visitati in un'intera partita IA contro IA & $618\,175$ & $36\,684$ \\
Tempo della partita (Python) & circa $2{,}9$\,s & circa $0{,}17$\,s \\
\bottomrule
\end{tabular}
\end{center}

Alpha-beta visita circa il 3\% dei nodi di Minimax, cioè ne pota circa il 97\%.

\subsection{Giocare contro l'IA}

L'ultima cella del notebook permette di sfidare l'algoritmo inserendo le mosse come
\texttt{riga colonna} (indici da 0). Se l'implementazione è corretta \kw{non si può mai vincere}:
al massimo si pareggia.

% ---------------------------------------------------------------------
\begin{riepilogo}
\begin{itemize}
  \item Teoria dei giochi = decisioni tra più agenti razionali; AGT = progettare algoritmi
        in ambienti strategici.
  \item Aree del corso: giochi non cooperativi (equilibrio di Nash), coalizionali (valore di
        Shapley), scelta sociale (paradosso di Condorcet), mechanism design, matching.
  \item Dilemma del prigioniero: efficienza (Pareto) e stabilità (Nash) possono non coincidere.
  \item Giochi combinatori: 2 giocatori, informazione perfetta, niente caso, finiti, a turni,
        a somma zero. Si rappresentano con l'albero di gioco.
  \item Minimax: valore di uno stato = max (MAX) o min (MIN) dei valori dei figli.
  \item Alpha-beta: stesso risultato, pota i rami con $\alpha \ge \beta$. Nel tris visita
        $18\,297$ nodi invece di $549\,946$.
\end{itemize}
\textbf{Prossima lezione}: teoria dell'utilità e decisioni in condizioni di incertezza;
giochi strategici e formalizzazione della \en{best response}.
\end{riepilogo}

% =====================================================================
%  Lezione 2 — Teoria dell'utilità e decisioni in condizioni di incertezza
%  Materiale: Lesson_2_Utility_Theory.pdf, Battleship_Blank.ipynb
% =====================================================================
\chapter{Teoria dell'utilità}

\parteteoria

Nella lezione precedente le utilità erano numeri dati ($+1$, $0$, $-1$ per il tris).
Qui si fa il passo indietro: \emph{da dove vengono} quei numeri? La risposta è che
non sono un dato primitivo, ma una \kw{rappresentazione} delle preferenze del giocatore.
Il percorso della lezione è:
\begin{enumerate}[label=\arabic*.]
  \item concetti di base: alternative e preferenze;
  \item rappresentare preferenze ordinali (utilità \kw{ordinale});
  \item oltre l'ordinale: lotterie e utilità \kw{cardinale} (VNM);
  \item atteggiamento verso il rischio;
  \item applicazioni e introduzione alla teoria delle decisioni.
\end{enumerate}

% ---------------------------------------------------------------------
\section{Concetti di base}

\subsection{Alternative}

Un agente sceglie all'interno di un insieme $X$ di \kw{alternative}, che devono essere:
\begin{itemize}
  \item \kw{mutuamente esclusive}: se ne realizza una sola;
  \item \kw{esaustive}: coprono tutti i casi possibili.
\end{itemize}

\begin{esempio}[scelta dei corsi]
Le opzioni sono $\{\text{Deep Learning},\ \text{Algorithmic Game Theory}\}$, ma queste
\emph{non} sono le alternative: si possono seguire entrambi i corsi o nessuno.
L'insieme corretto è
\[
  X = \{\, \mathrm{DL},\ \mathrm{AGT},\ \mathrm{DLAGT},\ \mathrm{N} \,\},
\]
dove DLAGT = entrambi e N = nessuno dei due.
\end{esempio}

\begin{osservazione}[è un passo di modellazione, non una formalità]
Costruire $X$ \emph{bene} è la prima cosa da fare in un modello. Con $n$ opzioni
combinabili liberamente le alternative sono $2^n$: già qui si vede che la dimensione
dello spazio delle alternative è un problema \emph{computazionale}, non solo concettuale.
\end{osservazione}

\subsection{Preferenze}

\begin{definizione}[Relazione di preferenza]
Le \kw{preferenze} sono una relazione $\pref$ su $X$, cioè un sottoinsieme di $X \times X$.
Si legge $x \pref y$ come ``$x$ è \emph{almeno buono quanto} $y$''. La relazione è:
\begin{itemize}
  \item \kw{completa} se $\forall x,y \in X$ vale $x \pref y$ oppure $y \pref x$;
  \item \kw{transitiva} se $\forall x,y,z \in X$, $[x \pref y \text{ e } y \pref z]
        \Rightarrow x \pref z$.
\end{itemize}
Una preferenza è una \kw{relazione di preferenza} se e solo se è completa e transitiva.
\end{definizione}

Da $\pref$ si derivano altre due relazioni:
\begin{align*}
  \text{preferenza \kw{stretta}:} && x \spref y &\iff [\,x \pref y \ \text{e}\ y \npref x\,];\\
  \text{\kw{indifferenza}:}       && x \indif y &\iff [\,x \pref y \ \text{e}\ y \pref x\,].
\end{align*}

\begin{osservazione}[completezza non vuol dire ``decidere sempre'']
La completezza non dice che l'agente preferisce sempre qualcosa a qualcos'altro: dice che
sa sempre \emph{confrontare} due alternative. Il caso ``mi sono indifferenti'' è ammesso
ed è proprio $x \indif y$. Quello che la completezza esclude è l'incomparabilità
(``non so dire'').
\end{osservazione}

% ---------------------------------------------------------------------
\section{Rappresentare le preferenze: utilità ordinale}

\begin{definizione}[Funzione di utilità]
Una relazione di preferenza $\pref$ è \kw{rappresentata} da una funzione di utilità
$u : X \to \R$ nel senso seguente:
\[
  x \pref y \iff u(x) \geq u(y) \qquad \forall x,y \in X.
\]
\end{definizione}

Se un giocatore trova $x$ almeno buono quanto $y$, allora $u(x)$ deve essere almeno
alto quanto $u(y)$. In questo senso il giocatore si comporta \emph{come se}
(\en{as if}) stesse massimizzando $u(\cdot)$.

\begin{osservazione}[la direzione della spiegazione]
L'ultima frase va presa alla lettera e sarà ripresa alla fine della lezione: l'utilità
è un \emph{modo di descrivere} le preferenze, non la loro causa. Non si sceglie $x$
perché ha utilità più alta; ha utilità più alta perché è ciò che si sceglie.
\end{osservazione}

\begin{teorema}[Rappresentazione ordinale]
Sia $X$ \emph{finito}. Una preferenza può essere rappresentata da una funzione di utilità
se e solo se è completa e transitiva. Inoltre, se $u : X \to \R$ rappresenta $\pref$ e
$f : \R \to \R$ è \kw{strettamente crescente}, allora anche $f \circ u$ rappresenta $\pref$.
\end{teorema}

Per l'ultima affermazione queste utilità si dicono \kw{ordinali}: conta solo
l'\emph{informazione d'ordine}.

\subsection{Le condizioni sono necessarie}

\paragraph{Transitività.}
Sia $X = \{a,b,c\}$ e si supponga il ciclo $a \spref b \spref c \spref a$.
Se esistesse $u$ si avrebbe $u(a) > u(b) > u(c) > u(a)$, cioè $u(a) > u(a)$: assurdo.

\begin{extra}[la \en{money pump}]
Filosofi ed economisti direbbero che violare la transitività genera una
\kw{pompa di soldi}. Se l'agente ha il ciclo sopra e possiede $c$, allora preferisce
strettamente $b$ a $c$: gli si può vendere $b$ in cambio di $c$ più un piccolo
importo $\varepsilon$. Poi preferisce $a$ a $b$: altro scambio, altro $\varepsilon$.
Poi preferisce $c$ ad $a$: terzo scambio. L'agente è tornato al punto di partenza
avendo pagato $3\varepsilon$, e il ciclo si può ripetere all'infinito.
È l'argomento ``pratico'' a favore della transitività: chi la viola si fa svuotare le tasche.
\end{extra}

\paragraph{Completezza.}
Si assuma di avere una funzione di utilità $u : X \to \R$. Dati due numeri $u(x)$ e $u(y)$,
sono ammissibili solo due possibilità: $u(x) \geq u(y)$ oppure $u(y) \geq u(x)$.
Quindi $x \pref y$ o $y \pref x$: la completezza è una necessità matematica, ereditata
dal fatto che l'ordine $\geq$ su $\R$ è totale.

\subsection{Le condizioni sono sufficienti}

\paragraph{Dimostrazione del teorema.}
Sia $X = \{x_1,\dots,x_n\}$. Grazie a completezza e transitività si possono partizionare
gli elementi di $X$ in $k$ \kw{classi di indifferenza} $C_1, \dots, C_k$ (ogni classe
raccoglie alternative tra loro indifferenti) ordinate in modo che
$C_1 \spref C_2 \spref \dots \spref C_k$. Si definisce allora
\[
  u(x) = k - i + 1 \qquad \text{per ogni } x \in C_i,
\]
cioè $u = k$ sulla classe migliore, $u = k-1$ sulla successiva, fino a $u = 1$ su $C_k$.
Per costruzione $u(x) \ge u(y) \iff x \pref y$. $\square$

\begin{attenzione}[dove serve che $X$ sia finito]
L'ipotesi ``$X$ finito'' non è decorativa: serve a poter costruire le classi di
indifferenza e a numerarle. Con $X$ infinito completezza e transitività \emph{non}
bastano. Il controesempio classico sono le preferenze \kw{lessicografiche} su
$X = \R^2$, dove $(a_1,a_2) \pref (b_1,b_2)$ se $a_1 > b_1$, oppure se $a_1 = b_1$ e
$a_2 \ge b_2$: la relazione è completa e transitiva ma non ammette alcuna
rappresentazione con una $u$ a valori reali.
\end{attenzione}

\subsection{Esempi: l'utilità ordinale non è interpretabile}

Con $X = \{a,b,c\}$ e preferenze $a \spref b \spref c$, tutte e tre queste funzioni
rappresentano le \emph{stesse} preferenze:
\begin{center}
\renewcommand{\arraystretch}{1.15}
\begin{tabular}{cccc}
\toprule
 & $u_1$ & $u_2$ & $u_3$ \\
\midrule
$u(a)$ & $-1$ & $30$ & $90$ \\
$u(b)$ & $-2$ & $20$ & $2$ \\
$u(c)$ & $-3$ & $10$ & $1$ \\
\bottomrule
\end{tabular}
\end{center}
Il valore numerico dell'utilità ordinale \kw{non è interpretabile}: non dice
con quale \emph{intensità} un giocatore preferisce $a$ a $b$. Con $u_2$ le
differenze sono uguali, con $u_3$ il salto tra $a$ e $b$ sembra enorme, ma le
preferenze rappresentate sono identiche.

\begin{osservazione}[perché non basta]
Finché si sceglie tra alternative \emph{certe}, l'ordine basta: si prende la migliore.
Ma appena entra in gioco l'\kw{incertezza} (una scelta dà $a$ con probabilità $0{,}5$ e
$c$ con probabilità $0{,}5$) serve sapere \emph{quanto} $b$ sta tra $a$ e $c$: serve
un'utilità \kw{cardinale}. È il problema che risolve la teoria di Von Neumann--Morgenstern.
\end{osservazione}

% ---------------------------------------------------------------------
\section{Oltre l'ordinale: le lotterie}

\begin{definizione}[Lotteria semplice]
Una \kw{lotteria semplice} è una tupla
\[
  L = (p_1, x_1;\ p_2, x_2;\ \dots;\ p_n, x_n)
\]
dove $x_1,\dots,x_n \in X \subseteq \R$ sono \kw{premi monetari} e
$(p_1,\dots,p_n)$ è una distribuzione di probabilità ($p_i \ge 0$, $\sum_i p_i = 1$).
Si indica con $\lott$ l'insieme delle lotterie semplici.
\end{definizione}

\begin{esempio}[una lotteria]
$L = (0{,}3,\ 10;\ 0{,}2,\ 5;\ 0{,}1,\ 0;\ 0{,}4,\ -5)$.
\begin{center}
\begin{tikzpicture}[>=Stealth, scale=0.9, every node/.style={font=\small\sffamily}]
  \node[circle, draw, inner sep=1.6pt] (r) at (0,0) {};
  \foreach \p/\x/\y [count=\k] in {0.3/10/1.5, 0.2/5/0.5, 0.1/0/-0.5, 0.4/-5/-1.5} {
    \node[circle, fill=black, inner sep=1.8pt] (n\k) at (2.6,\y) {};
    \draw[->, grigio] (r) -- node[above, pos=0.62, font=\scriptsize\sffamily]{\p} (n\k);
    \node[right=4pt of n\k] {$\x$};
  }
\end{tikzpicture}
\end{center}
\end{esempio}

\subsection{Lotterie composte}

Una lotteria può avere come premi altre lotterie. Siano
\[
  L_1 = (0{,}9{:}\ -1\$,\ 0{,}1{:}\ 5\$), \qquad
  L_2 = (0{,}75{:}\ -1\$,\ 0{,}25{:}\ 2\$).
\]
La \kw{lotteria composta} $L_c = (0{,}5{:}\ L_1,\ 0{,}5{:}\ L_2)$ equivale alla lotteria
semplice che si ottiene moltiplicando le probabilità lungo i rami e sommando quelle
che portano allo stesso premio:
\[
  L_e = \bigl(0{,}5 \times 0{,}9 + 0{,}5 \times 0{,}75 {:}\ -1\$,\quad
              0{,}5 \times 0{,}25 {:}\ 2\$,\quad
              0{,}5 \times 0{,}1 {:}\ 5\$\bigr)
      = (0{,}825{:}\ -1\$,\ 0{,}125{:}\ 2\$,\ 0{,}05{:}\ 5\$).
\]

\begin{osservazione}[ipotesi di riduzione]
Questo passaggio è un'\emph{assunzione} del modello: si suppone che al giocatore
interessi solo la distribuzione finale sui premi, non il modo in cui ci si arriva
(due estrazioni invece di una). Grazie a essa basta studiare $\lott$, l'insieme delle
lotterie \emph{semplici}.
\end{osservazione}

\subsection{Valore atteso}

\begin{definizione}[Valore atteso]
Per $L = (p_1,x_1;\dots;p_n,x_n)$ il \kw{valore atteso} è
\[
  \Ev(L) = \sum_{i=1}^{n} p_i\, x_i .
\]
\end{definizione}

Per la lotteria dell'esempio:
\[
  \Ev(L) = 0{,}3 \times 10 + 0{,}2 \times 5 + 0{,}1 \times 0 + 0{,}4 \times (-5) = 2 .
\]
Giocando ripetutamente quella lotteria ci si aspetta di vincere in media $2$.

\begin{attenzione}[il valore atteso non basta a decidere]
Verrebbe da usare $\Ev$ come criterio di scelta, ma è sbagliato. Si consideri il
\kw{paradosso di San Pietroburgo}: si lancia una moneta finché non esce testa; se esce
al lancio $k$ si vincono $2^k$ euro. Il valore atteso è
$\sum_{k \ge 1} 2^{-k} \cdot 2^{k} = \infty$, eppure nessuno pagherebbe più di poche
decine di euro per giocare. Serve applicare l'utilità \emph{ai premi}, non ai soldi.
\end{attenzione}

\subsection{Utilità attesa}

\begin{definizione}[Utilità attesa e funzione di Bernoulli]
Data una relazione di preferenza $\pref$ su $\lott$, una funzione di utilità
$U : \lott \to \R$ è una \kw{funzione di utilità attesa} se può essere scritta come
\[
  U(L) = \sum_{i=1}^{n} p_i\, u(x_i)
\]
per qualche funzione $u : \R \to \R$. La funzione $u$ si chiama
\kw{funzione di utilità di Bernoulli}.
\end{definizione}

\begin{attenzione}[due utilità diverse, due simboli diversi]
$u$ è definita sui \emph{premi} (numeri di euro), $U$ sulle \emph{lotterie}.
La forma di $U$ è vincolata: deve essere \kw{lineare nelle probabilità}. Non ogni
$U$ che rappresenta preferenze su $\lott$ ha questa forma; il teorema VNM dice
esattamente quando ce l'ha.
\end{attenzione}

% ---------------------------------------------------------------------
\section{L'utilità di Von Neumann--Morgenstern}

Von Neumann e Morgenstern (VNM) hanno fornito le condizioni sotto cui una relazione di
preferenza $\pref$ su $\lott$ è rappresentabile da una funzione di utilità attesa.

\begin{definizione}[Assiomi VNM]
\begin{description}[leftmargin=2.2cm, style=nextline, font=\normalfont\bfseries]
  \item[Assioma 1] (\kw{ordine di preferenza}) $\pref$ è completa e transitiva.
  \item[Assioma 2] (\kw{continuità}) se $L \spref M \spref N$, esiste $p \in [0,1]$
        tale che $pL + (1-p)N \indif M$.
  \item[Assioma 3] (\kw{indipendenza}) per ogni lotteria $N$ e ogni $p \in [0,1]$,
        \[ L \pref M \iff pL + (1-p)N \pref pM + (1-p)N. \]
\end{description}
\end{definizione}

\begin{osservazione}[come leggere gli assiomi 2 e 3]
La \kw{continuità} dice che non esistono premi ``infinitamente'' migliori di altri:
se $M$ sta tra $L$ e $N$, c'è una miscela di $L$ e $N$ esattamente equivalente a $M$.
L'\kw{indipendenza} dice che un esito \emph{comune} alle due lotterie non deve
cambiare il confronto: $N$ compare con la stessa probabilità $1-p$ da entrambe le parti,
quindi deve essere irrilevante. È l'assioma più forte dei tre e, come si vedrà con
Allais, è quello che gli esseri umani violano.
\end{osservazione}

\begin{teorema}[Von Neumann--Morgenstern]
Una relazione binaria $\pref$ su $\lott$ ammette una rappresentazione di utilità attesa
se e solo se soddisfa gli assiomi 1--3. Inoltre, se $U$ e $V$ sono due rappresentazioni
di utilità attesa di $\pref$, allora esistono costanti $a, b \in \R$ con $a > 0$ tali che
\[
  U(\cdot) = a\,V(\cdot) + b .
\]
\end{teorema}

L'ultima affermazione dice che la rappresentazione VNM è unica \kw{a meno di
trasformazioni affini}.

\begin{osservazione}[ordinale contro cardinale]
È qui la differenza con il Teorema 1. Lì si poteva applicare a $u$ \emph{qualunque}
funzione strettamente crescente; qui solo $a x + b$ con $a>0$. Il margine di libertà è
molto più stretto, e questo rende i numeri \emph{interpretabili}: sono fissati l'origine
e l'unità di misura, come per la temperatura in gradi Celsius. Ha senso dire che il salto
da $a$ a $b$ è il doppio del salto da $b$ a $c$; non ha senso dire che $a$ vale il doppio
di $b$. Si parla di \kw{scala a intervalli}.
\end{osservazione}

\subsection{Necessità degli assiomi}

\paragraph{Ordine di preferenza.} Già dimostrato per il Teorema 1.

\paragraph{Continuità.}
Sia $L \spref M \spref N$; allora $u(L) > u(M) > u(N)$. Si ponga
\[
  p = \frac{u(M) - u(N)}{u(L) - u(N)} \in (0,1).
\]
Per la definizione di utilità attesa:
\begin{align*}
  u\bigl(pL + (1-p)N\bigr)
    &= p\,u(L) + (1-p)\,u(N)\\
    &= p\,u(L) + u(N) - p\,u(N)\\
    &= p\bigl(u(L) - u(N)\bigr) + u(N)\\
    &= \frac{u(M) - u(N)}{u(L) - u(N)}\bigl(u(L) - u(N)\bigr) + u(N)\\
    &= u(M) - u(N) + u(N) = u(M).
\end{align*}
Quindi $pL + (1-p)N \indif M$: la continuità è necessaria.

\paragraph{Indipendenza.}
Siano $L_1 = pL + (1-p)N$ e $L_2 = pM + (1-p)N$. Per linearità dell'attesa:
\begin{align*}
  u(L_1) &= p\,u(L) + (1-p)\,u(N),\\
  u(L_2) &= p\,u(M) + (1-p)\,u(N).
\end{align*}
Sottraendo membro a membro, il termine comune in $N$ si cancella:
\[
  u(L_1) - u(L_2) = p\bigl(u(L) - u(M)\bigr).
\]
Se $L \spref M$ allora $u(L) > u(M)$, dunque $u(L_1) > u(L_2)$, cioè $L_1 \spref L_2$.
L'indipendenza è necessaria.

\begin{osservazione}[da dove viene la ``forza'' dell'indipendenza]
La cancellazione del termine $(1-p)u(N)$ funziona \emph{perché} $U$ è lineare nelle
probabilità. Indipendenza e linearità sono in sostanza la stessa proprietà vista da due
lati: uno comportamentale (l'esito comune non conta), l'altro algebrico.
\end{osservazione}

\subsection{Sufficienza: dimostrazione del teorema VNM}

La dimostrazione si divide in due parti:
\begin{itemize}
  \item \kw{Parte 1}: si mostra che $U(L) = \sum_{i} p_i\,u(x_i)$;
  \item \kw{Parte 2}: si mostra che $L \spref M \iff U(L) > U(M)$ per ogni $L,M \in \lott$.
\end{itemize}

\paragraph{Parte 1.}
Si supponga di avere $n$ esiti $o_1,\dots,o_n$. Per completezza e transitività si possono
ordinare dal peggiore al migliore: $o_1 \pref \dots \pref o_n$ (nel senso $o_n$ migliore).
Si fissano i due estremi della scala:
\[
  u(o_1) = 0, \qquad u(o_n) = 1 .
\]
Si definisce la famiglia di lotterie semplici
\[
  L(p) = p \cdot o_n + (1-p) \cdot o_1 ,
\]
cioè ``il premio migliore con probabilità $p$, il peggiore con probabilità $1-p$''.
Per l'assioma di continuità, per ogni esito $o_i$ esiste una probabilità $q_i \in [0,1]$
tale che $L(q_i) \indif o_i$; si pone $u(o_i) = q_i$. Ne segue che l'utilità di una
lotteria $M = \sum_i p_i o_i$ è l'attesa di $u$:
\[
  u(M) = u\Bigl(\sum_i p_i o_i\Bigr) = \sum_i p_i\,u(o_i) = \sum_i p_i\,q_i .
\]

\begin{osservazione}[che cosa è davvero $q_i$]
$q_i$ è un modo operativo di \emph{misurare} l'utilità: è la probabilità di vincere il
premio migliore che rende il giocatore indifferente a ricevere $o_i$ con certezza.
È una definizione che si può realizzare come esperimento, chiedendo alla persona
di aggiustare $q$ finché non è indifferente, ed è esattamente questo che rende la scala
cardinale invece che ordinale: $q_i$ posiziona $o_i$ \emph{tra} i due estremi.
\end{osservazione}

\paragraph{Parte 2.}
Per l'assioma di continuità il giocatore è indifferente tra l'esito $o_i$ e la lotteria
$q_i \cdot o_n + (1-q_i) \cdot o_1$. Insieme all'assioma di indipendenza (si sostituisce
un esito alla volta dentro la lotteria) il giocatore è indifferente tra le due lotterie
\[
  M = \sum_i p_i o_i
  \qquad\text{e}\qquad
  M' = \sum_i p_i \bigl[\,q_i \cdot o_n + (1-q_i)\cdot o_1\,\bigr].
\]
Ma $M'$ coinvolge solo i due esiti estremi, quindi si può riscrivere come lotteria semplice:
\[
  M' = \Bigl(\sum_i p_i q_i\Bigr)\cdot o_n + \Bigl(\sum_i p_i (1-q_i)\Bigr)\cdot o_1
     = U(M)\cdot o_n + \bigl(1 - U(M)\bigr)\cdot o_1 .
\]
Ogni lotteria è dunque equivalente a una lotteria sui soli due estremi, la cui unica
caratteristica è la probabilità $U(M)$ di ottenere il premio migliore.

Si supponga ora $L \spref M$ e si definiscano
\begin{align*}
  L' &= U(L)\cdot o_n + \bigl(1-U(L)\bigr)\cdot o_1,\\
  M' &= U(M)\cdot o_n + \bigl(1-U(M)\bigr)\cdot o_1 .
\end{align*}
Vale l'ordine $L' \indif L \spref M \indif M'$, quindi $L' \spref M'$. Ma $L'$ e $M'$
differiscono solo per la probabilità del premio migliore, e tra due lotterie di questa
famiglia si preferisce quella con la probabilità più alta: dunque $U(L) > U(M)$.
Vale anche il viceversa, quindi $L \spref M \iff U(L) > U(M)$. $\square$

% ---------------------------------------------------------------------
\section{Il paradosso di Allais}

Si risponda d'istinto a due domande. \emph{Prima}: quale lotteria si preferisce?
\[
  L = (1{:}\ 1.000.000\$), \qquad
  M = (0{,}89{:}\ 1.000.000\$,\ \ 0{,}10{:}\ 5.000.000\$,\ \ 0{,}01{:}\ 0\$).
\]
\emph{Seconda}: e tra queste?
\[
  N = (0{,}89{:}\ 0\$,\ \ 0{,}11{:}\ 1.000.000\$), \qquad
  O = (0{,}90{:}\ 0\$,\ \ 0{,}10{:}\ 5.000.000\$).
\]
La maggior parte delle persone sceglie $L$ nella prima coppia (il milione sicuro) e $O$
nella seconda. Questa combinazione è \kw{incompatibile} con la teoria dell'utilità attesa.

\begin{extra}[perché è una contraddizione]
Si normalizzi $u(0) = 0$ e si scrivano $u_1 = u(1\text{M})$, $u_5 = u(5\text{M})$.
\begin{itemize}
  \item $L \spref M$ significa
        $u_1 > 0{,}89\,u_1 + 0{,}10\,u_5$, cioè $0{,}11\,u_1 > 0{,}10\,u_5$;
  \item $O \spref N$ significa
        $0{,}10\,u_5 > 0{,}11\,u_1$.
\end{itemize}
Le due disuguaglianze si contraddicono: \emph{nessuna} funzione di utilità $u$ può
giustificare entrambe le scelte.

Visto con l'assioma di indipendenza è ancora più netto. Si scriva ogni lotteria come
miscela con $p = 0{,}11$:
\begin{center}
\renewcommand{\arraystretch}{1.25}
\begin{tabular}{lll}
\toprule
 & parte con probabilità $0{,}11$ & parte con probabilità $0{,}89$ \\
\midrule
$L$ & $1$M certo                                  & $1$M \\
$M$ & $(\tfrac{10}{11}{:}\,5\text{M},\ \tfrac{1}{11}{:}\,0)$ & $1$M \\[2pt]
$N$ & $1$M certo                                  & $0$ \\
$O$ & $(\tfrac{10}{11}{:}\,5\text{M},\ \tfrac{1}{11}{:}\,0)$ & $0$ \\
\bottomrule
\end{tabular}
\end{center}
Le coppie $(L,M)$ e $(N,O)$ differiscono nella prima colonna esattamente allo stesso modo;
la seconda colonna è un esito \emph{comune} (dentro ogni coppia). Per l'indipendenza
l'esito comune non deve contare, quindi chi preferisce $L$ a $M$ deve preferire $N$ a $O$.
\end{extra}

\begin{osservazione}[che cosa insegna il paradosso]
Allais non dimostra che la teoria VNM è sbagliata: dimostra che è \kw{normativa} e non
\kw{descrittiva}. Dice come si \emph{dovrebbe} decidere per essere coerenti, non come le
persone decidono davvero. L'effetto in gioco si chiama \en{certainty effect}: il salto da
$0{,}99$ a $1$ (certezza) pesa psicologicamente molto più del salto da $0{,}10$ a $0{,}11$.
Le teorie descrittive alternative (per esempio la \en{prospect theory} di Kahneman e
Tversky) nascono proprio da qui.
\end{osservazione}

% ---------------------------------------------------------------------
\section{Atteggiamento verso il rischio}

\begin{definizione}[Lotteria equa]
Una lotteria $L = (p{:}\ x,\ (1-p){:}\ y)$ è \kw{equa} se il suo valore atteso è nullo:
\[
  \Ev(L) = p\,x + (1-p)\,y = 0 .
\]
\end{definizione}

\begin{attenzione}[notazione delle slide]
Le slide scrivono ``$L = p \cdot x + (1-p)\cdot y = 0$''. La lotteria $L$ è un oggetto
(una distribuzione), non un numero: a essere nullo è il suo \emph{valore atteso}
$\Ev(L)$. Si tratta di un abuso di notazione, ma la sostanza è ``giocare questa lotteria
non costa e non rende nulla, in media''.
\end{attenzione}

\begin{definizione}[Atteggiamenti verso il rischio]
Un giocatore con funzione di utilità di Bernoulli $u$ è:
\begin{itemize}
  \item \kw{neutrale al rischio} se $u$ è \kw{lineare}, $u(x) = ax + b$ (con $a > 0$):
        è indifferente di fronte alle lotterie eque;
  \item \kw{avverso al rischio} se $u$ è \kw{concava}:
        se strettamente concava, non gioca \emph{nessuna} lotteria equa;
  \item \kw{amante del rischio} (\en{risk seeking}) se $u$ è \kw{convessa}:
        se strettamente convessa, gioca \emph{tutte} le lotterie eque.
\end{itemize}
\end{definizione}

\begin{figure}[H]
\centering
% Stessa scala per le tre curve: x_tikz = x/5, y_tikz = u/25.
% avverso  u(x) =  x^0.72  (x>0),  -(-x)^1.4 (x<0)   -> concava
% amante   u(x) =  x^1.4   (x>0),  -(-x)^0.72 (x<0)  -> convessa
% neutrale u(x) =  x
\begin{tikzpicture}[scale=1, >=Stealth, font=\footnotesize\sffamily]
  \draw[->, grigio] (-4.4,0) -- (4.4,0) node[below right]{$x$};
  \draw[->, grigio] (0,-2.9) -- (0,2.9) node[above left]{$u(x)$};
  \foreach \t in {-4,-2,2,4} \draw[grigio] (\t,0.07) -- (\t,-0.07);
  % a sinistra le etichette vanno sopra l'asse: sotto ci passano le curve
  \foreach \t/\l in {-4/-20, -2/-10}
    \node[grigio, anchor=south, font=\scriptsize\sffamily] at (\t,0.1) {\l};
  \foreach \t/\l in {2/10, 4/20}
    \node[grigio, anchor=north, font=\scriptsize\sffamily] at (\t,-0.1) {\l};
  \draw[grigio, thick, dashed] (-4,-0.8) -- (4,0.8);
  \draw[unical, very thick, domain=0:4,  samples=60, smooth] plot (\x, {pow(5*\x,0.72)/25});
  \draw[unical, very thick, domain=-4:0, samples=60, smooth] plot (\x, {-pow(-5*\x,1.4)/25});
  \draw[blu, very thick, domain=0:4,  samples=60, smooth] plot (\x, {pow(5*\x,1.4)/25});
  \draw[blu, very thick, domain=-4:0, samples=60, smooth] plot (\x, {-pow(-5*\x,0.72)/25});
  \node[blu,    anchor=west] at (4.15, 2.64) {amante del rischio};
  \node[grigio, anchor=west] at (4.15, 0.80) {neutrale};
  \node[unical, anchor=west] at (4.15, 0.34) {avverso al rischio};
\end{tikzpicture}
\caption{Le tre forme della funzione di utilità sui guadagni. L'avverso al rischio
soffre una perdita più di quanto goda un guadagno della stessa entità; per l'amante
del rischio vale il contrario.}
\end{figure}

\begin{extra}[perché concavità significa avversione al rischio]
È la \kw{disuguaglianza di Jensen}. Se $u$ è concava, per ogni lotteria $L$
\[
  \Ev\bigl[u(L)\bigr] \;\le\; u\bigl(\Ev[L]\bigr),
\]
cioè l'utilità del gioco è al più l'utilità di incassare con certezza il valore atteso.
Chi è avverso al rischio preferisce quindi sempre il certo all'aleatorio di pari media;
per una lotteria equa aggiunta alla propria ricchezza, non giocare è almeno altrettanto
buono che giocare.

Si definisce \kw{equivalente certo} (\en{certainty equivalent}) $CE(L)$ il premio certo
che dà la stessa utilità della lotteria, cioè $u(CE(L)) = \Ev[u(L)]$. La differenza
\[
  \Ev[L] - CE(L)
\]
è il \kw{premio per il rischio}: quanto il giocatore è disposto a rinunciare per
eliminare l'incertezza. Per un avverso al rischio è positivo, ed è esattamente la
quantità su cui si regge il mercato assicurativo.
\end{extra}

% ---------------------------------------------------------------------
\section{Applicazioni}

In entrambe le applicazioni la funzione di Bernoulli è $u(x) = \sqrt{x}$, che è
strettamente concava: il giocatore è avverso al rischio.

\subsection{Condivisione del rischio (\en{risk sharing})}

Due giocatori avversi al rischio, con $u(x) = \sqrt{x}$, possiedono due attività
rischiose \emph{indipendenti}
\[
  A_1 = A_2 = (0{,}5,\ 100;\ 0{,}5,\ 0).
\]
L'utilità di ciascuna attività, presa da sola, è
\[
  u(A_1) = u(A_2) = 0{,}5 \times \sqrt{100} + 0{,}5 \times \sqrt{0} = 5 .
\]
Se i due giocatori formano un \kw{fondo comune} in cui ciascuno possiede metà del totale,
ogni giocatore si ritrova con l'attività
\[
  A_m = (0{,}25,\ 100;\ 0{,}5,\ 50;\ 0{,}25,\ 0),
\]
la cui utilità è
\[
  u(A_m) = 0{,}25 \times \sqrt{100} + 0{,}5 \times \sqrt{50} + 0{,}25 \times \sqrt{0}
         \approx 2{,}5 + 3{,}54 = 6{,}04 > 5 .
\]
Entrambi i giocatori stanno \kw{meglio} pur non avendo cambiato il valore atteso.

\begin{extra}[da dove viene $A_m$ e perché funziona]
Le due attività sono indipendenti, quindi i quattro casi (vince $A_1$ sì/no, $A_2$ sì/no)
hanno probabilità $0{,}25$ ciascuno. Il montante totale è $200$, $100$, $100$, $0$; metà
del fondo vale quindi $100$, $50$, $50$, $0$, da cui
$A_m = (0{,}25,100;\ 0{,}5,50;\ 0{,}25,0)$.

Il valore atteso è rimasto $50$ in entrambi i casi: quello che è cambiato è la
\emph{varianza}, scesa da $2500$ a $1250$. Mettere insieme rischi indipendenti li
\kw{diversifica}, e per un giocatore con $u$ concava meno varianza a parità di media
significa più utilità. È il principio su cui si reggono i fondi comuni di investimento
e le mutue assicurative.
\end{extra}

\subsection{Assicurazione}

Un giocatore avverso al rischio con $u(x) = \sqrt{x}$ possiede l'attività
$A = (0{,}5,\ 100;\ 0{,}5,\ 0)$. Esiste una compagnia di assicurazione \kw{neutrale al
rischio} con molto capitale.

\paragraph{Quanto è disposto a pagare il giocatore?}
Il giocatore paga un premio $P$ e in cambio riceve con certezza $100$; il suo patrimonio
certo è quindi $100 - P$. Accetta se questo gli dà almeno l'utilità della lotteria:
\[
  u(100 - P) \ \geq\ 0{,}5 \times u(100) + 0{,}5 \times u(0)
  \ \Longrightarrow\ \sqrt{100 - P} \geq 5
  \ \Longrightarrow\ 100 - P \geq 25
  \ \Longrightarrow\ P \leq 75 .
\]

\paragraph{Quanto vuole la compagnia?}
Essendo neutrale al rischio, le basta che il premio copra il valore atteso del risarcimento:
\[
  P \ \geq\ 0{,}5 \times 100 + 0{,}5 \times 0 \ \Longrightarrow\ P \geq 50 .
\]

\paragraph{Conclusione.} Entrambe le parti guadagnano se il contratto si stipula a un
premio
\[
  P \in [50,\ 75].
\]

\begin{osservazione}[lettura con l'equivalente certo]
$\sqrt{100-P} \ge 5$ significa $100 - P \ge 25$: l'equivalente certo della lotteria è
$CE(A) = 25$, mentre il suo valore atteso è $\Ev(A) = 50$. Il premio per il rischio è
$50 - 25 = 25$, ed è esattamente l'ampiezza dell'intervallo $[50,75]$: il
\kw{surplus} che le due parti si dividono. L'assicurazione crea valore perché sposta il
rischio da chi lo odia a chi gli è indifferente.
\end{osservazione}

\begin{attenzione}[il patrimonio iniziale conta]
Qui si è assunto implicitamente che il giocatore non abbia altra ricchezza: l'utilità si
calcola su $100 - P$ e su $100$ o $0$. Se avesse un patrimonio iniziale $w$, le utilità
sarebbero $u(w + 100 - P)$ e $u(w+100)$, $u(w)$, e il premio massimo accettabile sarebbe
diverso (più basso in termini relativi: più si è ricchi, meno si è avversi al rischio
con $u(x) = \sqrt{x}$). La funzione di Bernoulli si applica alla \emph{ricchezza
finale}, non al guadagno della singola scommessa.
\end{attenzione}

% ---------------------------------------------------------------------
\section{Riepilogo sulla teoria VNM}

\begin{itemize}
  \item La teoria VNM fornisce un modo per costruire funzioni di utilità \kw{cardinali}:
        i valori numerici sono interpretabili su una \kw{scala a intervalli} e, in certe
        situazioni, su una scala di rapporti.
  \item La teoria VNM fornisce una teoria \kw{normativa} della decisione: massimizzare
        l'utilità attesa. È molto usata in applicazioni finanziarie ed economiche.
  \item La teoria VNM \kw{non descrive} il comportamento umano. Un'alternativa non è
        preferita perché ha un'utilità più alta; piuttosto, le si associa un'utilità più
        alta perché è l'alternativa preferita.
\end{itemize}

% =====================================================================
\partelab{Battaglia navale e ricerca bayesiana}
% =====================================================================
\fonte{Battleship\_Blank.ipynb}

\textbf{Obiettivo}: costruire un'IA per la \kw{battaglia navale}, partendo da strategie
euristiche semplici e arrivando a una strategia \kw{adattiva} che a ogni turno spara dove
l'\emph{utilità attesa} è massima. Quest'ultima tecnica si chiama \kw{ricerca bayesiana}
(\en{Bayesian search}) ed è uno degli algoritmi più efficaci per i problemi di ricerca e
di ``caccia al tesoro''.

Il collegamento con la teoria della lezione è diretto: non c'è un avversario strategico,
ma c'è \kw{incertezza}. Si deve scegliere l'azione (la cella su cui sparare) che
massimizza l'utilità attesa, dove la probabilità viene aggiornata a ogni colpo con le
informazioni raccolte.

\subsection{Regole e impostazione}

Il campo è una griglia $10 \times 10$ su cui sono disposte cinque navi, in orizzontale o
in verticale, senza sovrapposizioni e senza uscire dal bordo:

\begin{center}
\renewcommand{\arraystretch}{1.15}
\begin{tabular}{lc}
\toprule
Nave & Celle \\
\midrule
Portaerei (\en{Carrier})     & 5 \\
Corazzata (\en{Battleship})  & 4 \\
Incrociatore (\en{Cruiser})  & 3 \\
Sommergibile (\en{Submarine}) & 3 \\
Cacciatorpediniere (\en{Destroyer}) & 2 \\
\midrule
\textbf{Totale}              & \textbf{17} \\
\bottomrule
\end{tabular}
\end{center}

Non essendoci un avversario vero, le strategie non si valutano in partite vinte ma in
\kw{numero di turni necessari ad affondare tutte le navi}: 17 è il minimo teorico
(nessun colpo a vuoto), 100 il massimo.

Lo stato della conoscenza si tiene in una matrice \texttt{shot\_grid} con
\texttt{0} = cella inesplorata, \texttt{1} = colpo a vuoto, \texttt{2} = colpo a segno.

\begin{attenzione}[un bug nel generatore fornito]
La funzione \texttt{create\_random\_disposition} del notebook estrae le coordinate così:
\begin{lstlisting}
if horizontal:
    x = np.random.randint(0, grid_size - 1)        # riga
    y = np.random.randint(0, grid_size - length)   # colonna di partenza
else:
    x = np.random.randint(0, grid_size - length)
    y = np.random.randint(0, grid_size - 1)
\end{lstlisting}
\texttt{np.random.randint(0, n)} esclude \texttt{n}. Quindi per una nave orizzontale la
riga sta in $[0,8]$ (la riga 9 non esce mai) e la colonna di partenza in
$[0,\,9-\texttt{length}]$, cioè l'ultima cella occupata arriva al più alla colonna 8.
Simmetricamente per le navi verticali. Il risultato è che \kw{l'ultima riga e l'ultima
colonna non sono mai occupate}: su 2000 partite generate, le celle della riga 9 e della
colonna 9 risultano vuote in tutte e 2000 (contro circa 2800 occupazioni per la riga 0 e
per la colonna 0). Gli estremi corretti sarebbero
\texttt{randint(0, grid\_size)} e \texttt{randint(0, grid\_size - length + 1)}.

Conseguenza pratica: le cifre misurate sotto sono un po' ottimistiche, perché il 19\% del
campo è di fatto acqua garantita. Il confronto \emph{tra} le strategie resta però valido,
dato che tutte giocano sulle stesse partite.
\end{attenzione}

\subsection{Strategia 1: tiro casuale}

La strategia più semplice: si spara a caso, ignorando del tutto gli esiti. Non è una vera
strategia, ma serve da \kw{riferimento} (\en{baseline}): battere il caso è il minimo
sindacale.

\begin{lstlisting}
def random_attack_strategy(grid):
    shot_grid = np.zeros((10, 10), dtype=int)
    cells = [(i, j) for i in range(10) for j in range(10)]
    np.random.shuffle(cells)
    hits = turns = 0
    for (i, j) in cells:
        shot_grid[i, j] = 2 if grid[i, j] else 1
        turns += 1
        if grid[i, j]:
            hits += 1
            if hits == 17:              # tutte le navi affondate
                break
    return shot_grid, turns
\end{lstlisting}

\begin{osservazione}[mescolare invece di estrarre]
Sparare ``a caso'' si può implementare estraendo una cella e riprovando se è già stata
colpita, ma verso la fine della partita quasi tutte le celle sono esplorate e il ciclo
di riprova diventa lento. Mescolare una volta la lista delle 100 celle e scorrerla dà la
stessa distribuzione senza riprove.
\end{osservazione}

\subsection{Strategia 2: \en{hunt and target}}

La strategia casuale va male perché tratta tutte le celle come equiprobabili e
\emph{butta via} l'informazione dei colpi a segno. L'idea di \kw{hunt and target} è
alternare due modalità:
\begin{itemize}
  \item \kw{hunt} (caccia): si cerca a caso finché non si colpisce qualcosa;
  \item \kw{target} (bersaglio): appena si colpisce, si spara nelle celle adiacenti
        (sopra, sotto, destra, sinistra) non ancora esplorate, finché la nave è affondata.
\end{itemize}
È il classico compromesso tra \kw{esplorazione} e \kw{sfruttamento}.

In più, in modalità hunt non serve sparare ovunque: la nave più piccola occupa 2 celle,
quindi una nave non può nascondersi dentro una sola ``casella nera'' di una
\kw{scacchiera}. Basta sparare sulle celle con $(i+j)$ pari: 50 celle invece di 100, con
la garanzia di intercettare comunque ogni nave.

\begin{center}
\begin{tikzpicture}[scale=0.3]
  \foreach \i in {0,...,9} \foreach \j in {0,...,9} {
    \pgfmathparse{int(mod(\i+\j,2))}
    \ifnum\pgfmathresult=0
      \fill[blu!25] (\j,-\i) rectangle ++(1,1);
    \fi
    \draw[grigio!60] (\j,-\i) rectangle ++(1,1);
  }
  \fill[unical] (3,-4) rectangle ++(1,1);
  \fill[unical] (4,-4) rectangle ++(1,1);
\end{tikzpicture}
\end{center}
{\small\sffamily\color{grigio}\hspace*{\fill}Le celle della scacchiera (in blu) bastano:
qualunque nave da 2 celle (in rosso) ne tocca almeno una.\hspace*{\fill}}

\begin{lstlisting}
def hunt_and_target_strategy(grid):
    shot_grid = np.zeros((10, 10), dtype=int)
    hunt = [(i, j) for i in range(10) for j in range(10) if (i + j) % 2 == 0]
    np.random.shuffle(hunt)
    targets = []                       # pila delle celle da verificare
    hits = turns = 0
    while hits < 17:
        if targets:
            i, j = targets.pop()
        else:
            i, j = hunt.pop()
        if shot_grid[i, j]:            # gia' esplorata: salta il turno
            continue
        shot_grid[i, j] = 2 if grid[i, j] else 1
        turns += 1
        if grid[i, j]:
            hits += 1
            for di, dj in ((1, 0), (-1, 0), (0, 1), (0, -1)):
                a, b = i + di, j + dj
                if 0 <= a < 10 and 0 <= b < 10 and shot_grid[a, b] == 0:
                    targets.append((a, b))
\end{lstlisting}

\begin{attenzione}[la pila può contenere celle già colpite]
Una cella può finire nella pila \texttt{targets} due volte (da due colpi a segno diversi),
oppure essere colpita nel frattempo. Il controllo \texttt{if shot\_grid[i, j]: continue}
prima di sparare è indispensabile: senza, si conterebbero turni per colpi mai sparati e
il risultato sarebbe gonfiato.
\end{attenzione}

\begin{extra}[si può fare meglio in modalità target]
L'implementazione sopra è ``la più semplice'': spara alle quattro celle adiacenti senza
ragionare. Un miglioramento classico è sfruttare l'\kw{orientamento}: dopo due colpi a
segno allineati si sa che la nave è orizzontale o verticale, quindi si possono scartare
le celle adiacenti nell'altra direzione. La ricerca bayesiana della prossima sezione
ottiene questo effetto \emph{automaticamente}, senza regole scritte a mano.
\end{extra}

\subsection{Strategia 3: ricerca bayesiana}

Qui si usa davvero la teoria della lezione. A ogni turno si costruisce una
\kw{distribuzione a posteriori} sulla posizione delle navi e si spara dove la probabilità
di colpire è massima, cioè dove è massima l'\kw{utilità attesa} (con utilità $1$ per il
colpo a segno e $0$ per il colpo a vuoto, l'utilità attesa di una cella \emph{è} la sua
probabilità di contenere una nave).

Il conteggio è semplice: per ogni nave rimasta e per ogni sua possibile posa
(orizzontale o verticale) \emph{compatibile} con quanto si sa, si aggiunge un voto a
ciascuna cella coperta. Una posa è compatibile se non copre nessuna cella risultata a
vuoto. Si distinguono due regimi:
\begin{itemize}
  \item nessun colpo a segno ``aperto'' (\kw{hunt}): si contano tutte le pose compatibili;
  \item ci sono colpi a segno (\kw{target}): si contano solo le pose che \emph{spiegano}
        almeno un colpo a segno, pesandole con $10^{(\text{colpi coperti})}$ per premiare
        le pose che ne spiegano di più.
\end{itemize}
In entrambi i casi i voti si assegnano solo alle celle ancora \kw{inesplorate}: sparare
dove si è già sparato non ha alcun valore.

\begin{lstlisting}
SIZES = [5, 4, 3, 3, 2]

def update_posterior(shot_map, sizes):
    post = conta_pose(shot_map, sizes, targeting=(shot_map == 2).any())
    if post.max() == 0:                     # nessuna posa spiega i colpi a segno
        post = conta_pose(shot_map, sizes, targeting=False)
    if post.max() == 0:                     # nessuna posa compatibile: cella libera
        post = (shot_map == 0).astype(float)
    return post

def conta_pose(shot_map, sizes, targeting):
    post = np.zeros((10, 10))
    for length in sizes:
        for horizontal in (True, False):
            righe = 10 if horizontal else 10 - length + 1
            colonne = 10 - length + 1 if horizontal else 10
            for i in range(righe):
                for j in range(colonne):
                    cells = [(i, j + k) if horizontal else (i + k, j)
                             for k in range(length)]
                    if any(shot_map[a, b] == 1 for a, b in cells):
                        continue                 # la posa copre un colpo a vuoto
                    n_hit = sum(shot_map[a, b] == 2 for a, b in cells)
                    if targeting and n_hit == 0:
                        continue                 # non spiega nessun colpo a segno
                    peso = 10 ** n_hit
                    for a, b in cells:
                        if shot_map[a, b] == 0:
                            post[a, b] += peso
    return post

def get_best_coordinates(posterior_map):
    return np.unravel_index(np.argmax(posterior_map), posterior_map.shape)

def bayesian_search_strategy(grid):
    shot_grid = np.zeros((10, 10), dtype=int)
    hits = turns = 0
    while hits < 17:
        i, j = get_best_coordinates(update_posterior(shot_grid, SIZES))
        shot_grid[i, j] = 2 if grid[i, j] else 1
        turns += 1
        if grid[i, j]:
            hits += 1
    return shot_grid, turns
\end{lstlisting}

\begin{osservazione}[dov'è il teorema di Bayes]
Il nome ``bayesiana'' è giustificato anche se nel codice non compare la formula. Il
prior è uniforme su tutte le disposizioni possibili; una cella risultata a vuoto ha
verosimiglianza $0$ per ogni disposizione che la occupa, quindi tutte quelle
disposizioni vengono eliminate dal posteriore. Contare le pose sopravvissute è
esattamente calcolare la probabilità a posteriori marginale di ogni cella, a meno della
costante di normalizzazione, che non serve perché si cerca solo l'\emph{argmax}.
\end{osservazione}

\begin{attenzione}[il posteriore può azzerarsi, e il programma si blocca]
Scritta nel modo ovvio (solo il conteggio, senza i due \texttt{if} iniziali) questa
strategia entra in \kw{ciclo infinito}. Succede così: una nave da 2 celle viene colpita
tutta e si trova circondata da celle già esplorate. I suoi due colpi a segno restano
``aperti'' perché non si tiene traccia degli affondamenti, quindi si è in modalità target;
ma l'unica posa che li spiega non tocca alcuna cella inesplorata, e quindi non assegna
voti a nessuno. Se tutti i colpi a segno aperti sono in questa condizione il posteriore è
\emph{tutto zero}, \texttt{argmax} restituisce sempre $(0,0)$, si spara all'infinito sulla
stessa cella e il contatore dei colpi a segno non avanza più.

Il rimedio è il ripiego in due passi: se il posteriore in modalità target è nullo si
ricalcola in modalità hunt, e se è nullo anche quello si spara su una qualsiasi cella
inesplorata. Vale la pena saperlo perché il sintomo (il notebook che non finisce mai)
non assomiglia per niente alla causa.
\end{attenzione}

\begin{attenzione}[il peso $10^{n\_hit}$ è un'euristica, non Bayes]
Il filtro sui colpi a vuoto è un calcolo bayesiano esatto; il peso $10^{n\_hit}$ no, è
una scorciatoia per dire ``le pose che spiegano più colpi sono molto più plausibili''.
Il calcolo rigoroso richiederebbe di enumerare le disposizioni \emph{congiunte} delle
cinque navi compatibili con tutti i dati, che sono troppe per essere contate a ogni
turno. Vale la pena saperlo, perché in sede d'esame la differenza tra ``esatto'' e
``approssimato'' è il tipo di cosa su cui si viene interrogati.
\end{attenzione}

\begin{extra}[che cosa manca rispetto alle regole vere]
Nelle regole, quando una nave affonda l'avversario deve dichiarare \emph{quale} nave era.
L'implementazione sopra non usa questa informazione: tiene sempre tutte e cinque le
lunghezze in \texttt{SIZES}. Tenere traccia delle navi affondate e rimuoverle dalla lista
(e marcare come ``chiusi'' i loro colpi a segno) rende il posteriore più preciso e fa
guadagnare qualche turno. È il primo miglioramento da provare.
\end{extra}

\subsection{Risultati misurati}

Le tre strategie sono state eseguite sulle stesse partite generate casualmente
(2000 partite per le prime due, 300 per la ricerca bayesiana, che è molto più pesante:
ricostruisce il posteriore da zero a ogni turno).

\begin{center}
\renewcommand{\arraystretch}{1.25}
\begin{tabular}{lrrr}
\toprule
 & Casuale & Hunt and target & Ricerca bayesiana \\
\midrule
Partite simulate            & $2000$ & $2000$ & $300$ \\
\midrule
Turni medi                  & $95{,}3$ & $59{,}8$ & $54{,}2$ \\
Mediana                     & $97$ & $61$ & $54$ \\
Minimo                      & $66$ & $30$ & $29$ \\
Massimo                     & $100$ & $80$ & $77$ \\
Partite chiuse entro 50 turni & $0\%$ & $17{,}6\%$ & $35{,}0\%$ \\
\bottomrule
\end{tabular}
\end{center}

Il salto grosso è il primo: passare dal caso a \en{hunt and target} fa risparmiare
circa 35 turni, cioè oltre un terzo della partita. Il secondo salto è molto più piccolo,
circa 6 turni in media, ma cambia la \emph{forma} della distribuzione: la quota di partite
chiuse entro 50 turni raddoppia (dal 17{,}6\% al 35\%). La ricerca bayesiana non è tanto
più veloce nel caso tipico, quanto più \kw{affidabile}: sbaglia meno spesso in grande.

\begin{attenzione}[questi 54 turni si possono migliorare parecchio]
Le implementazioni pubblicate di ricerca bayesiana per la battaglia navale arrivano
attorno ai 42--45 turni. La differenza viene quasi tutta da due semplificazioni fatte
sopra: non si tiene traccia delle navi affondate (quindi il posteriore continua a cercare
navi che non ci sono più) e il peso $10^{n\_hit}$ è un'euristica. Il numero riportato è
quello \emph{effettivamente misurato con questo codice}, non il migliore ottenibile:
vale la pena saperlo prima di citarlo.
\end{attenzione}

\subsection{Che cosa portarsi dietro}

\begin{itemize}
  \item Le tre strategie sono tre livelli di uso dell'informazione: nessuno (casuale),
        locale (hunt and target), globale (bayesiana).
  \item La ricerca bayesiana non ha regole scritte a mano: si limita a massimizzare
        l'utilità attesa rispetto a una distribuzione aggiornata. Il comportamento
        ``intelligente'' (sparare a scacchiera all'inizio, seguire l'orientamento della
        nave dopo due colpi) \kw{emerge} dal calcolo.
  \item È lo schema che si ritrova nelle vere applicazioni di \en{search theory}: i paper
        in \texttt{labs/02-battleship/papers/} raccontano come la stessa idea sia stata usata per
        ritrovare il relitto del volo Air France 447.
\end{itemize}

% ---------------------------------------------------------------------
\begin{riepilogo}
\begin{itemize}
  \item Le preferenze sono una relazione $\pref$ su $X$; se è completa e transitiva è una
        \emph{relazione di preferenza}, e da essa si derivano $\spref$ e $\indif$.
  \item \textbf{Teorema 1}: con $X$ finito, $\pref$ è rappresentabile da $u : X \to \R$
        se e solo se è completa e transitiva. L'utilità è \emph{ordinale}: invariante per
        trasformazioni strettamente crescenti, quindi i numeri non si interpretano.
  \item Con l'incertezza si passa alle \emph{lotterie}. Una lotteria composta si riduce a
        una semplice moltiplicando le probabilità lungo i rami.
  \item \textbf{Teorema VNM}: $\pref$ su $\lott$ ammette utilità attesa
        $U(L) = \sum_i p_i u(x_i)$ se e solo se vale ordine di preferenza, continuità e
        indipendenza. L'utilità è \emph{cardinale}: unica a meno di trasformazioni affini
        $aU + b$ con $a>0$.
  \item Il \emph{paradosso di Allais} mostra che le persone violano l'indipendenza: VNM è
        una teoria normativa, non descrittiva.
  \item Atteggiamento verso il rischio = curvatura di $u$: lineare $\to$ neutrale,
        concava $\to$ avverso, convessa $\to$ amante. Con $u(x) = \sqrt{x}$ la
        condivisione del rischio e l'assicurazione creano valore per entrambe le parti
        (premio $P \in [50,75]$ nell'esempio).
  \item \textbf{Lab}: la ricerca bayesiana sceglie a ogni turno la cella di massima utilità
        attesa, aggiornando il posteriore con i colpi a vuoto e a segno.
\end{itemize}
\textbf{Prossima lezione}: giochi non cooperativi in forma normale, \en{best response} ed
equilibrio di Nash.
\end{riepilogo}

% =====================================================================
%  Lezione 3 — Teoria dei giochi non cooperativi
%  Materiale: Lesson_3_Non_Cooperative_Game_Theory.pdf,
%             Lab03-Non_Cooperative_Game_Theory.pdf,
%             NonCooperativeGameTheory_part1_blank.ipynb,
%             Potential_Games_Label_Propagation_blank.ipynb
% =====================================================================
\chapter{Giochi non cooperativi}

\parteteoria

Nella lezione 1 i giochi erano combinatori: due giocatori, a turni, a somma zero, con
informazione perfetta. Qui si toglie quasi tutto. I giocatori possono essere $n$, muovono
\emph{simultaneamente}, ognuno ha la propria utilità (non necessariamente opposta a quella
degli altri) e le utilità sono quelle di VNM della lezione 2.

Il percorso della lezione:
\begin{enumerate}[label=\arabic*.]
  \item giochi in forma normale;
  \item strategie dominanti;
  \item strategie pure e \en{best response};
  \item equilibrio di Nash in strategie pure;
  \item strategie miste e \en{best response};
  \item equilibrio di Nash in strategie miste;
  \item equilibri misti nei giochi $2\times2$;
  \item giochi a somma zero e minimax.
\end{enumerate}

% ---------------------------------------------------------------------
\section{Giochi strategici e forma normale}

Nei \kw{giochi strategici}:
\begin{itemize}
  \item i giocatori sono in \kw{competizione} tra loro;
  \item ognuno è \kw{egoista} (\en{self-interested}): vuole massimizzare il proprio payoff;
  \item ognuno deve decidere quale azione intraprendere.
\end{itemize}

Ci sono due modi di rappresentarli:
\begin{itemize}
  \item \kw{forma normale} (o \emph{forma matriciale}): elenca il payoff di ciascun giocatore
        in funzione delle azioni di tutti, \emph{come se} si muovesse simultaneamente;
  \item \kw{forma estesa}: include l'ordine delle mosse e ciò che ogni giocatore sa quando
        muove. Si rappresenta come un albero (lo si è visto nella lezione 1). Negli scacchi
        il bianco muove e il nero risponde; nel poker conta anche che cosa un giocatore
        \emph{vede} quando punta.
\end{itemize}

\begin{definizione}[Gioco in forma normale]
Un gioco in forma normale è una tupla $\langle N, A, u\rangle$ dove:
\begin{itemize}
  \item $N = \{1,2,\dots,n\}$ è l'insieme dei \kw{giocatori};
  \item $A = \{A_1, A_2, \dots, A_n\}$ raccoglie le \kw{azioni} di tutti: $A_i$ è
        l'insieme delle azioni del giocatore $i$, e
        $a = (a_1,\dots,a_n) \in A_1 \times \dots \times A_n$ è un \kw{profilo di azioni};
  \item $u_i : A_1\times\dots\times A_n \to \R$ è la \kw{funzione di utilità} (payoff) del
        giocatore $i$, e $u = (u_1,\dots,u_n)$ è il profilo delle utilità.
\end{itemize}
\end{definizione}

\subsection{Giochi a due giocatori come matrici}

Con due giocatori si usa una tabella: il giocatore 1 sceglie la \kw{riga}, il giocatore 2
la \kw{colonna}; in ogni cella il primo numero è il payoff della riga, il secondo quello
della colonna.

\begin{center}
\renewcommand{\arraystretch}{1.35}
\begin{tabular}{r|c|c|}
\multicolumn{1}{r}{} & \multicolumn{1}{c}{A} & \multicolumn{1}{c}{B}\\ \cline{2-3}
C & \pay{2}{1} & \pay{3}{4} \\ \cline{2-3}
D & \pay{1}{3} & \pay{4}{4} \\ \cline{2-3}
\end{tabular}
\qquad\qquad
$A = \begin{pmatrix} 2 & 3 \\ 1 & 4 \end{pmatrix},\qquad
 B = \begin{pmatrix} 1 & 4 \\ 3 & 4 \end{pmatrix}$
\end{center}

I giocatori sono 1 (riga) e 2 (colonna), gli insiemi di azioni sono $\{C,D\}$ e
$\{A,B\}$, e per esempio $u_1(C,A) = 2$, $u_2(C,A) = 1$. Equivalentemente si usano due
matrici: $A$ con i payoff della riga e $B$ con quelli della colonna. È la rappresentazione
che si usa nel codice.

% ---------------------------------------------------------------------
\section{Strategie dominanti}

Una \kw{strategia pura} per il giocatore $i$ è la scelta di un'azione $a \in A_i$.
Un giocatore razionale non gioca una strategia che non gli dà il risultato migliore
\emph{qualunque cosa facciano gli altri}.

\begin{definizione}[Strategia strettamente e debolmente dominata]
Siano $s_i'$ e $s_i''$ due strategie del giocatore $i$. La strategia $s_i'$ è
\kw{strettamente dominata} da $s_i''$ se, per \emph{ogni} combinazione delle strategie degli
altri giocatori, il payoff di $i$ con $s_i'$ è strettamente minore che con $s_i''$:
\[
  u_i(s_1,\dots,s_{i-1}, s_i', s_{i+1},\dots,s_n) <
  u_i(s_1,\dots,s_{i-1}, s_i'', s_{i+1},\dots,s_n).
\]
La strategia $s_i'$ è \kw{debolmente dominata} da $s_i''$ se vale la stessa cosa con
$\le$ al posto di $<$, \emph{e almeno una disuguaglianza è stretta}.
\end{definizione}

\begin{attenzione}[la dominanza debole nelle slide]
Le slide definiscono la dominanza debole con il solo $\le$ ovunque. Così ogni strategia
dominerebbe debolmente sé stessa, e due strategie con payoff identici si dominerebbero a
vicenda. La definizione standard richiede $\le$ per ogni profilo degli avversari e $<$ per
almeno uno.
\end{attenzione}

\subsection{Eliminazione iterata delle strategie strettamente dominate}

Una strategia strettamente dominata non verrà mai giocata da un giocatore razionale, quindi
la si può \emph{eliminare} dal gioco. Nel gioco ridotto possono comparire nuove strategie
dominate: si ripete. Il procedimento si chiama \kw{IESDS} (\en{Iterated Elimination of
Strictly Dominated Strategies}).

\begin{esempio}[il dilemma del prigioniero]
\begin{center}
\renewcommand{\arraystretch}{1.35}
\begin{tabular}{r|c|c|}
\multicolumn{1}{r}{} & \multicolumn{1}{c}{C} & \multicolumn{1}{c}{D}\\ \cline{2-3}
C & \pay{-1}{-1} & \pay{-4}{0} \\ \cline{2-3}
D & \pay{0}{-4}  & \pay{-3}{-3} \\ \cline{2-3}
\end{tabular}
\end{center}
Per il giocatore 1, \emph{Cooperate} è strettamente dominata da \emph{Defect}
($-1 < 0$ e $-4 < -3$): la si elimina. Lo stesso per il giocatore 2. Sopravvive solo
$(D,D)$: \emph{Defect} è una \kw{strategia strettamente dominante}. L'esito è razionale ma
non socialmente ottimo: $(C,C)$ farebbe stare meglio entrambi.

Un gioco che si risolve completamente con IESDS si dice \kw{risolvibile per dominanza}
(\en{dominance solvable}).
\end{esempio}

\begin{osservazione}[numeri diversi dalla lezione 1]
Nella lezione 1 il payoff di chi coopera da solo era $-5$, qui è $-4$. Non cambia nulla:
quello che conta nel dilemma del prigioniero è l'\emph{ordine} dei payoff,
$a > b > c > d$ con $a$ = tradire da soli, $b$ = cooperare entrambi, $c$ = tradire
entrambi, $d$ = cooperare da soli (qui $0 > -1 > -3 > -4$).
\end{osservazione}

\begin{esempio}[un gioco senza strategie dominate]
\begin{center}
\renewcommand{\arraystretch}{1.35}
\begin{tabular}{r|c|c|c|}
\multicolumn{1}{r}{} & \multicolumn{1}{c}{T} & \multicolumn{1}{c}{V} & \multicolumn{1}{c}{C}\\ \cline{2-4}
L & \pay{0}{4} & \pay{4}{0} & \pay{5}{3} \\ \cline{2-4}
M & \pay{4}{0} & \pay{0}{4} & \pay{5}{3} \\ \cline{2-4}
R & \pay{3}{5} & \pay{3}{5} & \pay{8}{8} \\ \cline{2-4}
\end{tabular}
\end{center}
Nessuna strategia si può eliminare. Per il giocatore 1, L è la migliore se l'altro gioca
V, M se gioca T, R se gioca C; per il giocatore 2, T è la migliore contro L, V contro M,
C contro R. Ogni strategia è la migliore in almeno un caso, quindi nessuna è dominata:
\kw{non tutti i giochi si risolvono con IESDS}.
\end{esempio}

% ---------------------------------------------------------------------
\section{Best response ed equilibrio di Nash in strategie pure}

Si immagini che il giocatore $i$ sappia che cosa giocheranno tutti gli altri:
\[
  a_{-i} = \langle a_1, \dots, a_{i-1}, a_{i+1}, \dots, a_n\rangle .
\]

\begin{definizione}[Best response pura]
La \kw{best response pura} ad $a_{-i}$ è l'azione che massimizza il payoff di $i$:
\[
  a_i^* \in \BR(a_{-i}) \iff \forall a_i \in A_i,\ u_i(a_i^*, a_{-i}) \ge u_i(a_i, a_{-i}).
\]
\end{definizione}
$\BR(a_{-i})$ è un \emph{insieme}: le best response possono essere più di una, in caso di
pareggio.

\begin{definizione}[Equilibrio di Nash puro]
Un profilo di azioni è un \kw{equilibrio di Nash puro} (PNE) se ogni giocatore sta giocando
una best response:
\[
  a = \langle a_1,\dots,a_n\rangle \text{ è un NE} \iff \forall i,\ a_i \in \BR(a_{-i}).
\]
\end{definizione}

\begin{esempio}[trovare l'equilibrio puro]
Nel gioco $3\times 3$ precedente si segnano, per ogni colonna, il miglior payoff del
giocatore 1 e, per ogni riga, il miglior payoff del giocatore 2:
\begin{center}
\renewcommand{\arraystretch}{1.35}
\begin{tabular}{r|c|c|c|}
\multicolumn{1}{r}{} & \multicolumn{1}{c}{T} & \multicolumn{1}{c}{V} & \multicolumn{1}{c}{C}\\ \cline{2-4}
L & $0,\,\textcolor{blu}{\mathbf{4}}$ & $\textcolor{unical}{\mathbf{4}},\,0$ & $5,\,3$ \\ \cline{2-4}
M & $\textcolor{unical}{\mathbf{4}},\,0$ & $0,\,\textcolor{blu}{\mathbf{4}}$ & $5,\,3$ \\ \cline{2-4}
R & $3,\,5$ & $3,\,5$ & $\textcolor{unical}{\mathbf{8}},\,\textcolor{blu}{\mathbf{8}}$ \\ \cline{2-4}
\end{tabular}
\end{center}
L'unica cella in cui \emph{entrambi} i giocatori stanno rispondendo al meglio all'altro è
$(R, C)$: è l'unico equilibrio di Nash puro.
\end{esempio}

\begin{osservazione}[il metodo è generale]
``Segnare le best response e cercare le celle con entrambi i segni'' funziona per qualunque
gioco a due giocatori in forma matriciale, ed è esattamente quello che fa il codice del
laboratorio: per ogni cella controlla che la riga sia ottima nella sua colonna e la colonna
ottima nella sua riga.
\end{osservazione}

\subsection{IESDS ed equilibrio di Nash}
\begin{itemize}
  \item Se IESDS elimina tutti i profili tranne uno, quel profilo è anche l'\kw{unico}
        equilibrio di Nash puro.
  \item Il viceversa non vale: alcune strategie sopravvivono a IESDS senza far parte di un
        equilibrio di Nash (nel gioco $3\times3$ sopravvivono tutte, ma l'equilibrio è uno).
\end{itemize}

\begin{extra}[perché le strategie eliminate non stanno in un equilibrio]
Se $s_i'$ è strettamente dominata da $s_i''$, allora contro \emph{qualunque} $a_{-i}$ vale
$u_i(s_i'', a_{-i}) > u_i(s_i', a_{-i})$, quindi $s_i'$ non è mai una best response e non
può comparire in un equilibrio di Nash. Applicando l'argomento a ogni passo
dell'eliminazione, nessun equilibrio di Nash viene mai eliminato. Se alla fine resta un
solo profilo, e un equilibrio esiste, deve essere quello.
\end{extra}

\subsection{Come leggere l'equilibrio di Nash}

Il concetto di equilibrio di Nash non dice \emph{come} i giocatori lo raggiungano, ma la
logica è questa:
\begin{itemize}
  \item se i giocatori non sono in equilibrio, almeno uno ha l'opportunità di deviare, e lo
        farà;
  \item se sono in equilibrio, nessuno ha interesse a cambiare strategia;
  \item questa logica ragiona \kw{un giocatore alla volta}: non considera deviazioni
        congiunte di più giocatori.
\end{itemize}
Per questo gli equilibri di Nash si dicono \kw{stabili} e \kw{auto-rinforzanti}.

\begin{esempio}[la battaglia dei sessi: due equilibri]
\begin{center}
\renewcommand{\arraystretch}{1.35}
\begin{tabular}{r|c|c|}
\multicolumn{1}{r}{} & \multicolumn{1}{c}{A} & \multicolumn{1}{c}{R}\\ \cline{2-3}
A & \pay{2}{1} & \pay{0}{0} \\ \cline{2-3}
R & \pay{0}{0} & \pay{1}{2} \\ \cline{2-3}
\end{tabular}
\end{center}
Ci sono due equilibri puri, $(A,A)$ e $(R,R)$, e ciascuno è sbilanciato a favore di uno dei
due giocatori. Un'opzione migliore per entrambi sarebbe \emph{accordarsi} e giocare
$(A,A)$ e $(R,R)$ con probabilità $0{,}5$ ciascuno: payoff atteso $1{,}5$ a testa.
\end{esempio}

\begin{attenzione}[l'accordo non è una strategia mista]
Giocare ``$(A,A)$ metà delle volte e $(R,R)$ l'altra metà'' richiede che i due giocatori
\emph{si coordinino}, per esempio lanciando insieme una moneta. Se ciascuno randomizzasse
per conto suo al 50\%, metà delle volte si ritroverebbero in $(A,R)$ o $(R,A)$ con payoff
nullo. Un accordo di questo tipo non è un equilibrio di Nash ma un \kw{equilibrio
correlato}: è l'argomento della lezione 4.
\end{attenzione}

\begin{esempio}[i calci di rigore: nessun equilibrio puro]
\begin{center}
\renewcommand{\arraystretch}{1.35}
\begin{tabular}{r|c|c|}
\multicolumn{1}{r}{} & \multicolumn{1}{c}{L} & \multicolumn{1}{c}{R}\\ \cline{2-3}
L & \pay{1}{-1} & \pay{-1}{1} \\ \cline{2-3}
R & \pay{-1}{1} & \pay{1}{-1} \\ \cline{2-3}
\end{tabular}
\end{center}
Il portiere (riga) vince se indovina il lato, il tiratore (colonna) se lo sbaglia. Da
qualunque cella uno dei due ha interesse a cambiare: non esiste un equilibrio puro. Si
``risolve'' solo randomizzando tra L e R con probabilità $0{,}5$.
\end{esempio}

\begin{attenzione}[``il miglior esito per entrambi'']
Le slide dicono che randomizzare al 50\% dà ``il miglior esito per entrambi''. In un gioco a
somma zero questa frase non ha senso: ciò che uno guadagna l'altro lo perde, quindi non
esiste un esito migliore per entrambi. Il senso corretto è che 50/50 è l'unico profilo
\emph{stabile}: ciascuno si garantisce payoff atteso 0 e nessuno dei due può fare meglio
cambiando da solo. È la strategia minimax, che si ritrova alla fine della lezione.
\end{attenzione}

\subsection{Equilibri con utilità analitiche: la tragedia dei beni comuni}

Le azioni non devono per forza essere un insieme finito. Due agricoltori pompano acqua
dalla stessa falda:
\begin{itemize}
  \item $x_1, x_2 \ge 0$ sono i livelli di estrazione e $X = x_1 + x_2$ il totale;
  \item il ricavo per unità d'acqua è $A - BX$ con $A, B > 0$: vale $A$ quando non si è
        estratto nulla e cala di $B$ per ogni unità estratta in totale;
  \item l'utilità dell'agricoltore $i$ è $u(x_i, x_j) = x_i\bigl(A - B(x_i + x_j)\bigr)$.
\end{itemize}

\paragraph{Best response.}
L'agricoltore $i$ sceglie $x_i$ prendendo $x_j$ come dato. Condizione del primo ordine:
\[
  \frac{\partial u}{\partial x_i}
  = \frac{\partial}{\partial x_i}\bigl(Ax_i - Bx_i^2 - Bx_ix_j\bigr)
  = A - 2Bx_i - Bx_j = 0
  \quad\Longrightarrow\quad
  \BR_i(x_j) = x_i = \frac{A - Bx_j}{2B}.
\]

\paragraph{Equilibrio di Nash.}
Il gioco è simmetrico, quindi si cerca $x_i^* = x_j^* = x^*$:
\[
  x^* = \frac{A - Bx^*}{2B} \ \Rightarrow\ 2Bx^* = A - Bx^* \ \Rightarrow\
  x^* = \frac{A}{3B},
  \qquad X^{\text{Nash}} = 2x^* = \frac{2A}{3B}.
\]

\paragraph{Ottimo sociale.}
Un pianificatore massimizza il benessere totale $W(X) = X(A - BX)$:
\[
  \frac{dW}{dX} = A - 2BX = 0 \ \Rightarrow\ X^{\text{Opt}} = \frac{A}{2B},
  \qquad x^{\text{plan}} = \frac{A}{4B} \text{ a testa}.
\]

\begin{center}
\renewcommand{\arraystretch}{1.4}
\begin{tabular}{lccc}
\toprule
 & Estrazione totale & Benessere totale & Utilità a testa \\
\midrule
Equilibrio di Nash & $\dfrac{2A}{3B}$ & $\dfrac{2A^2}{9B} \approx 0{,}222\,\dfrac{A^2}{B}$ & $\dfrac{A^2}{9B}$ \\[6pt]
Ottimo sociale     & $\dfrac{A}{2B}$  & $\dfrac{A^2}{4B} = 0{,}25\,\dfrac{A^2}{B}$ & $\dfrac{A^2}{8B}$ \\
\bottomrule
\end{tabular}
\end{center}

In competizione gli agricoltori estraggono \emph{più} acqua del necessario e ottengono
\emph{meno} utilità ciascuno. È la \kw{tragedia dei beni comuni}: i giocatori sono
incentivati a sfruttare una risorsa comune a danno di tutti. Il meccanismo è lo stesso del
dilemma del prigioniero, con un continuo di azioni al posto di due.

\begin{extra}[i passaggi che le slide saltano]
\begin{itemize}
  \item La condizione del primo ordine individua un \emph{massimo} perché
        $\partial^2 u/\partial x_i^2 = -2B < 0$: l'utilità è concava in $x_i$.
  \item La best response vale solo finché $A - Bx_j \ge 0$; se l'altro estrae più di
        $A/B$ la risposta ottima è $x_i = 0$. All'equilibrio il vincolo non è attivo.
  \item L'utilità a testa all'equilibrio è $x^*(A - BX^{\text{Nash}}) =
        \tfrac{A}{3B}\cdot\tfrac{A}{3} = \tfrac{A^2}{9B}$; con il piano è
        $\tfrac{A}{4B}\cdot\tfrac{A}{2} = \tfrac{A^2}{8B}$, più alta.
  \item Il piano non è un equilibrio: se l'altro estrae $A/4B$, la best response è
        $\BR(A/4B) = \tfrac{A - A/4}{2B} = \tfrac{3A}{8B} > \tfrac{A}{4B}$.
        Ognuno ha interesse a estrarre di più.
\end{itemize}
\end{extra}

\begin{attenzione}[``ottimo di Pareto'' nel titolo delle slide]
Le slide confrontano l'equilibrio con l'``ottimo di Pareto'', ma quello che calcolano è
l'ottimo \emph{utilitaristico}, cioè il massimo della somma delle utilità. L'ottimo
utilitaristico è Pareto-ottimo, ma non è l'unico esito Pareto-ottimo: per esempio
un agricoltore che estrae $A/2B$ da solo, con l'altro a zero, è anch'esso Pareto-ottimo.
\end{attenzione}

% ---------------------------------------------------------------------
\section{Strategie miste}

\begin{definizione}[Strategia mista]
Una \kw{strategia} $\sigma_i$ del giocatore $i$ è una distribuzione di probabilità su $A_i$.
\begin{itemize}
  \item È \kw{pura} se una sola azione ha probabilità 1.
  \item È \kw{mista} se più azioni hanno probabilità positiva. Le azioni con probabilità
        positiva formano il \kw{supporto} della strategia, $\supp(\sigma_i)$.
\end{itemize}
$S_i$ è l'insieme delle strategie di $i$, $S = S_1 \times \dots \times S_n$ quello dei
profili, e $\sigma = (\sigma_1,\dots,\sigma_n) \in S$ è un profilo di strategie miste.
\end{definizione}

Con le strategie miste le utilità non si leggono più nella tabella: si usa l'\kw{utilità
attesa} della lezione 2. La probabilità di un profilo di azioni è il prodotto delle
probabilità con cui ogni giocatore sceglie la propria azione (i giocatori randomizzano in
modo indipendente):
\[
  u_i(\sigma) = \sum_{a \in A} u_i(a)\,\Pr(a \mid \sigma),
  \qquad
  \Pr(a \mid \sigma) = \prod_{j \in N} \sigma_j(a_j).
\]
In forma matriciale, con $\sigma_r$ e $\sigma_c$ vettori riga:
$u_r = \sigma_r A \sigma_c^T$ e $u_c = \sigma_r B \sigma_c^T$.

\begin{esempio}[matching pennies]
\begin{center}
\renewcommand{\arraystretch}{1.35}
\begin{tabular}{r|c|c|}
\multicolumn{1}{r}{} & \multicolumn{1}{c}{H} & \multicolumn{1}{c}{T}\\ \cline{2-3}
H & \pay{1}{-1} & \pay{-1}{1} \\ \cline{2-3}
T & \pay{-1}{1} & \pay{1}{-1} \\ \cline{2-3}
\end{tabular}
\end{center}
Con $\sigma_r = \sigma_c = (\tfrac12, \tfrac12)$ ogni cella ha probabilità $0{,}25$:
\[
  u_r = 1\cdot 0{,}25 - 1\cdot 0{,}25 - 1\cdot 0{,}25 + 1\cdot 0{,}25 = 0 .
\]
Con $\sigma_r = \sigma_c = (\tfrac45, \tfrac15)$:
\[
  u_r = 1\cdot(0{,}8\cdot0{,}8) - 1\cdot(0{,}8\cdot0{,}2) - 1\cdot(0{,}2\cdot0{,}8)
        + 1\cdot(0{,}2\cdot0{,}2) = 0{,}64 - 0{,}16 - 0{,}16 + 0{,}04 = 0{,}36 .
\]
\end{esempio}

\begin{attenzione}[perché servono le utilità di VNM]
Calcolare il valore atteso dei payoff ha senso solo se i payoff sono utilità
\emph{cardinali}, cioè di VNM (lezione 2). Con utilità solo ordinali un equilibrio misto
non è nemmeno definito: trasformando i numeri con una funzione crescente non lineare,
l'equilibrio misto cambierebbe, mentre gli equilibri puri resterebbero gli stessi.
\end{attenzione}

\section{Best response ed equilibrio di Nash misti}

\begin{definizione}[Best response mista ed equilibrio di Nash misto]
La \kw{best response mista} a $\sigma_{-i}$ è la strategia che massimizza il payoff
atteso di $i$:
\[
  \sigma_i^* \in \BR(\sigma_{-i}) \iff
  \forall \sigma_i \in S_i,\ u_i(\sigma_i^*, \sigma_{-i}) \ge u_i(\sigma_i, \sigma_{-i}).
\]
Un profilo $\sigma = \langle\sigma_1,\dots,\sigma_n\rangle$ è un \kw{equilibrio di Nash
misto} (MNE) se $\forall i,\ \sigma_i \in \BR(\sigma_{-i})$.
\end{definizione}

\begin{attenzione}[un refuso nella definizione]
Le slide scrivono ``$\forall \sigma_i \in A_i$''. Le strategie miste stanno in $S_i$
(l'insieme delle distribuzioni su $A_i$), non in $A_i$.
\end{attenzione}

\begin{teorema}[Nash, 1950]
Ogni gioco finito (numero finito di giocatori e di azioni) ha almeno un equilibrio di Nash
misto.
\end{teorema}
La dimostrazione, tramite il teorema del punto fisso di Brouwer, è nella lezione 4.

\subsection{Condizione generale per una best response}

\begin{teorema}[Condizione del supporto]
In un gioco a due giocatori $(A,B) \in \R^{m\times n}$, la strategia $\sigma_r$ della riga è
una best response a $\sigma_c$ se e solo se
\[
  \sigma_r(i) > 0 \ \Longrightarrow\ (A\sigma_c^T)_i = u = \max_k (A\sigma_c^T)_k
  \qquad \forall i \in A_1 .
\]
\end{teorema}
In parole: \kw{una strategia mista è una best response se e solo se mette probabilità
solo su azioni pure che sono a loro volta best response}.

\paragraph{Dimostrazione.}
$(A\sigma_c^T)_i$ è l'utilità che la riga ottiene giocando la sua $i$-esima azione. Sia
$u = \max_k (A\sigma_c^T)_k$. L'utilità di $\sigma_r$ è
\begin{align*}
  \sigma_r A \sigma_c^T
  &= \sum_{i=1}^m \sigma_r(i)\,(A\sigma_c^T)_i
   = \sum_{i=1}^m \sigma_r(i)\bigl[u - u + (A\sigma_c^T)_i\bigr]\\
  &= \sum_{i=1}^m \sigma_r(i)\,u - \sum_{i=1}^m \sigma_r(i)\bigl(u - (A\sigma_c^T)_i\bigr)
   = u - \sum_{i=1}^m \sigma_r(i)\bigl(u - (A\sigma_c^T)_i\bigr).
\end{align*}
Ogni termine $u - (A\sigma_c^T)_i$ è $\ge 0$, quindi l'utilità massima possibile è $u$, e
si raggiunge se e solo se la somma è nulla, cioè se ogni azione con $\sigma_r(i) > 0$ ha
$(A\sigma_c^T)_i = u$. $\square$

\begin{osservazione}[due conseguenze pratiche]
\begin{itemize}
  \item Per verificare se una strategia è una best response non serve confrontarla con
        tutte le infinite strategie miste: basta calcolare il vettore $A\sigma_c^T$
        (le utilità delle azioni pure) e guardare il supporto. È l'algoritmo del laboratorio.
  \item All'equilibrio misto ogni giocatore è \emph{indifferente} tra le azioni del
        proprio supporto. È questa indifferenza che permette di calcolare gli equilibri nei
        giochi $2\times2$.
\end{itemize}
\end{osservazione}

\begin{esempio}[best response in sasso-carta-forbici]
La colonna non gioca mai forbici e gioca sasso o carta con probabilità $0{,}5$:
\[
  A = \begin{pmatrix} 0 & -1 & 1 \\ 1 & 0 & -1 \\ -1 & 1 & 0 \end{pmatrix},
  \qquad \sigma_c = \bigl(\tfrac12, \tfrac12, 0\bigr),
  \qquad A\sigma_c^T = \begin{pmatrix} -0{,}5 \\ 0{,}5 \\ 0 \end{pmatrix}
  \ \Longrightarrow\ \sigma_r^* = (0, 1, 0).
\]
La seconda riga (carta) ha l'utilità massima: l'unica best response è giocare sempre carta.
\end{esempio}

\begin{esempio}[best response nella battaglia dei sessi]
Con $A = \begin{pmatrix}2&0\\0&1\end{pmatrix}$, $B = \begin{pmatrix}1&0\\0&2\end{pmatrix}$
e $\sigma_c = (\tfrac13, \tfrac23)$:
\[
  A\sigma_c^T = \bigl(\tfrac23, \tfrac23\bigr)^T .
\]
Le due azioni della riga danno la stessa utilità, quindi \emph{qualunque} strategia della
riga è una best response: $(\tfrac23,\tfrac13)$, $(\tfrac12,\tfrac12)$, $(1,0)$, \ldots
\begin{itemize}
  \item Con $\sigma_r = (\tfrac23, \tfrac13)$: $\sigma_r B = (\tfrac23, \tfrac23)$, quindi anche
        $\sigma_c$ è una best response a $\sigma_r$. Le due strategie sono una best response
        l'una all'altra: $(\sigma_r, \sigma_c)$ è un equilibrio di Nash.
  \item Con $\sigma_r = (\tfrac12, \tfrac12)$: $\sigma_r B = (\tfrac12, 1)$. La colonna
        dovrebbe giocare solo la seconda azione, ma $\sigma_c$ mette probabilità $\tfrac13$
        sulla prima: $\sigma_c \notin \BR(\sigma_r)$ e il profilo non è un equilibrio,
        anche se $\sigma_r \in \BR(\sigma_c)$.
\end{itemize}
\end{esempio}

% ---------------------------------------------------------------------
\section{Equilibri misti nei giochi \texorpdfstring{$2\times2$}{2x2}}

Nella battaglia dei sessi
\begin{center}
\renewcommand{\arraystretch}{1.35}
\begin{tabular}{r|c|c|}
\multicolumn{1}{r}{} & \multicolumn{1}{c}{A} & \multicolumn{1}{c}{C}\\ \cline{2-3}
A & \pay{2}{1} & \pay{0}{0} \\ \cline{2-3}
C & \pay{0}{0} & \pay{1}{2} \\ \cline{2-3}
\end{tabular}
\end{center}
si cerca un equilibrio in cui entrambi randomizzano. Per la condizione del supporto, ognuno
dev'essere indifferente tra le proprie due azioni.

\paragraph{La colonna rende indifferente la riga.}
La colonna gioca A con probabilità $p$. La riga è indifferente se
\[
  u_r(A) = u_r(C) \ \Longleftrightarrow\ 2p + 0(1-p) = 0\cdot p + 1(1-p)
  \ \Longrightarrow\ p = \tfrac13 .
\]

\paragraph{La riga rende indifferente la colonna.}
La riga gioca A con probabilità $q$. La colonna è indifferente se
\[
  u_c(A) = u_c(C) \ \Longleftrightarrow\ q\cdot 1 + 0(1-q) = 0\cdot q + 2(1-q)
  \ \Longrightarrow\ q = \tfrac23 .
\]

L'equilibrio misto è quindi $\sigma_r = (\tfrac23, \tfrac13)$, $\sigma_c = (\tfrac13, \tfrac23)$.

\begin{osservazione}[il trucco controintuitivo]
La probabilità con cui gioca un giocatore è determinata dai payoff dell'\emph{altro}. La
riga sceglie $q$ non per massimizzare il proprio payoff (all'equilibrio le sue due azioni le
rendono uguale) ma per rendere l'avversario indifferente: se non lo fosse, l'avversario
smetterebbe di randomizzare e l'equilibrio crollerebbe.
\end{osservazione}

\begin{extra}[quanto vale l'equilibrio misto]
All'equilibrio misto della battaglia dei sessi ciascuno ottiene payoff atteso $\tfrac23$:
meno di $1$, il peggior payoff dei due equilibri puri. I giocatori si mancano con
probabilità $\tfrac23\cdot\tfrac23 + \tfrac13\cdot\tfrac13 = \tfrac59$. L'equilibrio misto è
``equo'' (simmetrico nei payoff) ma inefficiente: è il motivo per cui si cercano
alternative come gli equilibri correlati (lezione 4).
\end{extra}

% ---------------------------------------------------------------------
\section{Giochi a somma zero e minimax}

\begin{definizione}[Gioco a somma zero]
Un gioco a due giocatori $(A,B)$ è \kw{a somma zero} se e solo se $A = -B$. Basta quindi una
sola matrice $A$: la riga massimizza $\sigma_r A\sigma_c^T$, la colonna lo minimizza.
\end{definizione}

\paragraph{Il punto di vista della colonna.}
Data $A \in \R^{m\times n}$ e la strategia $\sigma_c$ della colonna, la best response della
riga ottiene
\[
  \max_{\sigma_r \in S_r} \sigma_r A\sigma_c^T = \max_i (A\sigma_c^T)_i
  = \min\{ v \in \R \mid A\sigma_c^T \le \uno v \}.
\]
La prima uguaglianza è la condizione del supporto (il massimo si raggiunge su un'azione
pura); la seconda dice che il massimo di un vettore è il più piccolo $v$ che lo maggiora
componente per componente. Con la sua scelta di $\sigma_c$ la colonna \emph{fissa} questo
tetto $v$ al guadagno della riga; essendo il gioco a somma zero, vuole renderlo il più basso
possibile.

\begin{definizione}[I due programmi lineari del minimax]
\[
\begin{array}{lll}
\text{Colonna:} & \displaystyle\min_{v,\,\sigma_c} v
  & \text{s.t. } A\sigma_c^T \le \uno v,\ \ \sigma_c \in S_c\\[10pt]
\text{Riga:} & \displaystyle\max_{u,\,\sigma_r} u
  & \text{s.t. } \sigma_r A \ge u\uno^T,\ \ \sigma_r \in S_r
\end{array}
\]
\end{definizione}
La riga cerca la strategia che massimizza il payoff \emph{garantito} $u$ contro qualunque
colonna (\kw{maximin}); la colonna quella che minimizza il payoff massimo concesso $v$
(\kw{minimax}).

\begin{teorema}[Minimax, von Neumann 1928]
In un gioco a somma zero $A \in \R^{m\times n}$ con insiemi di strategie miste $S_r, S_c$
esistono $\sigma_r \in S_r$, $\sigma_c \in S_c$ e $v \in \R$ tali che
\[
  A\sigma_c^T \le \uno v, \qquad \sigma_r A \ge v\uno^T .
\]
Allora $v$ è l'unico \kw{valore del gioco} e $\sigma_r, \sigma_c$ sono strategie ottime dei
giocatori.
\end{teorema}

In altre parole $\max\min = \min\max$: il payoff che la riga si può garantire coincide con
quello che la colonna può impedirle di superare.

\paragraph{Dimostrazione (traccia delle slide).}
\begin{enumerate}
  \item \emph{Debolezza} ($u \le v$). Per soluzioni ammissibili qualsiasi dei due programmi,
        \[
          u = u\,\uno^T\sigma_c^T \le \sigma_r A\sigma_c^T \le \sigma_r\uno v = v .
        \]
        Quindi gli obiettivi sono limitati e si possono restringere agli insiemi compatti
        $S_c\times[u,v]$ e $S_r\times[u,v]$; per continuità di $\sigma_r A\sigma_c^T$ gli
        ottimi $u$ e $v$ esistono.
  \item \emph{Uguaglianza} ($u = v$), per induzione su $m+n$. Con $m+n = 2$ (matrice
        $1\times1$) è ovvio. Siano $\sigma_c, v$ e $\sigma_r, u$ ottimi. Se
        $A\sigma_c^T = \uno v$ e $\sigma_r A = u\uno^T$, allora
        $u = \sigma_r A\sigma_c^T = v$. Altrimenti una disuguaglianza è stretta, per esempio
        la riga $k$: $(A\sigma_c^T)_k < v$.
  \item Si cancella la riga $k$ ottenendo $\bar A$ ($(m-1)\times n$), con ottimi
        $\bar u, \bar v$. Per induzione $\bar u = \bar v$. Togliere una riga può solo
        peggiorare la riga e aiutare la colonna: $\bar u \le u$ e $\bar v \le v$.
  \item Se fosse $\bar v < v$, la strategia $\sigma_c' = (1-\epsilon)\sigma_c +
        \epsilon\bar\sigma_c$ darebbe $(A\sigma_c'^T)_i < v$ su tutte le righe diverse da
        $k$ (per ogni $\epsilon \in (0,1]$) e anche sulla riga $k$ (per $\epsilon$ piccolo,
        perché lì la disuguaglianza era già stretta): contraddice la minimalità di $v$.
        Quindi $\bar v = v$ e
        \[
          \bar u \le u \le v = \bar v = \bar u \ \Longrightarrow\ u = v . \qquad\square
        \]
\end{enumerate}

\begin{attenzione}[due dettagli della dimostrazione]
\begin{itemize}
  \item Al passo 4 le slide scelgono $\epsilon$ solo in modo che
        $v(1-\epsilon) + \bar v\epsilon < v$; per la riga cancellata $k$ serve anche
        $\epsilon$ \emph{abbastanza piccolo}, perché $(A\bar\sigma_c^T)_k$ può essere
        grande. Esiste sempre, dato che $(A\sigma_c^T)_k < v$ strettamente.
  \item Le slide trattano solo il caso in cui la disuguaglianza stretta è in
        $A\sigma_c^T \le \uno v$. Se tutte quelle sono uguaglianze ma una di
        $\sigma_r A \ge u\uno^T$ è stretta, si ragiona allo stesso modo cancellando una
        \emph{colonna}.
\end{itemize}
\end{attenzione}

\begin{extra}[il legame con la lezione 1]
Minimax per i giochi combinatori (lezione 1) è il caso particolare in cui l'informazione
è perfetta e i giocatori muovono a turno: lì il max e il min si scambiano senza bisogno di
strategie miste, e il valore si calcola risalendo l'albero. Nei giochi simultanei come
sasso-carta-forbici, invece, senza randomizzare il max-min e il min-max \emph{non}
coincidono (in pure: $-1$ contro $+1$); le strategie miste sono proprio ciò che serve a
farli coincidere.
\end{extra}

\subsection{Esempio: sasso-carta-forbici come programma lineare}

Con $A = \begin{pmatrix} 0&-1&1\\ 1&0&-1\\ -1&1&0 \end{pmatrix}$ e $x = \sigma_r$, il
programma della riga $\max u$ s.t.\ $xA \ge \uno u$, $x \in S_r$ diventa:
\[
\begin{aligned}
  \max_{u,x}\ & u\\
  \text{s.t. }\ & 0x_1 + 1x_2 - 1x_3 \ge u\\
  & -1x_1 + 0x_2 + 1x_3 \ge u\\
  & 1x_1 - 1x_2 + 0x_3 \ge u\\
  & x_1 + x_2 + x_3 = 1
\end{aligned}
\]
Ogni vincolo dice ``contro la colonna $j$ la riga ottiene almeno $u$''.

\paragraph{Forma standard.}
I risolutori (come \texttt{linprog} di SciPy) vogliono
$\min c\,x$ s.t.\ $M_{ub}\,x \le b_{ub}$, $M_{eq}\,x = b_{eq}$. Si rinomina $u = x_4$, si
porta tutto a sinistra e si cambia segno:
\[
\begin{aligned}
  \min_x\ & -x_4\\
  \text{s.t. }\ & 0x_1 - 1x_2 + 1x_3 + 1x_4 \le 0\\
  & 1x_1 + 0x_2 - 1x_3 + 1x_4 \le 0\\
  & -1x_1 + 1x_2 + 0x_3 + 1x_4 \le 0\\
  & 1x_1 + 1x_2 + 1x_3 + 0x_4 = 1
\end{aligned}
\qquad
\begin{aligned}
  M_{ub} &= \begin{pmatrix} 0&-1&1&1\\ 1&0&-1&1\\ -1&1&0&1 \end{pmatrix},\
  b_{ub} = \begin{pmatrix}0\\0\\0\end{pmatrix},\\
  M_{eq} &= (1,1,1,0),\ b_{eq} = 1,\ c = (0,0,0,-1).
\end{aligned}
\]
Per una matrice qualsiasi $A \in \R^{m\times n}$:
\[
  M_{ub} = \bigl[\,-A^T \mid \uno\,\bigr] \in \R^{n\times(m+1)},\quad
  b_{ub} = \mathbf{0}_n,\quad
  c = (\underbrace{0,\dots,0}_{m}, -1),\quad
  M_{eq} = (\underbrace{1,\dots,1}_{m}, 0),\quad b_{eq} = 1 .
\]

\begin{attenzione}[il valore del gioco può essere negativo]
Le slide scrivono la forma standard con il vincolo $x \ge 0$ su \emph{tutte} le variabili.
Ma $x_4 = u$ è il valore del gioco, che può essere negativo. Il notebook lo dice
correttamente ($x_i \ge 0$ solo per $i \le m$), ed è un dettaglio che conta: \texttt{linprog}
impone di default $x_i \ge 0$ a \emph{tutte} le variabili. Se non si dichiara $u$ libera
(\texttt{bounds=[(0, None)]*m + [(None, None)]}), su un gioco di valore negativo il
programma diventa \kw{impossibile} e il risolutore non restituisce nessuna soluzione. Nel
laboratorio lo si verifica su un esempio; tutti i test del notebook hanno valore
$\ge 0$ e non se ne accorgono.
\end{attenzione}

% =====================================================================
\partelab{Best response, equilibri e giochi potenziali}
% =====================================================================
\fonte{NonCooperativeGameTheory\_part1\_blank.ipynb, Potential\_Games\_Label\_Propagation\_blank.ipynb}

\textbf{Obiettivo}: implementare gli strumenti della lezione per giochi a due giocatori
(verifica di best response ed equilibri, minimax per i giochi a somma zero, equilibri misti
$2\times2$) e poi un'applicazione dei \kw{giochi potenziali}: trovare le comunità di un
grafo facendo giocare i nodi. Tutto il codice sotto è stato eseguito: i risultati attesi
scritti nel notebook sono riprodotti esattamente.

\begin{osservazione}[il notebook della lezione 4]
Il notebook del laboratorio 4 riprende integralmente questa prima parte e la estende. Le
celle di quel notebook contengono già gli output della soluzione del docente, ed è con
quelli che sono stati confrontati i risultati di questa sezione.
\end{osservazione}

\subsection{Best response ed equilibrio di Nash}

Le strategie sono vettori di probabilità (\texttt{[0, 1]} = seconda azione pura). La stessa
funzione serve per strategie pure e miste, grazie alla condizione del supporto: si calcolano
le utilità delle azioni pure contro l'avversario e si controlla che il supporto contenga solo
azioni di utilità massima.

\begin{lstlisting}
import numpy as np

def is_best_response(P1, P2, s1, s2):
    s1, s2 = np.asarray(s1, float), np.asarray(s2, float)
    for s in (s1, s2):
        assert np.all(s >= 0) and np.isclose(s.sum(), 1), "strategia non valida"
    u_row = P1 @ s2               # utilita' di ogni riga contro s2
    u_col = s1 @ P2               # utilita' di ogni colonna contro s1
    br_row = np.all(u_row[s1 > 0] >= u_row.max() - 1e-9)
    br_col = np.all(u_col[s2 > 0] >= u_col.max() - 1e-9)
    return bool(br_row), bool(br_col)

def is_nash_equilibrium(P1, P2, s1, s2):
    return all(is_best_response(P1, P2, s1, s2))
\end{lstlisting}

\begin{center}
\renewcommand{\arraystretch}{1.2}
\small
\begin{tabular}{llll}
\toprule
Gioco & Profilo & Atteso & Ottenuto \\
\midrule
Dilemma del prigioniero & (D, C) & (True, False) & (True, False) \\
Dilemma del prigioniero & (D, D) & NE & True \\
Gioco $C/D \times A/B$  & (C, A) & (True, False) & (True, False) \\
Gioco $C/D \times A/B$  & (D, B) & NE & True \\
Battaglia dei sessi     & $(\frac23,\frac13), (\frac13,\frac23)$ & (True, True) & (True, True) \\
Battaglia dei sessi     & $(\frac23,\frac13), (\frac12,\frac12)$ & (False, True) & (False, True) \\
Battaglia dei sessi     & (A, A) & NE & True \\
Battaglia dei sessi     & $(0{,}8;\ 0{,}8)$ & errore & AssertionError \\
\bottomrule
\end{tabular}
\end{center}

\begin{attenzione}[confronti esatti con le probabilità]
Le strategie del tipo $(\tfrac23, \tfrac13)$ non sono rappresentabili esattamente in
floating point, e le utilità che devono essere ``uguali'' possono differire di
$10^{-16}$. Un test come \texttt{u\_row[i] == u\_row.max()} può allora dare risultati
casuali. Per questo il codice confronta con una tolleranza (\texttt{- 1e-9}) e verifica la
somma con \texttt{np.isclose}. Lo stesso vale per qualunque verifica di equilibrio nel
progetto.
\end{attenzione}

\subsection{Minimax nei giochi a somma zero}

\begin{lstlisting}
import scipy.optimize

def minimax_LP(A):
    A = np.asarray(A, float)
    m, n = A.shape
    c = np.zeros(m + 1); c[-1] = -1                  # max u = min -u
    M_ub = np.hstack([-A.T, np.ones((n, 1))])        # u - (xA)_j <= 0
    b_ub = np.zeros(n)
    M_eq = np.hstack([np.ones((1, m)), np.zeros((1, 1))])
    b_eq = [1]
    bounds = [(0, None)] * m + [(None, None)]        # u libera!
    res = scipy.optimize.linprog(c, A_ub=M_ub, b_ub=b_ub,
                                 A_eq=M_eq, b_eq=b_eq, bounds=bounds)
    return res.x[:m]
\end{lstlisting}

\begin{center}
\renewcommand{\arraystretch}{1.2}
\small
\begin{tabular}{llc}
\toprule
Gioco & Strategia minimax della riga & Valore \\
\midrule
Sasso-carta-forbici & $(\frac13, \frac13, \frac13)$ & $0$ \\
RPS modificato ($4\times3$) & $(0;\ 0{,}125;\ 0{,}5;\ 0{,}375)$ & $0{,}125$ \\
RPS modificato, $A^T$ ($3\times4$) & $(0{,}25;\ 0{,}4167;\ 0{,}3333)$ & $0{,}0833$ \\
Matching pennies & $(\frac12, \frac12)$ & $0$ \\
Casuale $5\times5$, seme 42 & $(0;\ 0{,}5629;\ 0{,}4371;\ 0;\ 0)$ & $50{,}28$ \\
\bottomrule
\end{tabular}
\end{center}

\begin{attenzione}[il commento del notebook sull'RPS modificato]
Nel notebook, sopra l'RPS modificato, c'è il commento ``It should print [1/3, 1/3, 1/3]''.
È un residuo copiato dalla cella precedente: la matrice è $4\times3$ e la soluzione corretta
(confermata dall'output del docente nel notebook della lezione 4) è quella in tabella.
\end{attenzione}

\paragraph{Il bug dei limiti, verificato.}
Il gioco $A = \begin{pmatrix} -2 & -1 \\ -1 & -3 \end{pmatrix}$ ha valore
$-\tfrac53$ e strategia ottima $(\tfrac23, \tfrac13)$. Con \texttt{u} libera,
\texttt{linprog} restituisce esattamente questo. Con i limiti di default ($u \ge 0$)
restituisce \texttt{status = 2} (programma impossibile) e \texttt{res.x = None}: il codice
successivo si rompe con un \texttt{TypeError} che non dice nulla sulla causa.

\subsection{Equilibri misti nei giochi \texorpdfstring{$2\times2$}{2x2}}

Dalla condizione di indifferenza della riga
$qA_{11} + (1-q)A_{12} = qA_{21} + (1-q)A_{22}$ si ricava la probabilità $q$ con cui la
\emph{colonna} gioca la prima azione, e simmetricamente $p$ per la riga:
\[
  q = \frac{A_{22} - A_{12}}{A_{11} - A_{12} - A_{21} + A_{22}},
  \qquad
  p = \frac{B_{22} - B_{21}}{B_{11} - B_{12} - B_{21} + B_{22}} .
\]

\begin{lstlisting}
def mixed_nash_equilibrium_2x2(A, B):
    A, B = np.asarray(A, float), np.asarray(B, float)
    den_q = A[0, 0] - A[0, 1] - A[1, 0] + A[1, 1]
    den_p = B[0, 0] - B[0, 1] - B[1, 0] + B[1, 1]
    if den_q == 0 or den_p == 0:
        raise ValueError("nessun equilibrio completamente misto")
    q = (A[1, 1] - A[0, 1]) / den_q      # rende indifferente la riga
    p = (B[1, 1] - B[1, 0]) / den_p      # rende indifferente la colonna
    if not (0 <= p <= 1 and 0 <= q <= 1):
        raise ValueError("niente equilibrio misto interno")
    return np.array([p, 1 - p]), np.array([q, 1 - q])
\end{lstlisting}

\begin{center}
\renewcommand{\arraystretch}{1.2}
\small
\begin{tabular}{lll}
\toprule
Gioco & Riga & Colonna \\
\midrule
Battaglia dei sessi & $(\frac23, \frac13)$ & $(\frac13, \frac23)$ \\
Attraversamento ($A = [[-10,5],[0,-1]]$) & $(0{,}375;\ 0{,}625)$ & $(0{,}375;\ 0{,}625)$ \\
Matching pennies & $(\frac12, \frac12)$ & $(\frac12, \frac12)$ \\
$A = [[4,5],[6,3]],\ B = [[2,1],[0,3]]$ & $(0{,}75;\ 0{,}25)$ & $(0{,}5;\ 0{,}5)$ \\
$A = [[3,1],[0,2]],\ B = [[2,1],[0,3]]$ & $(0{,}75;\ 0{,}25)$ & $(0{,}25;\ 0{,}75)$ \\
$A = [[2,3],[3,-1]],\ B = [[-5,-6],[-2,0]]$ & $(\frac23, \frac13)$ & $(0{,}8;\ 0{,}2)$ \\
\bottomrule
\end{tabular}
\end{center}
Tutti coincidono con il notebook, e \texttt{is\_nash\_equilibrium} li conferma.

\begin{attenzione}[la formula non funziona sempre]
La formula presuppone che esista un equilibrio \emph{completamente misto}. Non sempre c'è:
\begin{itemize}
  \item nel \kw{dilemma del prigioniero} il denominatore di $q$ è
        $-1 + 4 - 0 - 3 = 0$: divisione per zero. Una strategia dominante non lascia
        spazio all'indifferenza;
  \item nel gioco $C/D\times A/B$ di inizio lezione esce $p = -0{,}5$: una
        ``probabilità'' negativa. L'unico equilibrio è quello puro $(D, B)$.
\end{itemize}
Una versione da usare nel progetto deve controllare il denominatore e che $p, q \in [0,1]$,
come il codice sopra, e in quei casi cercare gli equilibri puri.
\end{attenzione}

\subsection{Giochi potenziali e rilevamento di comunità}
\fonte{Lab03-Non\_Cooperative\_Game\_Theory.pdf}

\begin{definizione}[Gioco potenziale esatto]
Un gioco $\langle N, S, u\rangle$ è un \kw{gioco potenziale esatto} se esiste una
\kw{funzione potenziale} $\Phi : S \to \R$ tale che, per ogni giocatore $i$, ogni profilo
$s_{-i}$ degli altri e ogni coppia di strategie $s_i, s_i' \in S_i$:
\[
  u_i(s_i', s_{-i}) - u_i(s_i, s_{-i}) = \Phi(s_i', s_{-i}) - \Phi(s_i, s_{-i}).
\]
\end{definizione}
In parole: quando un giocatore cambia strategia da solo, la variazione del \emph{suo} payoff
è esattamente la variazione del potenziale \emph{comune} $\Phi$.
\begin{itemize}
  \item I giochi potenziali (finiti) hanno sempre un \kw{equilibrio di Nash puro}.
  \item Hanno la \kw{proprietà del miglioramento finito} (FIP): ogni sequenza di
        deviazioni unilaterali che migliorano il payoff di chi devia è finita.
\end{itemize}

\begin{extra}[perché le due proprietà valgono]
Ogni deviazione che migliora il payoff di chi devia aumenta $\Phi$ della stessa quantità.
Siccome $S$ è finito, $\Phi$ assume un numero finito di valori e non può crescere per
sempre: ogni sequenza di miglioramenti termina (FIP). Dove termina, nessuno può migliorare
da solo: è un equilibrio di Nash puro. In particolare il profilo che \emph{massimizza}
$\Phi$ è sempre un equilibrio.
\end{extra}

\paragraph{Dinamica di best response (BRD).}
Dalla FIP segue un algoritmo che converge sempre a un equilibrio puro:
\begin{lstlisting}[style=pseudo]
while not converged:
    for each player $i \in N$:
        if $s_i \notin BR(s_{-i})$:
            $s_i$ = BR($s_{-i}$)
\end{lstlisting}

\paragraph{Il gioco delle comunità.}
Dato un grafo $G = (V,E)$, si vogliono trovare gruppi di nodi densamente connessi al loro
interno e poco connessi tra loro (\kw{community detection}). Il problema è NP-difficile e
ha molte euristiche (Louvain, Infomap, Walktrap, \emph{label propagation}). Il laboratorio
mostra che la label propagation è una dinamica di best response:
\begin{itemize}
  \item ogni nodo è un giocatore: $N = V$;
  \item ogni nodo sceglie un cluster: $S_i = \{1, 2, \dots, |V|\}$;
  \item l'utilità di un nodo è il numero di vicini nello stesso cluster:
        $u_i(s_i, s_{-i}) = \sum_{j \in \Gamma(i)} \mathbb{I}(s_i = s_j)$;
  \item il potenziale è $\Phi(s) = \tfrac12 \sum_{i \in V} \sum_{j \in \Gamma(i)}
        \mathbb{I}(s_i = s_j)$, cioè il numero di archi con estremi nello stesso cluster.
\end{itemize}

\begin{extra}[verifica che $\Phi$ è un potenziale]
Ogni arco $\{i,j\}$ compare due volte nella doppia somma (da $i$ e da $j$), da cui il
$\tfrac12$. Quando $i$ cambia cluster cambiano solo gli archi incidenti a $i$, e ciascuno
contribuisce a $\Phi$ con $1$ se gli estremi sono nello stesso cluster: la variazione di
$\Phi$ è la variazione del numero di vicini di $i$ nel suo cluster, cioè esattamente
$\Delta u_i$.
\end{extra}

Non serve calcolare $\Phi$: è solo la garanzia matematica che il ciclo termina. Basta
assegnare a ogni nodo un cluster iniziale (il proprio identificativo) e poi, finché qualcosa
cambia, spostare ogni nodo nel cluster più frequente tra i suoi vicini.

\begin{lstlisting}
import random
from collections import Counter

def best_response(node, G, action_profile):
    counts = Counter(action_profile[nb] for nb in G.neighbors(node))
    if not counts:
        return action_profile[node]
    best = max(counts.values())
    if counts.get(action_profile[node], 0) == best:
        return action_profile[node]          # gia' in un cluster ottimo: resta
    return random.choice([c for c, k in counts.items() if k == best])

def BRD(G):
    profile = {node: node for node in G.nodes()}   # ognuno nel proprio cluster
    nodes = list(G.nodes())
    converged = False
    while not converged:
        converged = True
        random.shuffle(nodes)
        for node in nodes:
            new = best_response(node, G, profile)
            if new != profile[node]:
                profile[node] = new
                converged = False
    return [profile[n] for n in sorted(G.nodes())]
\end{lstlisting}

\begin{attenzione}[i pareggi possono impedire la convergenza]
La garanzia di terminazione vale solo se un nodo si sposta quando \emph{migliora
strettamente}. Un nodo che ha due cluster ugualmente frequenti tra i vicini e cambia anche a
pari merito non aumenta $\Phi$, e la dimostrazione non vale più. Verificato su un grafo con
due cricche da 5 nodi e un nodo collegato a entrambe: con la regola ``resta se sei già in un
cluster ottimo'' la dinamica termina al primo turno; con la regola ``passa all'altro
cluster ottimo'' non termina mai. Il controllo nella \texttt{best\_response} sopra serve a
questo.
\end{attenzione}

\paragraph{Risultati.}
Il grafo di test è un \kw{modello a blocchi stocastico}: 4 comunità da 100 nodi, probabilità
di arco $0{,}2$ dentro una comunità e $0{,}01$ tra comunità diverse. Su 20 grafi generati con
semi diversi:
\begin{itemize}
  \item in 19 casi su 20 la dinamica trova \emph{esattamente} le 4 comunità vere, con
        purezza 1 (nessun nodo assegnato al blocco sbagliato);
  \item converge in media in $3{,}8$ turni (tra 3 e 5), con 485--595 spostamenti di nodo;
  \item il potenziale cresce strettamente a ogni spostamento, come previsto (verificato
        su ogni mossa di ogni esecuzione);
  \item nel caso rimanente due comunità vengono fuse: 3 cluster, purezza $0{,}75$.
\end{itemize}

\begin{osservazione}[il massimo di $\Phi$ è la soluzione sbagliata]
Il profilo con il potenziale \emph{massimo} è ``tutti i nodi nello stesso cluster'':
$\Phi = |E|$, perché ogni arco ha gli estremi nello stesso gruppo. È un equilibrio di Nash
(nessuno vuole lasciare un gruppo che contiene tutti i suoi vicini), ma non dice nulla sulle
comunità. Infatti l'unica esecuzione che ha fuso due comunità è anche quella con il
potenziale relativo più alto: l'89\% degli archi finisce dentro un cluster, contro l'86--88\%
delle esecuzioni che trovano le 4 comunità giuste. La label propagation funziona perché si
ferma in un \emph{massimo locale}
del potenziale, raggiunto partendo da etichette tutte diverse, e non perché lo massimizza.
È un buon esempio della differenza tra ``trovare un equilibrio'' e ``trovare l'equilibrio
giusto''.
\end{osservazione}

% ---------------------------------------------------------------------
\begin{riepilogo}
\begin{itemize}
  \item Un gioco in forma normale è $\langle N, A, u\rangle$; con due giocatori si
        rappresenta con due matrici $A$ (riga) e $B$ (colonna).
  \item Una strategia strettamente dominata non è mai giocata. IESDS la elimina; se resta
        un solo profilo è l'unico equilibrio di Nash. Non tutti i giochi si risolvono così.
  \item Equilibrio di Nash: ognuno gioca una best response agli altri. Può essere unico,
        multiplo (battaglia dei sessi) o assente in strategie pure (rigori).
  \item Tragedia dei beni comuni: all'equilibrio si estrae $\tfrac{2A}{3B}$ invece
        dell'ottimo $\tfrac{A}{2B}$, e ognuno ottiene $\tfrac{A^2}{9B}$ invece di
        $\tfrac{A^2}{8B}$.
  \item Strategie miste: utilità attesa $\sigma_r A\sigma_c^T$. \textbf{Nash (1950)}: ogni
        gioco finito ha un equilibrio misto.
  \item Condizione del supporto: una strategia è best response se e solo se mette
        probabilità solo su azioni pure ottime. All'equilibrio ognuno è indifferente sul
        proprio supporto; nei $2\times2$ l'indifferenza dà l'equilibrio
        ($\tfrac23,\tfrac13$; $\tfrac13,\tfrac23$ nella battaglia dei sessi).
  \item Somma zero: \textbf{teorema minimax}, $\max\min = \min\max$; strategie ottime e
        valore si calcolano con un programma lineare (attenzione a lasciare libera la
        variabile del valore).
  \item \textbf{Lab}: verifica di best response ed equilibri, minimax con \texttt{linprog},
        equilibri $2\times2$; label propagation come dinamica di best response in un gioco
        potenziale (4 comunità ritrovate in 19 grafi su 20).
\end{itemize}
\textbf{Prossima lezione}: esistenza dell'equilibrio di Nash, complessità del calcolo degli
equilibri, equilibri correlati e apprendimento nei giochi.
\end{riepilogo}

% =====================================================================
%  Lezione 4 — Equilibri alternativi e apprendimento nei giochi
%  Materiale: Lesson_4_Non_Cooperative_Game_Theory_II.pdf,
%             Lesson4-NonCooperativeGameTheory_with_Learning_in_Games.ipynb
% =====================================================================
\chapter{Equilibri alternativi e apprendimento}

\parteteoria

La lezione 3 ha lasciato quattro concetti di soluzione, ognuno con un limite:
\begin{center}
\renewcommand{\arraystretch}{1.25}
\begin{tabularx}{\linewidth}{@{}l X@{}}
\toprule
Concetto & Limite \\
\midrule
Eliminazione iterata (IESDS) & non tutti i giochi sono risolvibili per dominanza \\
Equilibrio di Nash puro & non tutti i giochi ne hanno uno \\
Equilibrio di Nash misto & esiste sempre nei giochi finiti (Nash, 1950), ma è difficile da
calcolare oltre i giochi piccoli \\
Minimax & vale per i giochi a somma zero, dove si calcola in modo efficiente \\
\bottomrule
\end{tabularx}
\end{center}

Questa lezione risponde a tre domande:
\begin{enumerate}[label=\arabic*.]
  \item \emph{perché} l'equilibrio di Nash misto esiste sempre;
  \item quanto è \emph{difficile} calcolarlo;
  \item che cosa si può fare in alternativa: \kw{equilibri correlati} e \kw{apprendimento}
        nei giochi ripetuti.
\end{enumerate}

% ---------------------------------------------------------------------
\section{Esistenza dell'equilibrio di Nash misto}

\begin{teorema}[Nash, 1950]
Ogni gioco finito ha un equilibrio di Nash misto.
\end{teorema}

\textbf{Idea della dimostrazione}: \emph{ridurre} l'esistenza di un equilibrio misto
all'esistenza di un \kw{punto fisso} di una funzione.

\subsection{Il teorema del punto fisso di Brouwer}

\begin{teorema}[Brouwer, 1911]
Ogni funzione \kw{continua} $f$ che manda un insieme euclideo \kw{compatto} e \kw{convesso}
in sé stesso ha un punto $x_0$ tale che $f(x_0) = x_0$.
\end{teorema}

Ricordando le definizioni:
\begin{itemize}
  \item un insieme di $\R^d$ è \kw{compatto} se è chiuso e limitato;
  \item un insieme $C$ è \kw{convesso} se per ogni $x, y \in C$ e ogni
        $\epsilon \in [0,1]$ vale $\epsilon x + (1-\epsilon)y \in C$: contiene il segmento tra
        due suoi punti qualsiasi.
\end{itemize}

Ogni ipotesi serve. Le slide mostrano un controesempio per ciascuna, sull'intervallo
(o sull'anello) e con il grafico di $f$ confrontato con la bisettrice $y = x$: un punto fisso
è un punto in cui il grafico incontra la bisettrice.

\begin{center}
\renewcommand{\arraystretch}{1.3}
\begin{tabularx}{\linewidth}{@{}l X@{}}
\toprule
Ipotesi tolta & Funzione senza punto fisso \\
\midrule
Insieme limitato & $f(x) = x + 0{,}1$ su $\R$: una retta parallela alla bisettrice \\
Insieme chiuso & $f(x) = x^2$ sull'aperto $(0,1)$: i punti fissi di $x^2$ sono $0$ e $1$,
 entrambi fuori dall'insieme \\
Insieme convesso & una rotazione dell'anello attorno al centro: il centro, che sarebbe il punto
fisso, non sta nell'insieme \\
Funzione continua & $f(x) = x + \tfrac12$ per $x < \tfrac12$, $f(x) = x - \tfrac12$ per
 $x \ge \tfrac12$ su $[0,1]$: il grafico ``salta'' oltre la bisettrice \\
\bottomrule
\end{tabularx}
\end{center}

\subsection{La costruzione di Nash}

Il piano della dimostrazione:
\begin{enumerate}
  \item mostrare che l'insieme $S = S_1 \times S_2 \times \dots \times S_n$ dei profili di
        strategie miste è compatto e convesso;
  \item costruire una funzione continua $f : S \to S$ con $f(\sigma) = \sigma$ se e solo se
        $\sigma$ è un equilibrio misto;
  \item applicare Brouwer a $f$.
\end{enumerate}

\paragraph{Passo 1: l'insieme dei profili.}
Le strategie miste del giocatore $i$ formano il \kw{simplesso standard}
\[
  S_i = \bigl\{ x \in \R^{|A_i|} \;:\; x \ge \mathbf{0},\ \uno^T x = 1 \bigr\},
\]
che è compatto e convesso. Il prodotto cartesiano di compatti è compatto e il prodotto di
convessi è convesso, quindi anche $S$ lo è.

\begin{attenzione}[``compatto e chiuso'']
Le slide scrivono che $S$ è ``convesso, compatto e chiuso'' nel piano della dimostrazione e
poi solo ``compatto e chiuso'' nel passo 1. In $\R^d$ un compatto è già chiuso; ciò che
serve per Brouwer, e che i due teoremi enunciati subito dopo dimostrano, è
\emph{compatto e convesso}.
\end{attenzione}

\paragraph{Passo 2: la funzione.}
L'idea è una funzione che prende un profilo e ne restituisce uno \emph{diverso}, a meno che
il profilo non sia un equilibrio. La best response sarebbe la scelta naturale, ma non è una
funzione (può avere più valori) e può essere discontinua. Si usa invece il \kw{guadagno}
che $i$ otterrebbe passando dalla sua strategia mista all'azione pura $a_i^j$:
\[
  g\bigl(a_i^j, \sigma\bigr) = \max\Bigl\{0,\ u_i\bigl(a_i^j, \sigma_{-i}\bigr) - u_i(\sigma)\Bigr\},
\]
e si sposta la probabilità verso le azioni che fanno guadagnare:
\[
  f(\sigma) = \sigma', \qquad
  \sigma'_{i,a_j} = \frac{\sigma_{i,a_j} + g(a_i^j, \sigma)}
                         {\sum_{k \in A_i} \bigl(\sigma_{i,a_k} + g(a_i^k, \sigma)\bigr)} .
\]
La funzione è continua (composizione di funzioni continue, con denominatore $\ge 1$), manda
$S$ in sé (il risultato è ancora una distribuzione di probabilità) e ha un punto fisso se e
solo se $\sigma$ è un equilibrio misto. Per Brouwer un punto fisso esiste, quindi esiste un
equilibrio. $\square$

\begin{extra}[il passaggio che le slide non dimostrano]
Che un equilibrio sia un punto fisso è facile: all'equilibrio nessuna azione pura migliora,
tutti i $g$ sono 0 e $f(\sigma) = \sigma$. Il viceversa è il passo delicato. Sia $\sigma$
un punto fisso e sia $G_i = \sum_k g(a_i^k, \sigma)$; allora per ogni azione
\[
  \sigma_{i,a_j} = \frac{\sigma_{i,a_j} + g(a_i^j,\sigma)}{1 + G_i}
  \quad\Longleftrightarrow\quad
  \sigma_{i,a_j}\,G_i = g(a_i^j, \sigma).
\]
L'utilità $u_i(\sigma)$ è la media delle $u_i(a_i^j, \sigma_{-i})$ pesata con $\sigma_i$,
quindi almeno un'azione del supporto ha $u_i(a_i^j,\sigma_{-i}) \le u_i(\sigma)$, cioè
$g(a_i^j,\sigma) = 0$. Per quell'azione $\sigma_{i,a_j} > 0$ e $\sigma_{i,a_j} G_i = 0$,
quindi $G_i = 0$: \emph{nessuna} azione pura fa guadagnare $i$, e $\sigma_i$ è una best
response. Vale per ogni $i$, quindi $\sigma$ è un equilibrio.
\end{extra}

\begin{attenzione}[refusi]
Il nome è \emph{Brouwer}, non ``Brower''; e $f$ manda $S$ \emph{in} sé stesso (\en{into}),
non necessariamente \emph{su} tutto $S$ (\en{onto}): Brouwer non richiede la suriettività.
\end{attenzione}

% ---------------------------------------------------------------------
\section{Complessità del calcolo degli equilibri}

\subsection{Equilibri misti: PPAD-completezza}

\begin{itemize}
  \item La riduzione appena vista mostra che trovare un equilibrio misto non è più difficile
        che trovare un punto fisso.
  \item Una riduzione (molto più elaborata) mostra il contrario: trovare un punto fisso non
        è più difficile che trovare un equilibrio misto.
  \item Insieme: trovare un equilibrio misto è \emph{esattamente} difficile quanto trovare
        un punto fisso. Il problema è \kw{PPAD-completo}.
\end{itemize}

\begin{definizione}[La classe PPAD]
\kw{PPAD} (\en{Polynomial Parity Argument on Directed graphs}) contiene i problemi di punto
fisso. Il suo problema rappresentativo è \kw{End of Line}: è dato implicitamente un grafo
orientato $G = (V,E)$ con $|V| = 2^n$ nodi, in cui ogni nodo ha grado entrante e uscente al
più 1, tramite due circuiti
\[
  S : \{0,1\}^n \to \{0,1\}^n \ (\text{successore}),\qquad
  P : \{0,1\}^n \to \{0,1\}^n \ (\text{predecessore}).
\]
Dato un nodo sorgente $s$, si deve trovare un pozzo $t$ (la ``fine della linea'').
\end{definizione}

\begin{osservazione}[perché è difficile e perché la soluzione esiste]
La soluzione esiste sempre: seguendo i successori da $s$ si arriva prima o poi a un nodo
senza successore, perché i nodi sono finiti e non si possono ripetere. Ma il grafo ha
$2^n$ nodi e la linea può essere lunga in modo esponenziale; i circuiti $S$ e $P$ non
dicono dove finisce. È lo stesso tipo di problema di Nash: si sa che la soluzione c'è,
non si sa come trovarla in fretta.
\end{osservazione}

\paragraph{Il lemma di Sperner.}
È la controparte discreta di Brouwer. Si prende un reticolo triangolato e se ne colorano i
nodi di bordo con tre colori secondo la regola delle slide (rosso sul lato sinistro, giallo
in basso, blu in alto e a destra). Comunque si colorino i nodi interni, esiste un
\kw{triangolo tricromatico}, anzi ne esiste un numero \emph{dispari}. Trovarne uno è
PPAD-completo.

L'argomento che lo dimostra è proprio un ``End of Line'': si considerano i triangoli come
nodi e le ``porte'' rosso-giallo come archi, e si attraversa ogni porta tenendo il rosso a
sinistra. Partendo dalla porta sul bordo si cammina di triangolo in triangolo, e il cammino
può finire solo in un triangolo tricromatico. Le slide riassumono la catena di riduzioni:
\[
  \textsc{EoL} \le_p \textsc{Sperner} \le_p \textsc{Brouwer} \le_p \textsc{Nash}.
\]

\begin{osservazione}[come leggere la catena]
$X \le_p Y$ significa ``$X$ si riduce a $Y$'': chi sa risolvere $Y$ sa risolvere $X$, quindi
$Y$ è almeno difficile quanto $X$. La catena dice che Nash è almeno difficile quanto End of
Line, cioè PPAD-\emph{difficile}. La costruzione di Nash (Brouwer applicato a $f$) dà la
direzione opposta, l'appartenenza a PPAD. Insieme: PPAD-completo.
\end{osservazione}

\subsection{Equilibri puri: NP-difficoltà}

\kw{Trovare un equilibrio di Nash puro è NP-difficile.} Lo si dimostra riducendo a esso
il problema della \kw{3-colorabilità}: dato un grafo $G = (V,E)$, decidere se si possono
colorare i nodi con tre colori senza che due nodi adiacenti abbiano lo stesso colore.

Si usano i \kw{giochi su grafo}: i giocatori sono i nodi e il payoff di ciascuno dipende
solo dalle interazioni locali (la sua azione e quelle dei vicini).

\paragraph{Primo tentativo.}
Nodi = giocatori, azioni = $\{R, G, B\}$, payoff $+10$ se il colore scelto non è in
conflitto con nessun vicino, $-10$ se lo è. Se $G$ è 3-colorabile c'è un equilibrio con
benessere massimo; se non lo è c'è comunque un equilibrio, solo con benessere più basso.
Questo dimostra che trovare l'equilibrio \emph{di benessere massimo} è NP-difficile, ma non
che lo sia trovarne \emph{uno}.

\paragraph{La costruzione giusta.}
Si aggiunge un \kw{nodo sfidante} collegato a tutti gli altri:
\begin{itemize}
  \item azioni dei nodi del grafo: $\{R, G, B, X, Y\}$; azioni dello sfidante: $\{X, Y\}$;
  \item payoff di un nodo: $+10$ se sceglie un colore senza conflitti, $-10$ se il colore è
        in conflitto, $0$ se sceglie $X$/$Y$ \emph{copiando} lo sfidante, $-5$ se sceglie
        $X$/$Y$ senza copiarlo;
  \item payoff dello sfidante: $+10$ se nessun nodo lo copia, $-10$ se qualcuno lo copia.
\end{itemize}
Allora il gioco ha un equilibrio puro \emph{se e solo se} $G$ è 3-colorabile.

\begin{extra}[perché funziona]
\begin{itemize}
  \item \emph{Se $G$ è 3-colorabile}: i nodi giocano una colorazione corretta ($+10$
        ciascuno) e lo sfidante una lettera qualsiasi ($+10$, nessuno lo copia). Nessun
        nodo ha interesse a passare a una lettera ($\le 0 < 10$), lo sfidante è già al
        massimo: equilibrio.
  \item \emph{Se non lo è}, si supponga per assurdo un equilibrio. Un nodo che gioca una
        lettera senza copiare prende $-5$ e preferirebbe copiare ($0$); quindi in
        equilibrio chi gioca una lettera copia lo sfidante. Ma se almeno un nodo lo copia,
        lo sfidante prende $-10$ e, cambiando lettera, nessuno lo copierebbe più ($+10$):
        devierebbe. Quindi nessun nodo gioca lettere: tutti giocano colori. Poiché $G$ non
        è 3-colorabile, qualche nodo è in conflitto ($-10$) e preferirebbe copiare lo
        sfidante ($0$). Contraddizione.
\end{itemize}
Verificato con una ricerca esaustiva su tutti i profili: il triangolo e il ciclo da 4 nodi
(3-colorabili) hanno un equilibrio puro, il grafo completo $K_4$ (non 3-colorabile) no.
\end{extra}

\begin{attenzione}[NP-difficile rispetto a che cosa]
Se un gioco è dato in forma normale \emph{esplicita} (la tabella completa dei payoff),
trovare un equilibrio puro è facile: basta controllare ogni cella, in tempo polinomiale
nella dimensione della tabella. La difficoltà nasce perché nei giochi su grafo la
rappresentazione è \emph{compatta}: con $|V|$ nodi e 5 azioni la tabella avrebbe
$5^{|V|}\cdot 2$ celle, ma il gioco si descrive con poche regole locali. Il risultato
vale rispetto a questa descrizione compatta, ed è quella che conta nei problemi reali.
\end{attenzione}

\subsection{Il quadro complessivo}

\begin{itemize}
  \item Trovare un equilibrio di Nash misto è \kw{PPAD-completo}.
  \item Trovare un equilibrio di Nash puro è \kw{NP-difficile}.
  \item Trovare l'equilibrio che massimizza il benessere sociale è \kw{NP-difficile}.
\end{itemize}
Conseguenza: se calcolare un equilibrio è difficile per un algoritmo centralizzato e
ottimizzato, non ci si può aspettare che un sistema di agenti autonomi che interagiscono
lo raggiunga da solo nei giochi generici. È la motivazione per le nozioni alternative di
equilibrio e per l'apprendimento.

% ---------------------------------------------------------------------
\section{Equilibri correlati}

\subsection{L'intuizione: il semaforo}

Il gioco dell'\kw{attraversamento}: due auto arrivano a un incrocio e ciascuna sceglie se
passare (\emph{Go}) o fermarsi (\emph{Stop}).
\begin{center}
\renewcommand{\arraystretch}{1.35}
\begin{tabular}{r|c|c|}
\multicolumn{1}{r}{} & \multicolumn{1}{c}{Go} & \multicolumn{1}{c}{Stop}\\ \cline{2-3}
Go   & \pay{-10}{-10} & \pay{5}{0} \\ \cline{2-3}
Stop & \pay{0}{5}     & \pay{-1}{-1} \\ \cline{2-3}
\end{tabular}
\end{center}
\begin{itemize}
  \item Due equilibri puri: (Go, Stop) con $u = (5, 0)$ e (Stop, Go) con $u = (0, 5)$.
        Sbilanciati.
  \item Un equilibrio misto: entrambi giocano Go con probabilità $\tfrac38$, con
        $u_1 = u_2 = -\tfrac58$. Equo ma pessimo: ci si scontra con probabilità
        $\tfrac{9}{64}$.
\end{itemize}

Nella realtà il problema si risolve con un \kw{semaforo}: un \emph{dispositivo di
coordinamento} che raccomanda un'azione a ciascun giocatore. Lo si modella come una
distribuzione di probabilità sui profili:
\begin{center}
\renewcommand{\arraystretch}{1.35}
\begin{tabular}{r|c|c|}
\multicolumn{1}{r}{} & \multicolumn{1}{c}{Go} & \multicolumn{1}{c}{Stop}\\ \cline{2-3}
Go   & $0$         & $\tfrac12$ \\ \cline{2-3}
Stop & $\tfrac12$  & $0$ \\ \cline{2-3}
\end{tabular}
\end{center}
Se entrambi seguono la raccomandazione ottengono $u_1 = u_2 = 2{,}5$. E \kw{deviare dalla
raccomandazione non conviene}, assumendo che l'altro la segua: chi riceve ``Stop'' sa che
l'altro ha ricevuto ``Go'' e, passando, prenderebbe $-10$ invece di $0$; chi riceve ``Go'' sa
che l'altro si ferma e, fermandosi, prenderebbe $-1$ invece di $5$.

\subsection{Definizione}

\begin{definizione}[Equilibrio correlato]
Un \kw{equilibrio correlato} (CE) è una distribuzione di probabilità sui profili di azioni
\[
  p : A_1 \times A_2 \times \dots \times A_n \to [0,1]
\]
tale che per ogni giocatore $i$ e ogni coppia di azioni $a_i, a_i' \in A_i$
\[
  \sum_{a_{-i} \in A_{-i}} p(a_i, a_{-i})\, u_i(a_i, a_{-i})
  \ \ge\ \sum_{a_{-i} \in A_{-i}} p(a_i, a_{-i})\, u_i(a_i', a_{-i}).
\]
\end{definizione}
In parole: se il dispositivo raccomanda $a_i$ al giocatore $i$, la probabilità condizionata
degli altri profili è proporzionale a $p(a_i, \cdot)$, e contro quella distribuzione seguire
la raccomandazione è almeno buono quanto passare a qualunque $a_i'$.

\begin{teorema}[Esistenza]
Ogni gioco finito ha un equilibrio correlato.
\end{teorema}

\begin{extra}[perché esiste, e che rapporto c'è con Nash]
Ogni equilibrio di Nash misto $\sigma$ è anche un equilibrio correlato: basta prendere
$p(a) = \prod_j \sigma_j(a_j)$, la distribuzione dei profili quando ognuno randomizza per
conto suo. L'esistenza segue quindi dal teorema di Nash. Il viceversa non vale: il semaforo
è un equilibrio correlato ma non è il prodotto di due strategie indipendenti (ha
probabilità nulla su (Go,Go) e su (Stop,Stop), cosa impossibile per un prodotto con
entrambe le marginali $(\tfrac12,\tfrac12)$). Gli equilibri correlati sono un insieme più
grande, e quindi possono contenere esiti migliori.
\end{extra}

\subsection{Le disuguaglianze nell'esempio}

Per il giocatore 1 (riga) le due raccomandazioni danno:
\begin{align*}
  \text{raccomandato Go:}\quad
  & p(G,G)\cdot(-10) + p(G,S)\cdot 5 \ \ge\ p(G,G)\cdot 0 + p(G,S)\cdot(-1)\\
  \text{raccomandato Stop:}\quad
  & p(S,G)\cdot 0 + p(S,S)\cdot(-1) \ \ge\ p(S,G)\cdot(-10) + p(S,S)\cdot 5
\end{align*}
e per il giocatore 2 (colonna):
\begin{align*}
  \text{raccomandato Go:}\quad
  & p(G,G)\cdot(-10) + p(S,G)\cdot 5 \ \ge\ p(G,G)\cdot 0 + p(S,G)\cdot(-1)\\
  \text{raccomandato Stop:}\quad
  & p(G,S)\cdot 0 + p(S,S)\cdot(-1) \ \ge\ p(G,S)\cdot(-10) + p(S,S)\cdot 5
\end{align*}
Semplificando:
\[
  p(G,S) \ge \tfrac53\,p(G,G), \quad
  p(S,G) \ge \tfrac35\,p(S,S), \quad
  p(S,G) \ge \tfrac53\,p(G,G), \quad
  p(G,S) \ge \tfrac35\,p(S,S).
\]
Il semaforo $p(G,G) = p(S,S) = 0$, $p(G,S) = p(S,G) = \tfrac12$ le soddisfa tutte.

\subsection{Calcolo con la programmazione lineare}

Le condizioni sono \kw{lineari} nelle incognite $p(a)$: un equilibrio correlato si trova
risolvendo un programma lineare, cioè in tempo polinomiale nella dimensione del gioco.
Gli equilibri correlati sono quindi \emph{più facili} da calcolare di quelli di Nash. Si può
scegliere l'obiettivo che si preferisce:
\[
\begin{aligned}
  \max\ & \text{[obiettivo a scelta]}\\
  \text{s.t. }\ & \sum_{a_{-i}} p(a_i, a_{-i})\bigl(u_i(a_i, a_{-i}) - u_i(a_i', a_{-i})\bigr) \ge 0
  && \forall i,\ \forall a_i, a_i' \in A_i\\
  & \textstyle\sum_{a \in A} p(a) = 1,\qquad p(a) \ge 0 .
\end{aligned}
\]
Due obiettivi tipici:
\begin{itemize}
  \item \kw{benessere totale}: $\max \sum_{a \in A} p(a) \sum_i u_i(a)$;
  \item \kw{regola maximin}: $\max \min_i \sum_{a \in A} p(a)\,u_i(a)$, cioè rendere il più
        alto possibile il payoff atteso del giocatore che sta peggio.
\end{itemize}

\begin{attenzione}[tre errori nelle slide del programma lineare]
\begin{itemize}
  \item Nel vincolo, a destra le slide scrivono $p(a_i', a_{-i})\,u_i(a_i', a_{-i})$. La
        probabilità corretta è $p(a_i, a_{-i})$, la stessa del lato sinistro: si sta
        confrontando che cosa succede \emph{con la stessa raccomandazione} se si gioca
        $a_i$ oppure $a_i'$. Manca anche la sommatoria a destra. La definizione data
        qualche slide prima è giusta; lo è anche il notebook del docente.
  \item L'obiettivo ``$\max \sum_i p(a)\,u_i(a)$'' manca della somma sui profili $a$.
  \item ``$\max\min_i p(a)\,u_i(a)$'' non è un programma lineare così com'è (il minimo non
        è lineare). Lo diventa con una variabile ausiliaria $t$:
        $\max t$ s.t.\ $t \le \sum_a p(a)\,u_i(a)$ per ogni $i$.
\end{itemize}
\end{attenzione}

\section{Altri concetti di soluzione}

\begin{itemize}
  \item \kw{Trattare il gioco come a somma zero} e calcolare una strategia minimax: approccio
        conservativo, utile se si crede che l'avversario voglia danneggiarci o se non si
        conosce la sua utilità.
  \item \kw{Minimizzare il massimo rimpianto}. Il \kw{rimpianto} (\en{regret}) di un'azione
        è quanto si perde rispetto alla migliore risposta, a posteriori:
        \[
          R(a_i, a_{-i}) = \Bigl[\max_{a_i' \in A_i} u_i(a_i', a_{-i})\Bigr] - u_i(a_i, a_{-i}),
          \qquad
          a_i^* \in \argmin_{a_i \in A_i}\ \max_{a_{-i} \in A_{-i}} R(a_i, a_{-i}).
        \]
        È meno conservativo del minimax, utile quando l'avversario è imprevedibile più che
        ostile.
\end{itemize}

\begin{extra}[minimax e minimo rimpianto nell'attraversamento]
Per il giocatore 1: contro Go la risposta migliore vale $0$ (Stop), contro Stop vale $5$
(Go). I rimpianti sono:
\begin{center}
\renewcommand{\arraystretch}{1.2}
\begin{tabular}{lccc}
\toprule
 & contro Go & contro Stop & rimpianto massimo \\
\midrule
Go   & $0 - (-10) = 10$ & $5 - 5 = 0$ & $10$ \\
Stop & $0 - 0 = 0$ & $5 - (-1) = 6$ & $\mathbf{6}$ \\
\bottomrule
\end{tabular}
\end{center}
Il minimo rimpianto massimo è Stop. Anche il criterio maximin sceglie Stop (caso peggiore
$-1$ contro $-10$). Qui i due criteri coincidono, ma in generale no: il maximin guarda solo
i payoff assoluti, il rimpianto li confronta con quello che si sarebbe potuto ottenere.
\end{extra}

% ---------------------------------------------------------------------
\section{Apprendimento nei giochi}

Si immagini ora che i giocatori giochino lo stesso gioco $G$ \emph{ripetutamente}, per $T$
turni. Si vogliono algoritmi adattivi con cui i giocatori \kw{imparino} a giocare strategie
di equilibrio. Poiché calcolare un equilibrio di Nash è PPAD-completo, non ci si aspetta di
impararlo in \emph{ogni} gioco. Due algoritmi classici:
\begin{itemize}
  \item \kw{fictitious play} (Robinson, 1951);
  \item \kw{regret matching} (Hart e Mas-Colell, 2000).
\end{itemize}

\subsection{Fictitious play}

Ogni giocatore stima la strategia degli avversari dalle frequenze osservate e gioca la best
response a quella stima.
\begin{enumerate}
  \item A ogni turno aggiorna il conteggio di quante volte gli avversari hanno giocato
        ciascuna azione: $\text{count}_{a_{-i}}[a_j]$.
  \item Normalizza il conteggio in una distribuzione, la stima della strategia avversaria:
        \[
          \sigma_{-i}^t[a_j] = \frac{\text{count}_{a_{-i}}[a_j]}{t - 1}.
        \]
  \item Gioca una best response alla stima:
        $a_i^t \in \argmax_{a_i} \Ev\bigl[u_i(a_i, a_{-i}) \mid \sigma_{-i}^t\bigr]$.
\end{enumerate}

Proprietà:
\begin{itemize}
  \item aggiorna le credenze assumendo che gli altri giocatori siano \kw{stazionari}
        (giochino sempre la stessa strategia mista), il che non è vero: anche loro imparano;
  \item non è garantito che converga; ma \emph{se} le frequenze $\sigma_{-i}^t$ convergono,
        convergono a un equilibrio di Nash;
  \item converge sicuramente a un equilibrio di Nash nei giochi a due giocatori a somma zero,
        nei giochi risolvibili per eliminazione iterata delle strategie strettamente
        dominate e nei giochi a due giocatori $2\times N$.
\end{itemize}

\begin{osservazione}[che cosa converge]
Le \emph{azioni} giocate in fictitious play sono sempre pure e possono ciclare per sempre:
in sasso-carta-forbici l'algoritmo gioca sasso per un po', poi carta, poi forbici, \ldots\ Ciò
che converge all'equilibrio $(\tfrac13,\tfrac13,\tfrac13)$ sono le \emph{frequenze
empiriche}, cioè le medie nel tempo. Il grafico delle slide mostra proprio queste
oscillazioni sempre più strette attorno a $\tfrac13$.
\end{osservazione}

\subsection{Regret matching}

Ogni giocatore decide quanto spesso giocare un'azione in base a \emph{quanto avrebbe
guadagnato in più} giocandola sempre, col senno di poi.
\begin{enumerate}
  \item A ogni turno calcola il \kw{rimpianto cumulato} di ciascuna azione $a$ rispetto a
        ciò che ha effettivamente giocato:
        \[
          R_i^t(a) = \sum_{s \le t} \Bigl( u_i\bigl(a, \sigma_{-i}^s\bigr)
                     - u_i\bigl(\sigma_i^s, \sigma_{-i}^s\bigr) \Bigr).
        \]
  \item Aggiorna la strategia in proporzione alla \emph{parte positiva} del rimpianto:
        \[
          \sigma_i^{t+1}(a) = \frac{\max(0, R_i^t(a))}{\sum_{a'} \max(0, R_i^t(a'))},
        \]
        e gioca in modo uniforme se tutti i rimpianti sono $\le 0$.
\end{enumerate}

\begin{attenzione}[la formula delle slide è un'altra cosa]
Le slide scrivono $R_i^t(a) = \sum_{s \le t} \max\bigl(0,\ u_i(a,\sigma_{-i}^s) -
u_i(\sigma^s)\bigr)$: la somma delle parti positive di ogni turno, invece della parte
positiva della somma. Non è la stessa cosa: con la versione delle slide i rimpianti non
possono mai diminuire, e un'azione che è stata buona in passato e poi pessima continua a
essere giocata. Il regret matching di Hart e Mas-Colell, quello con le garanzie
teoriche, usa $\max(0, \sum_s \dots)$.

L'effetto si misura nel laboratorio. In sasso-carta-forbici entrambe le versioni sembrano
funzionare, perché per simmetria finiscono comunque attorno a $\tfrac13$. Su un gioco a
somma zero asimmetrico, con strategia minimax $(0{,}33;\ 0{,}56;\ 0{,}11)$, la versione
corretta converge a $(0{,}34;\ 0{,}56;\ 0{,}10)$, quella delle slide a
$(0{,}41;\ 0{,}35;\ 0{,}24)$, che non è un equilibrio. Le slide inoltre non dicono che
cosa fare quando tutti i rimpianti sono nulli (denominatore zero): si gioca uniforme.
\end{attenzione}

Proprietà secondo le slide:
\begin{itemize}
  \item nei giochi generali la media nel tempo converge all'insieme degli
        \kw{equilibri correlati};
  \item converge a un equilibrio di Nash nelle stesse classi di giochi di fictitious play
        (somma zero a due giocatori, risolvibili per dominanza, $2\times N$).
\end{itemize}

\begin{attenzione}[che cosa converge, e a che cosa]
Due precisazioni sull'enunciato delle slide, utili se all'esame si chiede ``a che cosa
converge il regret matching''.
\begin{itemize}
  \item \emph{Che cosa converge.} Le slide parlano della media delle strategie,
        $\bar\sigma^T = \tfrac1T\sum_t \sigma_t$. Ma un equilibrio correlato è una
        distribuzione sui \emph{profili congiunti}, non una strategia per giocatore: ciò
        che converge è la \emph{distribuzione empirica dei profili giocati}, cioè la
        frazione di turni in cui si è giocato ciascun $(a_1,\dots,a_n)$.
  \item \emph{A che cosa.} La versione vista sopra confronta ogni azione con la strategia
        giocata nel suo complesso (rimpianto \emph{incondizionato}, o ``esterno''). La
        teoria garantisce allora la convergenza agli \kw{equilibri correlati grossolani}
        (\en{coarse correlated equilibria}): distribuzioni da cui nessun giocatore vorrebbe
        sganciarsi \emph{prima} di ricevere la raccomandazione, per giocare sempre una
        stessa azione fissa. La convergenza agli equilibri correlati veri e propri vale per
        la variante del paper di Hart e Mas-Colell con rimpianti \emph{condizionati}, che
        confronta ogni azione effettivamente giocata con ciascuna alternativa (``avrei
        fatto meglio a giocare $k$ \emph{ogni volta che} ho giocato $j$?'').
\end{itemize}
\end{attenzione}

\begin{extra}[perché nei giochi a somma zero si arriva comunque a Nash]
Nei giochi a due giocatori a somma zero la distinzione non conta: le marginali di qualunque
equilibrio correlato, anche grossolano, sono strategie minimax. Se nessuno dei due ha
rimpianti verso un'azione fissa, ciascuno ottiene almeno il valore del gioco; poiché la somma
dei payoff è zero, entrambi ottengono esattamente il valore, e le frequenze di ciascuno sono
una strategia ottima. Analogamente, un giocatore senza rimpianti gioca sempre più di rado
un'azione strettamente dominata da un'altra azione pura (ogni volta che la gioca accumula
rimpianto verso quella che la domina), il che spiega il caso dei giochi risolvibili per
dominanza.
\end{extra}

% =====================================================================
\partelab{Equilibri correlati e apprendimento}
% =====================================================================
\fonte{Lesson4-NonCooperativeGameTheory\_with\_Learning\_in\_Games.ipynb}

Il notebook ripete la prima parte del laboratorio 3 (best response, equilibri, minimax,
equilibri $2\times2$: vedi lezione 3) e aggiunge tre esercizi: il programma lineare per gli
equilibri correlati, fictitious play e regret matching. Il notebook contiene gli output
della soluzione del docente: tutti i numeri sotto li riproducono.

\subsection{Equilibri correlati con \texttt{linprog}}

Le incognite sono le $m\cdot n$ probabilità $p(i,j)$, appiattite riga per riga. Ogni coppia
(azione raccomandata, deviazione) dà un vincolo. \texttt{linprog} vuole vincoli $\le 0$ e un
problema di minimo, quindi si cambia segno.

\begin{lstlisting}
import numpy as np
from scipy.optimize import linprog

def build_ce_lp(A1, A2):
    A1, A2 = np.asarray(A1, float), np.asarray(A2, float)
    m, n = A1.shape
    idx = lambda i, j: i * n + j
    rows = []
    for i in range(m):                    # giocatore 1: raccomandato i, devia a k
        for k in range(m):
            if k != i:
                r = np.zeros(m * n)
                for j in range(n):
                    r[idx(i, j)] = -(A1[i, j] - A1[k, j])
                rows.append(r)
    for j in range(n):                    # giocatore 2: raccomandato j, devia a l
        for l in range(n):
            if l != j:
                r = np.zeros(m * n)
                for i in range(m):
                    r[idx(i, j)] = -(A2[i, j] - A2[i, l])
                rows.append(r)
    A_ub, b_ub = np.array(rows), np.zeros(len(rows))
    A_eq, b_eq = np.ones((1, m * n)), np.array([1.0])
    c = -(A1 + A2).flatten()              # max benessere = min -benessere
    bounds = [(0, None)] * (m * n)
    return c, A_ub, b_ub, A_eq, b_eq, bounds, (m, n)

def solve_ce_max_social_welfare(A1, A2, method='highs'):
    c, A_ub, b_ub, A_eq, b_eq, bounds, (m, n) = build_ce_lp(A1, A2)
    res = linprog(c, A_ub=A_ub, b_ub=b_ub, A_eq=A_eq, b_eq=b_eq,
                  bounds=bounds, method=method)
    dist = res.x.reshape(m, n)
    return {"distribution": dist,
            "social_welfare": float((dist * (A1 + A2)).sum()),
            "linprog_result": res}
\end{lstlisting}

\paragraph{Risultati.}
\begin{itemize}
  \item \emph{Attraversamento}: benessere massimo $5{,}0$, distribuzione
        $\bigl(\begin{smallmatrix}0&1\\0&0\end{smallmatrix}\bigr)$, cioè l'equilibrio puro
        (Go, Stop). La matrice dei vincoli è identica a quella del notebook:
        \[
          A_{ub} = \begin{pmatrix}
            10 & -6 & 0 & 0\\ 0 & 0 & -10 & 6\\ 10 & 0 & -6 & 0\\ 0 & -10 & 0 & 6
          \end{pmatrix}
          \qquad\text{(variabili: } p_{GG}, p_{GS}, p_{SG}, p_{SS}\text{)}.
        \]
  \item \emph{Rigori}: benessere $0$, distribuzione uniforme $0{,}25$ su ogni cella.
\end{itemize}

\begin{osservazione}[il programma lineare sceglie un vertice]
Anche il semaforo ha benessere $5$ (verificato: è un equilibrio correlato), ma il risolutore
restituisce (Go, Stop): quando più soluzioni hanno lo stesso valore ottimo, ne restituisce
una qualsiasi, di solito un vertice del poliedro. Se si vuole la soluzione \emph{equa}
bisogna chiederla, con l'obiettivo maximin: risolvendo $\max t$ s.t.\
$t \le \Ev_p[u_i]$ si ottiene esattamente il semaforo,
$\bigl(\begin{smallmatrix}0&\frac12\\\frac12&0\end{smallmatrix}\bigr)$, con $2{,}5$ a testa.
\end{osservazione}

\begin{osservazione}[il commento sui rigori]
Il notebook motiva la distribuzione uniforme con ``il suo equilibrio di Nash è
$[0{,}5; 0{,}5]$''. Il motivo vero è più netto: nei rigori i quattro vincoli diventano
$p_{12} \le p_{11} \le p_{21} \le p_{22} \le p_{12}$, una catena chiusa che forza tutte le
probabilità a essere uguali. L'equilibrio correlato è \emph{unico}, e l'obiettivo non conta
(in un gioco a somma zero il benessere è $0$ per qualunque distribuzione).
\end{osservazione}

\subsection{Fictitious play}

\begin{lstlisting}
def fictitious_play(P, iterations=1000):
    p1_counts, p2_counts = np.ones(3), np.ones(3)   # 1 per evitare divisioni per 0
    sigma_p1, sigma_p2 = [], []
    for _ in range(iterations):
        p1_prob = p1_counts / p1_counts.sum()        # frequenze empiriche
        p2_prob = p2_counts / p2_counts.sum()
        a1 = np.argmax(P @ p2_prob)                  # BR alla stima sull'altro
        a2 = np.argmax(-P.T @ p1_prob)               # somma zero: payoff -P
        p1_counts[a1] += 1
        p2_counts[a2] += 1
        sigma_p1.append(p1_counts / p1_counts.sum())
        sigma_p2.append(p2_counts / p2_counts.sum())
    return p1_prob, p2_prob, np.array(sigma_p1), np.array(sigma_p2)
\end{lstlisting}

In sasso-carta-forbici le frequenze si avvicinano a $\tfrac13$, con oscillazioni che si
stringono:
\begin{center}
\renewcommand{\arraystretch}{1.2}
\begin{tabular}{rcc}
\toprule
Turni & Frequenze del giocatore 1 & Scarto massimo da $\tfrac13$ \\
\midrule
100    & $(0{,}359;\ 0{,}301;\ 0{,}340)$ & $0{,}032$ \\
1000   & $(0{,}347;\ 0{,}319;\ 0{,}334)$ & $0{,}014$ \\
10000  & $(0{,}332;\ 0{,}336;\ 0{,}333)$ & $0{,}002$ \\
\bottomrule
\end{tabular}
\end{center}

\begin{osservazione}[\texttt{argmax} e i pareggi]
\texttt{np.argmax} restituisce sempre il \emph{primo} massimo: è una regola di pareggio
deterministica. In sasso-carta-forbici la matrice è antisimmetrica ($-P^T = P$), quindi i due
giocatori calcolano la stessa best response a partire dagli stessi conteggi: giocano sempre
la stessa azione e hanno frequenze identiche. Non è un errore, ma è un artefatto da tenere presente se si usa fictitious play
per confronti nel progetto.
\end{osservazione}

\subsection{Regret matching}

\begin{lstlisting}
def regret_matching(P, T=10000, seed=0):
    rng = np.random.default_rng(seed)
    R1, R2 = np.zeros(3), np.zeros(3)          # rimpianti cumulati
    p1 = p2 = np.ones(3) / 3
    counts1 = np.zeros(3)
    for t in range(T):
        a1, a2 = rng.choice(3, p=p1), rng.choice(3, p=p2)
        counts1[a1] += 1
        R1 += P[:, a2] - P[a1, a2]             # quanto avrei preso con ogni azione
        R2 += -P[a1, :] + P[a1, a2]            # payoff della colonna = -P
        pos1, pos2 = np.maximum(R1, 0), np.maximum(R2, 0)
        p1 = pos1 / pos1.sum() if pos1.sum() > 0 else np.ones(3) / 3
        p2 = pos2 / pos2.sum() if pos2.sum() > 0 else np.ones(3) / 3
    return counts1 / T
\end{lstlisting}

Qui le azioni sono campionate e il rimpianto si calcola sull'azione effettivamente giocata
dall'avversario; le slide usano le strategie miste (utilità attese). Il rimpianto campionato
è una stima non distorta di quello atteso, quindi le due versioni si comportano allo stesso
modo in media.

\paragraph{Risultati.}
In sasso-carta-forbici, dopo 10\,000 turni, la frequenza del giocatore 1 è
$(0{,}335;\ 0{,}332;\ 0{,}333)$, con scarto massimo da $\tfrac13$ tra $0{,}002$ e $0{,}010$
su 5 semi. La versione con la formula delle slide dà scarti simili su questo gioco, ed è per
questo che il difetto non si vede nel notebook. Il confronto su un gioco asimmetrico
\[
  M = \begin{pmatrix} 2 & -1 & 1 \\ -1 & 1 & 0 \\ 1 & 0 & -1 \end{pmatrix},
  \qquad \text{minimax della riga } (0{,}333;\ 0{,}556;\ 0{,}111),\ \text{valore } 0{,}222,
\]
dopo 20\,000 turni su 3 semi:
\begin{center}
\renewcommand{\arraystretch}{1.2}
\begin{tabular}{lc}
\toprule
Versione & Frequenze della riga \\
\midrule
$\max(0, \sum_s r_s)$ (Hart, Mas-Colell) & $(0{,}34;\ 0{,}56;\ 0{,}10)$, $(0{,}34;\ 0{,}56;\ 0{,}11)$, $(0{,}35;\ 0{,}57;\ 0{,}08)$ \\
$\sum_s \max(0, r_s)$ (slide) & $(0{,}41;\ 0{,}35;\ 0{,}24)$, $(0{,}41;\ 0{,}35;\ 0{,}24)$, $(0{,}41;\ 0{,}35;\ 0{,}24)$ \\
\bottomrule
\end{tabular}
\end{center}
La seconda converge in modo stabile, ma al punto sbagliato: contro la seconda colonna la
strategia $(0{,}41;\ 0{,}35;\ 0{,}24)$ ottiene $-0{,}06$, mentre il minimax garantisce
$0{,}222$ contro qualunque colonna.

% ---------------------------------------------------------------------
\begin{riepilogo}
\begin{itemize}
  \item \textbf{Nash (1950)} via \textbf{Brouwer}: l'insieme dei profili misti è compatto
        e convesso, la funzione dei ``guadagni'' è continua e i suoi punti fissi sono
        esattamente gli equilibri.
  \item Trovare un equilibrio misto è \textbf{PPAD-completo} (End of Line, Sperner);
        trovare un equilibrio puro, o quello di benessere massimo, è \textbf{NP-difficile}
        (riduzione da 3-colorabilità con un nodo sfidante).
  \item \textbf{Equilibrio correlato}: una distribuzione sui profili da cui a nessuno
        conviene deviare data la propria raccomandazione. Ogni Nash è un CE; il semaforo è
        un CE che non è Nash. Si calcola con un programma lineare.
  \item Altri criteri: minimax (avversario ostile), minimo rimpianto massimo
        (avversario imprevedibile).
  \item \textbf{Fictitious play}: best response alle frequenze osservate; converge nei
        giochi a somma zero, risolvibili per dominanza, $2\times N$.
        \textbf{Regret matching}: si gioca in proporzione al rimpianto cumulato positivo;
        la distribuzione dei profili converge agli equilibri correlati (grossolani, per la
        versione con rimpianti incondizionati), e al minimax nei giochi a somma zero.
  \item \textbf{Lab}: CE con \texttt{linprog} (semaforo con l'obiettivo maximin),
        fictitious play e regret matching su RPS; la formula del rimpianto delle slide
        converge al punto sbagliato su un gioco asimmetrico.
\end{itemize}
\textbf{Prossima lezione}: giochi coalizionali (core, nucleolo, valore di Shapley).
\end{riepilogo}

% =====================================================================
%  Lezione 5 — Giochi coalizionali e concetti di soluzione
%  Materiale: Lesson_5_Coalitional_Game_Theory_I.pdf,
%             Lab_5-Coalitional_Game_Theory_Core_Shapley.pdf,
%             CoalitionalGames_Frozensets_blank.ipynb
% =====================================================================
\chapter{Giochi coalizionali}

\parteteoria

Il percorso della lezione:
\begin{enumerate}[label=\arabic*.]
  \item introduzione ai giochi coalizionali;
  \item come rappresentarli;
  \item quali coalizioni si formano;
  \item come dividere il payoff: \kw{stabilità} (core) e \kw{equità} (nucleolo, valore di
        Shapley).
\end{enumerate}

% ---------------------------------------------------------------------
\section{Dai giochi non cooperativi a quelli coalizionali}

\begin{itemize}
  \item I giochi non cooperativi modellano giocatori egoisti e in competizione.
  \item Giochi come il dilemma del prigioniero o l'attraversamento mostrano che la
        competizione può portare a esiti subottimi.
  \item Gli equilibri correlati modellano (e rendono possibile) la cooperazione, ma
        introducono un dispositivo di coordinamento, un \emph{mediatore}.
  \item Nella realtà, però, le persone formano spontaneamente \kw{coalizioni} e prendono
        \kw{accordi vincolanti}.
\end{itemize}

La situazione tipica che si modella con un gioco coalizionale è questa:
\begin{itemize}
  \item i giocatori possono trarre vantaggio dal formare coalizioni;
  \item a ogni coalizione è associato un \kw{valore};
  \item il valore totale va \kw{distribuito} tra i giocatori.
\end{itemize}

Servono due scelte di modellazione:
\begin{center}
\renewcommand{\arraystretch}{1.3}
\begin{tabularx}{\linewidth}{@{}l X X@{}}
\toprule
 & Domanda & Modello \\
\midrule
Coalizioni & il risultato di una coalizione è indipendente dalle altre? &
 sì $\to$ coalizioni singole; no $\to$ più coalizioni insieme (giochi in funzione di
 partizione) \\
Valore & l'utilità è trasferibile? & sì $\to$ un solo numero per coalizione (TU);
 no $\to$ un insieme di vettori di utilità (NTU) \\
\bottomrule
\end{tabularx}
\end{center}

\begin{definizione}[Utilità trasferibile e non trasferibile]
L'utilità è \kw{trasferibile} (TU) se è la stessa per tutti i membri della coalizione e si
può ridistribuire tra loro: i payoff sono monetari (o convertibili in denaro), i giocatori
possono farsi pagamenti laterali, e la ridistribuzione non cambia il valore totale.
Esempi: dividere profitti, ripartire costi, assegnare seggi in parlamento.

L'utilità è \kw{non trasferibile} (NTU) se ogni membro ha la propria e la coalizione non
può convertirla o ridistribuirla liberamente. Esempi: il prestigio di un articolo scientifico
scritto insieme, i benefici personali di una collaborazione, i vantaggi legati a indicatori
di prestazione individuali.
\end{definizione}

\begin{attenzione}[denaro non vuol dire utilità trasferibile]
Il notebook lo dice bene: poter trasferire \emph{denaro} non implica che sia trasferibile
l'\emph{utilità}. Un giocatore ricco e uno povero possono trarre utilità molto diverse dalla
stessa somma (lezione 2: la funzione di Bernoulli è concava). Assumere TU significa assumere
che tutti valutino il denaro allo stesso modo, cioè utilità lineari nel denaro.
\end{attenzione}

\section{Giochi in funzione caratteristica}

\begin{definizione}[Gioco in funzione caratteristica]
Un gioco in funzione caratteristica è una coppia $(N, v)$ dove:
\begin{itemize}
  \item $N = \{1,\dots,n\}$ è l'insieme dei giocatori, detto anche \kw{grande coalizione};
  \item $v : 2^N \to \R$ è la \kw{funzione caratteristica}: per ogni sottoinsieme $C$ di
        giocatori, $v(C)$ è quanto possono guadagnare i membri di $C$ lavorando insieme.
        L'utilità $v(C)$ è trasferibile.
\end{itemize}
Di solito si assume che $v$ sia
\begin{itemize}
  \item \kw{normalizzata}: $v(\emptyset) = 0$;
  \item \kw{non negativa}: $v(C) \ge 0$ per ogni $C \subseteq N$;
  \item \kw{monotona}: $v(C) \le v(D)$ se $C \subseteq D$.
\end{itemize}
\end{definizione}

\begin{esempio}[il gelato]
Tre persone vogliono comprare del gelato e possono mettere insieme i soldi.
\begin{center}
\renewcommand{\arraystretch}{1.2}
\begin{tabular}{lc@{\qquad\qquad}lc}
\toprule
Giocatore & Soldi & Vaschetta & Prezzo \\
\midrule
A & \$6 & 500 g & \$7 \\
B & \$3 & 750 g & \$9 \\
C & \$3 & 1000 g & \$11 \\
\bottomrule
\end{tabular}
\end{center}
Il valore di una coalizione è la quantità massima di gelato che può comprare:
\[
\begin{aligned}
  v(\emptyset) &= v(\{A\}) = v(\{B\}) = v(\{C\}) = 0
    &&\text{(nessuno ha \$7 da solo)}\\
  v(\{A,B\}) &= v(\{A,C\}) = 750
    &&\text{(\$9: la vaschetta da 750 g)}\\
  v(\{B,C\}) &= 0 &&\text{(\$6: non basta per nulla)}\\
  v(\{A,B,C\}) &= 1000 &&\text{(\$12: la vaschetta da 1000 g).}
\end{aligned}
\]
\end{esempio}

\paragraph{Oltre la funzione caratteristica.}
I giochi in funzione caratteristica modellano utilità trasferibile e risultati di una
coalizione indipendenti da ciò che fanno gli altri. Se i risultati dipendono dalle altre
coalizioni si usano i \kw{giochi in funzione di partizione}, con valori $v(C, P)$ dove $C$
è una coalizione della partizione $P$; se l'utilità non è trasferibile si usano i
\kw{giochi NTU}, dove a ogni coalizione è associato un \emph{insieme} di vettori di utilità.
Ogni gioco TU si può rappresentare come un gioco NTU con un continuo di azioni, quindi
\[
  \text{NTU} \supset \text{TU} \supset \text{giochi in funzione caratteristica};
\]
nella pratica si usano soprattutto questi ultimi, che spesso si chiamano semplicemente
``giochi TU''.

\paragraph{Le due domande.}
I giochi coalizionali servono a rispondere a due domande:
\begin{enumerate}
  \item quali coalizioni si formeranno?
  \item come si divide il payoff all'interno di una coalizione?
\end{enumerate}
Nel gioco del gelato la coalizione con il benessere più alto è $\{A,B,C\}$. Ma dividere in
parti uguali non è equo: $A$ è più decisivo di $B$ e $C$, perché senza di lui non si compra
nulla.

% ---------------------------------------------------------------------
\section{Struttura di coalizione}

\begin{definizione}[Generazione della struttura di coalizione]
Dato un gioco $G = (N,v)$ si cerca la \kw{struttura di coalizione} $CS(G)$, cioè una
\emph{partizione} di $N$ in coalizioni disgiunte, che massimizza il benessere sociale
$\sum_{C \in CS(G)} v(C)$.
\end{definizione}

\begin{attenzione}[partizione, non ricoprimento]
Le slide dicono che una struttura di coalizione è un ``ricoprimento'' (\en{cover}) di $N$.
Un ricoprimento può avere insiemi sovrapposti, e un giocatore che sta in due coalizioni
verrebbe contato due volte. Quello che si intende, e che le slide stesse scrivono più
avanti, è una \kw{partizione}: le coalizioni coprono $N$ e sono a due a due disgiunte.
\end{attenzione}

\begin{esempio}[quando non conviene la grande coalizione]
\[
\begin{aligned}
  v(\emptyset) &= 0, & v(\{A\}) &= 0, & v(\{B\}) &= 300, & v(\{C\}) &= 0,\\
  v(\{A,B\}) &= 750, & v(\{A,C\}) &= 750, & v(\{B,C\}) &= 400, & v(\{A,B,C\}) &= 1000.
\end{aligned}
\]
La struttura di coalizione che massimizza il benessere è $\{\{A,C\}, \{B\}\}$, con
$750 + 300 = 1050 > 1000$: qui la grande coalizione \emph{non} è la scelta migliore.
\end{esempio}

\subsection{Il problema è difficile}

Le partizioni di un insieme di $n$ elementi sono il \kw{numero di Bell} $B_n$, che soddisfa
\[
  \left(\frac{n}{4}\right)^{n/2} \le B_n < n^n .
\]
Enumerarle tutte è impraticabile. Gli algoritmi usati sono la programmazione dinamica, il
grafo delle partizioni e le euristiche per giochi strutturati.

\subsection{Programmazione dinamica}

Sia $v^*(S)$ il massimo benessere ottenibile con una struttura di coalizione dei giocatori
di $S$, e $CS^*(S)$ la struttura che lo realizza: $v(CS^*(S)) = v^*(S)$.

\begin{teorema}[Equazione della programmazione dinamica]
Per ogni $C \subseteq N$:
\[
  v^*(C) = \max\Bigl\{ v(C),\ \max_{\substack{C' \cup C'' = C\\ C' \cap C'' = \emptyset}}
            \bigl(v^*(C') + v^*(C'')\bigr) \Bigr\}.
\]
\end{teorema}
O la coalizione $C$ resta intera, o la si divide in due parti, ciascuna delle quali va poi
organizzata nel modo migliore.

\paragraph{Dimostrazione.}
($\ge$) $v^*(C) \ge v(C)$, perché $\{C\}$ è una struttura di coalizione di $C$; e
$v^*(C) \ge v^*(C') + v^*(C'')$, perché unendo le strutture ottime delle due parti si ottiene
una struttura di $C$.
($\le$) Se $CS^*(C) = \{C\}$ allora $v^*(C) = v(C)$. Altrimenti sia $C'$ una coalizione di
$CS^*(C)$ e $C'' = C \setminus C'$. Il resto, $CS^*(C) \setminus \{C'\}$, è una struttura
di coalizione di $C''$, quindi $v^*(C'') \ge v^*(C) - v(C')$; inoltre
$v^*(C') \ge v(C')$. Sommando, $v^*(C') + v^*(C'') \ge v^*(C)$. $\square$

\begin{esempio}[la tabella della programmazione dinamica]
Con la funzione caratteristica dell'esempio precedente, procedendo per dimensione crescente:
\begin{center}
\renewcommand{\arraystretch}{1.25}
\begin{tabular}{lll}
\toprule
$C$ & $v^*(C)$ & $CS^*(C)$ \\
\midrule
$\{A\}$ & $0$ & $\{A\}$ \\
$\{B\}$ & $300$ & $\{B\}$ \\
$\{C\}$ & $0$ & $\{C\}$ \\
$\{A,B\}$ & $\max(v^*(A)+v^*(B),\ v(AB)) = \max(300, 750) = 750$ & $\{A,B\}$ \\
$\{A,C\}$ & $\max(0, 750) = 750$ & $\{A,C\}$ \\
$\{B,C\}$ & $\max(300, 400) = 400$ & $\{B,C\}$ \\
$\{A,B,C\}$ & $\max\bigl(1000,\ 750 + 0,\ 750 + 300,\ 400 + 0\bigr) = 1050$ & $\{\{A,C\},\{B\}\}$ \\
\bottomrule
\end{tabular}
\end{center}
Per $\{A,B,C\}$ le divisioni in due parti sono tre: $\{A,B\}|\{C\}$, $\{A,C\}|\{B\}$,
$\{B,C\}|\{A\}$. La migliore è la seconda.
\end{esempio}

\begin{attenzione}[un errore nella tabella delle slide]
Le slide scrivono $v^*(\{B,C\}) = \max(300, 400) = 750$ e riportano 750 in tabella. Il
massimo è 400. L'errore non cambia il risultato finale (1050 viene dalla divisione
$\{A,C\}|\{B\}$, che non usa $v^*(\{B,C\})$), ma nel passo successivo le slide elencano due
volte la divisione $\{A,B\}|\{C\}$ e mai $\{B,C\}|\{A\}$, proprio quella che userebbe il
valore sbagliato. Il codice del laboratorio conferma la tabella corretta.
\end{attenzione}

L'approccio è \kw{esatto}, ma il tempo è $O(3^n)$: non è praticabile per $n \ge 20$.

\begin{extra}[da dove viene $3^n$]
Si considerano tutte le coppie (coalizione $C$, sua parte $C'$). Ogni giocatore, in una
coppia, può stare in $C'$, in $C \setminus C'$ o fuori da $C$: tre possibilità per ciascuno
degli $n$ giocatori, cioè $3^n$ coppie. Con $n = 20$ sono circa $3{,}5$ miliardi.
\end{extra}

\paragraph{Grafo delle partizioni.}
Le partizioni si possono disporre su livelli, dal livello 1 (la sola grande coalizione) al
livello $n$ (tutti singoli); due partizioni sono collegate se una si ottiene dall'altra
fondendo due coalizioni. Le slide mostrano il grafo per 4 giocatori senza svilupparlo; l'idea
degli algoritmi che lo usano è esplorare per primi pochi livelli scelti bene e fermarsi
quando si ha una garanzia sufficiente sulla distanza dall'ottimo.

\paragraph{In pratica.}
Si usano euristiche e rappresentazioni compatte della funzione caratteristica: funzioni
caratteristiche analitiche, algoritmi di clustering (formazione delle coalizioni per
similarità o distanza, partizionamento di grafi), metaeuristiche (algoritmi genetici,
\en{simulated annealing}, \ldots).

% ---------------------------------------------------------------------
\section{Dividere il payoff}

\begin{definizione}[Esito di un gioco TU]
Un \kw{esito} di un gioco TU $G = (N,v)$ è una coppia $(CS, x)$ dove:
\begin{itemize}
  \item $CS = (C_1, \dots, C_k)$ è una struttura di coalizione, cioè una partizione di $N$:
        $\bigcup_i C_i = N$, $C_i \cap C_j = \emptyset$ per $i \ne j$;
  \item $x = (x_1,\dots,x_n)$ è un \kw{vettore di payoff} che distribuisce il valore di ogni
        coalizione di $CS$: $\sum_{i \in C} x_i = v(C)$ per ogni $C \in CS$
        (\kw{efficienza}).
\end{itemize}
Un esito è un'\kw{imputazione} se soddisfa la \kw{razionalità individuale}:
$x_i \ge v(\{i\})$ per ogni $i \in N$.
\end{definizione}

\begin{esempio}[esiti]
Con $v(\{1,2,3\}) = 9$ e $v(\{4,5\}) = 4$:
\begin{itemize}
  \item $\bigl((\{1,2,3\},\{4,5\}),\ (3,3,3,3,1)\bigr)$ è un esito: $3+3+3 = 9$, $3+1 = 4$;
  \item $\bigl((\{1,2,3\},\{4,5\}),\ (2,3,2,3,3)\bigr)$ \emph{non} lo è: la prima
        coalizione distribuisce $7$ e la seconda $6$. Non sono ammessi trasferimenti
        \emph{tra} coalizioni.
\end{itemize}
\end{esempio}

\subsection{Superadditività}

\begin{definizione}[Gioco superadditivo]
Un gioco $(N,v)$ è \kw{superadditivo} se
\[
  v(C_1 \cup C_2) \ge v(C_1) + v(C_2)
  \qquad \forall C_1, C_2 \subset N \text{ con } C_1 \cap C_2 = \emptyset .
\]
\end{definizione}

Il gioco del gelato è superadditivo: per esempio $v(\{A,B\}) = 750 \ge v(\{A\}) + v(\{B\}) = 0$,
$v(\{B,C\}) = 0 \ge 0 + 0$, $v(\{A,B,C\}) = 1000 \ge 750 + 0$.

Nei giochi superadditivi due coalizioni possono sempre fondersi senza perdere nulla: si può
assumere che si formi la \kw{grande coalizione}. Non c'è il problema della struttura di
coalizione e la notazione si semplifica: un esito è solo un vettore
$x = (x_1,\dots,x_n)$ con $\sum_{i\in N} x_i = v(N)$.

\begin{osservazione}[un'avvertenza delle slide]
Molti articoli definiscono gli esiti in questo modo, come semplici vettori con
$\sum x_i = v(N)$, \emph{anche} quando il gioco non è superadditivo. Va tenuto presente
quando si leggono: nel gioco della programmazione dinamica, per esempio, dividere $v(N)=1000$
sarebbe sbagliato, perché la struttura migliore produce $1050$.
\end{osservazione}

Da qui in poi si considerano solo giochi TU superadditivi.

% ---------------------------------------------------------------------
\section{Stabilità: il core}

Nel gioco del gelato si deve dividere $1000$. Se la divisione è $(200, 200, 600)$, $A$ e $B$
possono comprare da soli la vaschetta da 750 g e dividersela a metà, $375$ a testa: meglio
di $200$. L'esito $(200,200,600)$ \kw{non è stabile}.

\begin{definizione}[Core]
Il \kw{core} di un gioco è l'insieme degli esiti stabili, quelli da cui nessuna coalizione
vuole sganciarsi:
\[
  \core(G) = \Bigl\{ x \ \Big|\ x_i \ge 0\ \forall i \in N,\
  \sum_{i \in N} x_i = v(N),\
  \sum_{i \in C} x_i \ge v(C)\ \forall C \subseteq N \Bigr\}.
\]
Ogni coalizione riceve almeno quanto guadagnerebbe da sola.
\end{definizione}

\begin{esempio}[esiti nel core]
\begin{itemize}
  \item Con $v(\{1,2,3\}) = 9$, $v(\{4,5\}) = 4$ e $v(\{2,4\}) = 7$, l'esito
        $(3,3,3,3,1)$ \emph{non} è nel core: i giocatori 2 e 4 ricevono $3 + 3 = 6 < 7$ e
        preferiscono formare la coalizione $\{2,4\}$.
  \item Nel gelato, $(200,200,600)$ non è nel core perché
        $v(\{A,B\}) = 750 > x_A + x_B = 400$.
  \item Nel gelato, $(500,250,250)$ \emph{è} nel core: $x_A + x_B = 750 \ge 750$,
        $x_A + x_C = 750 \ge 750$, $x_B + x_C = 500 \ge 0$. Nessuna deviazione dà di più.
\end{itemize}
\end{esempio}

\begin{osservazione}[il primo esempio]
L'esempio con $(3,3,3,3,1)$ usa una struttura di coalizione $\{\{1,2,3\},\{4,5\}\}$, mentre la
definizione di core appena data è quella per la grande coalizione. Il concetto è lo stesso
(nessuna coalizione deve poter fare meglio da sola), ma nella versione con strutture di
coalizione l'efficienza si chiede coalizione per coalizione.
\end{osservazione}

\begin{extra}[com'è fatto il core del gelato]
I vincoli che contano sono due. Da $x_A + x_B \ge 750$ e $x_A + x_B + x_C = 1000$ segue
$x_C \le 250$; analogamente $x_B \le 250$. Il core è quindi
\[
  \core = \{ x \ge 0 : x_B \le 250,\ x_C \le 250,\ x_A = 1000 - x_B - x_C \},
\]
e in particolare $x_A \ge 500$ in ogni esito stabile. Si vede bene il potere di $A$: se $B$
chiedesse più di 250, $A$ potrebbe accordarsi con $C$ e dividere 750 dando a entrambi più di
quanto ricevono; lo stesso vale a parti invertite. $B$ e $C$ sono in concorrenza tra loro
per essere il partner di $A$.
\end{extra}

\subsection{Due cattive notizie}

\begin{esempio}[il core può essere vuoto]
\[
  v(\{1\}) = v(\{2\}) = v(\{3\}) = 0,\quad
  v(\{1,2\}) = 7,\ v(\{1,3\}) = 6,\ v(\{2,3\}) = 5,\quad
  v(\{1,2,3\}) = 8.
\]
Per stare nel core serve stabilità, $x_1 + x_2 \ge 7$, $x_1 + x_3 \ge 6$, $x_2 + x_3 \ge 5$,
ed efficienza, $x_1 + x_2 + x_3 = 8$. Sommando le tre disuguaglianze:
$2(x_1+x_2+x_3) \ge 18$, cioè $x_1+x_2+x_3 \ge 9 > 8$. Impossibile: il core è vuoto. E il
gioco è superadditivo, quindi la superadditività non basta a garantire un core non vuoto.
\end{esempio}

\begin{esempio}[il core può contenere esiti iniqui]
Nel gelato $(1000, 0, 0)$ è nel core: $x_A + x_B = 1000 \ge 750$, $x_A + x_C = 1000 \ge 750$,
$x_B + x_C = 0 \ge 0$, e la somma è $1000$. $B$ e $C$ mettono metà dei soldi e non ricevono
nulla, ma nessuna coalizione può fare meglio da sola. \kw{Stabilità non vuol dire equità.}
\end{esempio}

% ---------------------------------------------------------------------
\section{Il nucleolo}

\begin{definizione}[Eccesso]
L'\kw{eccesso} di una coalizione $S$ rispetto all'esito $x$ è
\[
  e(S, x) = v(S) - \sum_{i \in S} x_i .
\]
Misura l'\kw{insoddisfazione} della coalizione: quanto potrebbe guadagnare da sola rispetto
a quello che riceve. Un esito è nel core se e solo se tutti gli eccessi sono $\le 0$.
\end{definizione}

\begin{esempio}[eccessi]
Con $v(\{i\}) = 0$, $v(\{i,j\}) = 1$ per ogni coppia e $v(\{1,2,3\}) = 3$:
\begin{itemize}
  \item $x = (0,0,3)$: $e(\{1,2\},x) = 1 - 0 = 1$, la coppia $\{1,2\}$ è insoddisfatta;
  \item $x = (1,2,0)$: $e(\{1,2\},x) = 1 - 3 = -2$, la coppia è più che soddisfatta.
\end{itemize}
\end{esempio}

\begin{definizione}[Nucleolo]
Sia $\theta(x)$ il \kw{vettore degli eccessi} di tutte le coalizioni (diverse da $\emptyset$
e da $N$), ordinato dal più grande al più piccolo. Il \kw{nucleolo} è l'imputazione
\[
  \eta \in X(G) \ :\ \nexists\, y \in X(G) \text{ con } \theta(y) \prec_{lex} \theta(\eta),
\]
dove $X(G)$ è l'insieme delle imputazioni e $\prec_{lex}$ l'ordine lessicografico.
\end{definizione}
In parole: è l'esito che rende \kw{lessicograficamente minimo} il vettore delle
insoddisfazioni. Prima si rende la coalizione più insoddisfatta il meno insoddisfatta
possibile; a parità, la seconda; e così via.

\begin{esempio}[nucleolo del gioco simmetrico]
Nel gioco precedente:
\begin{itemize}
  \item $x = (1,2,0)$: $\theta(x) = (0, 0, -1, -1, -2, -2)$;
  \item $\eta = (1,1,1)$: $\theta(\eta) = (-1,-1,-1,-1,-1,-1)$.
\end{itemize}
Il primo elemento di $\theta(\eta)$ è $-1 < 0$: $\eta$ è lessicograficamente migliore, ed è
il nucleolo (lo conferma il calcolo con l'algoritmo seguente).
\end{esempio}

\begin{extra}[il nucleolo esiste sempre ed è unico]
Il nucleolo esiste ed è unico per ogni gioco con almeno un'imputazione. E se il core non è
vuoto il nucleolo sta nel core: un esito del core ha tutti gli eccessi $\le 0$, quindi il
nucleolo, che minimizza il più grande, ne ha anch'esso il più grande $\le 0$. È per questo
che lo si considera la versione ``più centrale'' del core.
\end{extra}

\subsection{L'algoritmo di Kopelowitz}

Il vettore degli eccessi si minimizza lessicograficamente con una \emph{sequenza} di
programmi lineari. Il primo minimizza l'eccesso massimo $\epsilon_1$:
\[
LP_1:\quad
\begin{aligned}
  \min\ & \epsilon_1\\
  \text{s.t. }\ & v(S) - \textstyle\sum_{i \in S} x_i \le \epsilon_1
    && \forall S \subset N,\ S \notin W_0 = \{\emptyset\}\\
  & x_i \ge v(\{i\}),\quad \textstyle\sum_{i \in N} x_i = v(N).
\end{aligned}
\]
Le coalizioni che all'ottimo hanno eccesso esattamente $\epsilon_1^*$ formano l'insieme
$W_1$ e il loro eccesso viene \emph{fissato}. Il programma $k$-esimo minimizza il massimo
eccesso delle coalizioni rimaste:
\[
LP_k:\quad
\begin{aligned}
  \min\ & \epsilon_k\\
  \text{s.t. }\ & v(S) - \textstyle\sum_{i \in S} x_i = \epsilon_r^*
    && \forall S \in W_r,\ r = 1,\dots,k-1\\
  & v(S) - \textstyle\sum_{i \in S} x_i \le \epsilon_k
    && \forall S \notin W_0 \cup W_1 \cup \dots \cup W_{k-1}\\
  & x_i \ge v(\{i\}),\quad \textstyle\sum_{i \in N} x_i = v(N),
\end{aligned}
\]
con $W_r = \{S \mid e(S,x) = \epsilon_r^*\}$. Ci si ferma quando tutte le coalizioni hanno
l'eccesso fissato.

\begin{itemize}
  \item Ogni programma ha $O(2^n)$ vincoli, uno per coalizione.
  \item Nel caso peggiore si fissa una coalizione per volta: fino a $2^n$ iterazioni.
  \item In generale è \kw{molto poco pratico}.
\end{itemize}

\begin{attenzione}[quali coalizioni fissare]
``Le coalizioni con eccesso pari a $\epsilon_k^*$'' va inteso come ``le coalizioni il cui
eccesso vale $\epsilon_k^*$ in \emph{tutte} le soluzioni ottime del programma $k$''. Una
soluzione ottima qualsiasi può avere per caso altre coalizioni con eccesso
$\epsilon_k^*$ che in un'altra soluzione ottima potrebbero scendere: fissarle porterebbe a un
risultato sbagliato. Nell'implementazione del laboratorio lo si controlla con un programma
lineare in più per ciascuna coalizione candidata (si prova ad abbassarne l'eccesso restando
all'ottimo).
\end{attenzione}

\subsection{Il nucleolo può essere iniquo}

\begin{esempio}[un gioco in cui conta solo A]
\begin{gather*}
  v(\emptyset) = v(\{A\}) = v(\{B\}) = v(\{C\}) = 0,\\
  v(\{A,B\}) = v(\{A,C\}) = 750,\quad v(\{B,C\}) = 0,\quad v(\{A,B,C\}) = 750.
\end{gather*}
Il nucleolo è $\eta = (750, 0, 0)$: assegna tutto ad $A$, anche se almeno uno tra $B$ e $C$
è necessario per ottenere qualcosa. Anche qui stabilità non implica equità.
\end{esempio}

\begin{attenzione}[il vettore degli eccessi nelle slide]
Le slide riportano $\theta(\eta) = (0,0,0,0,0,0)$. Gli eccessi di $\eta = (750,0,0)$ sono:
$e(\{B\}) = e(\{C\}) = 0$, $e(\{A,B\}) = e(\{A,C\}) = 750 - 750 = 0$, $e(\{B,C\}) = 0$, ma
$e(\{A\}) = 0 - 750 = -750$. Quindi $\theta(\eta) = (0,0,0,0,0,-750)$. Il nucleolo resta
quello: il core di questo gioco contiene \emph{solo} $(750,0,0)$ (da $x_A + x_B \ge 750$ e
$x_A + x_B + x_C = 750$ segue $x_C = 0$, e lo stesso per $x_B$), e il nucleolo sta nel core.
\end{attenzione}

% ---------------------------------------------------------------------
\section{Equità: il valore di Shapley}

\textbf{Idea}: un esito equo dovrebbe ricompensare ogni giocatore in proporzione al
\kw{contributo} che porta agli altri. Se il giocatore $i$ entra nella coalizione $C$, il suo
\kw{contributo marginale} è il surplus che crea:
\[
  v(C \cup \{i\}) - v(C).
\]
Nel gioco precedente, per esempio, $v(\{A,B\}) - v(\{A\}) = 750$ e
$v(\{A,B\}) - v(\{B\}) = 750$: il contributo di un giocatore dipende da chi c'è già. Un
esito equo si ottiene premiando ciascuno con il suo contributo \emph{medio}.

\begin{definizione}[Valore di Shapley]
Il valore di Shapley considera tutti gli ordini in cui la grande coalizione si può formare e
dà a ciascuno la stessa importanza:
\[
  \phi(i, v) = \frac{1}{|N|!} \sum_{\pi \in \Pi_N}
     \Bigl( v\bigl(B(\pi, i) \cup \{i\}\bigr) - v\bigl(B(\pi, i)\bigr) \Bigr),
\]
dove $\Pi_N$ è l'insieme delle permutazioni di $N$ e $B(\pi,i)$ l'insieme dei giocatori che
precedono $i$ nella permutazione $\pi$.
\end{definizione}

\begin{esempio}[Shapley nel gioco del gelato]
Contributi marginali di $A$ nei sei ordini:
\begin{center}
\renewcommand{\arraystretch}{1.2}
\begin{tabular}{lll}
\toprule
Ordine & Contributo di $A$ & \\
\midrule
$A,B,C$ & $v(\{A\}) - v(\emptyset)$ & $= 0$ \\
$A,C,B$ & $v(\{A\}) - v(\emptyset)$ & $= 0$ \\
$B,A,C$ & $v(\{A,B\}) - v(\{B\})$ & $= 750$ \\
$B,C,A$ & $v(\{A,B,C\}) - v(\{B,C\})$ & $= 1000$ \\
$C,A,B$ & $v(\{A,C\}) - v(\{C\})$ & $= 750$ \\
$C,B,A$ & $v(\{A,B,C\}) - v(\{B,C\})$ & $= 1000$ \\
\bottomrule
\end{tabular}
\end{center}
\[
  \phi(A) = \tfrac16(0 + 0 + 750 + 1000 + 750 + 1000) = 583{,}\overline{3}.
\]
Per $B$ i contributi sono $750, 250, 0, 0, 250, 0$, quindi
$\phi(B) = \tfrac{1250}{6} = 208{,}\overline{3}$, e per simmetria
$\phi(C) = 208{,}\overline{3}$. Somma: $1000$.
\end{esempio}

\begin{esempio}[Shapley nel gioco in cui conta solo A]
\begin{center}
\renewcommand{\arraystretch}{1.2}
\begin{tabular}{lccc}
\toprule
Ordine & contributo di $A$ & contributo di $B$ & contributo di $C$ \\
\midrule
$A,B,C$ & 0   & 750 & 0   \\
$A,C,B$ & 0   & 0   & 750 \\
$B,A,C$ & 750 & 0   & 0   \\
$B,C,A$ & 750 & 0   & 0   \\
$C,A,B$ & 750 & 0   & 0   \\
$C,B,A$ & 750 & 0   & 0   \\
\midrule
Valore di Shapley & $500$ & $125$ & $125$ \\
\bottomrule
\end{tabular}
\end{center}
A differenza del nucleolo, $B$ e $C$ ricevono qualcosa: chi arriva subito dopo $A$ è quello
che ``sblocca'' il valore.
\end{esempio}

\begin{extra}[equo, ma non stabile]
$(500, 125, 125)$ non è nel core di questo gioco: il core contiene solo $(750,0,0)$, e infatti
$x_A + x_B = 625 < 750 = v(\{A,B\})$. $A$ e $B$ da soli prenderebbero di più. Equità e
stabilità sono richieste diverse e possono essere incompatibili. Nel gioco del gelato,
invece, il valore di Shapley $(583{,}3;\ 208{,}3;\ 208{,}3)$ \emph{è} nel core (rispetta
$x_B, x_C \le 250$). Il risultato è coerente con l'esempio delle tre imprese della lezione 1,
dove con valore 100 si otteneva $(66{,}67;\ 16{,}67;\ 16{,}67)$: è lo stesso gioco in scala.
\end{extra}

\subsection{Caratterizzazione assiomatica}

\begin{teorema}[Shapley]
Il valore di Shapley è l'\kw{unico} schema di ripartizione che soddisfa:
\begin{enumerate}
  \item \kw{efficienza}: $\phi_1 + \phi_2 + \dots + \phi_n = v(N)$;
  \item \kw{giocatore nullo}: se $v(C \cup \{i\}) = v(C)$ per ogni $C$, allora $\phi_i = 0$;
  \item \kw{simmetria}: se $i$ e $j$ sono intercambiabili
        ($v(C \cup \{i\}) = v(C \cup \{j\})$ per ogni $C$ che non li contiene), allora
        $\phi_i = \phi_j$;
  \item \kw{additività}: $\phi_i(G_1 + G_2) = \phi_i(G_1) + \phi_i(G_2)$.
\end{enumerate}
Inoltre si può calcolare con la formula sui sottoinsiemi:
\[
  \phi(i, v) = \sum_{\substack{C \subseteq N\\ i \in C}}
     \frac{(|N| - |C|)!\,(|C| - 1)!}{|N|!}\,\bigl( v(C) - v(C \setminus \{i\}) \bigr).
\]
\end{teorema}

\begin{osservazione}[da dove viene il coefficiente]
Negli ordini in cui $i$ entra trovando già esattamente la coalizione $C \setminus \{i\}$,
i $|C|-1$ predecessori possono stare in $(|C|-1)!$ ordini e i $|N|-|C|$ successori in
$(|N|-|C|)!$. Il coefficiente è la frazione di permutazioni in cui $i$ dà proprio il
contributo $v(C) - v(C\setminus\{i\})$. Si passa da $n!$ permutazioni a $2^{n-1}$
coalizioni contenenti $i$: è la formula da usare nel codice.
\end{osservazione}

\begin{attenzione}[``giocatore fittizio'' e ``giocatore nullo'']
Le slide chiamano l'assioma 2 \emph{dummy player} e chiedono $\phi_i = 0$. In molti testi
il \emph{dummy} è un giocatore che aggiunge sempre esattamente $v(\{i\})$, e riceve
$\phi_i = v(\{i\})$; quello che riceve zero è il \emph{null player}, che non aggiunge mai
nulla. Con $v(\{i\}) = 0$ le due nozioni coincidono. Le slide scrivono anche ``Shapely''
invece di Shapley.
\end{attenzione}

% =====================================================================
\partelab{Core e valore di Shapley}
% =====================================================================
\fonte{CoalitionalGames\_Frozensets\_blank.ipynb, Lab\_5-Coalitional\_Game\_Theory\_Core\_Shapley.pdf}

\textbf{Obiettivo}: rappresentare un gioco coalizionale in Python, verificare la
superadditività, controllare se un esito sta nel core e calcolare il valore di Shapley.
Il notebook contiene gli output attesi; tutti sono riprodotti. In più sono stati implementati
e verificati la programmazione dinamica per la struttura di coalizione e il nucleolo.

\subsection{Rappresentare la funzione caratteristica}

Si usa un dizionario indicizzato da \texttt{frozenset}: i \texttt{frozenset} sono insiemi
immutabili, quindi si possono usare come chiavi, e hanno tutte le operazioni degli insiemi
(appartenenza, inclusione, unione, intersezione).

\begin{lstlisting}
from itertools import combinations
from math import factorial

v = {
    frozenset(['A']): 0, frozenset(['B']): 0, frozenset(['C']): 0,
    frozenset(['A', 'B']): 750, frozenset(['A', 'C']): 750,
    frozenset(['B', 'C']): 0,
    frozenset(['A', 'B', 'C']): 1000,
}

def powerset(players):   # tutte le coalizioni NON vuote: 2^n - 1
    return [frozenset(c) for k in range(1, len(players) + 1)
            for c in combinations(players, k)]
\end{lstlisting}

\begin{attenzione}[la coalizione vuota non c'è]
Il dizionario del notebook non contiene $\emptyset$, e la \texttt{powerset} fornita lo
esclude (il commento dice ``$2^{n-1}$ frozen sets'', ma sono $2^n - 1$). Ogni formula che
usa $v(\emptyset)$, come la formula di Shapley con $C = \{i\}$, deve trattarlo a parte:
\texttt{v.get(S, 0)} invece di \texttt{v[S]}, altrimenti un \texttt{KeyError}.
\end{attenzione}

\subsection{Superadditività}

\begin{lstlisting}
def is_superadditive(v):
    for C1 in v:
        for C2 in v:
            if not (C1 & C2):                       # disgiunte
                if v.get(C1 | C2, 0) < v[C1] + v[C2]:
                    return False
    return True
\end{lstlisting}
Il gioco del gelato dà \texttt{True}; il gioco del notebook con $v(A) = 6$, $v(B) = 12$,
$v(A,B) = 12$ dà \texttt{False} (perché $12 < 6 + 12$). Entrambi come atteso.

\subsection{Core}

\begin{lstlisting}
def is_stable(x, v, tol=1e-9):
    return all(sum(x[i] for i in C) >= v[C] - tol for C in v)

def is_in_the_core(x, v, tol=1e-9):
    N = max(v)                                       # la grande coalizione
    return (all(xi >= -tol for xi in x.values())
            and abs(sum(x.values()) - v[N]) <= tol
            and is_stable(x, v, tol))
\end{lstlisting}

\begin{center}
\renewcommand{\arraystretch}{1.2}
\begin{tabular}{lcl}
\toprule
Esito & Nel core? & Perché \\
\midrule
$(333{,}3334;\ 333{,}3333;\ 333{,}3333)$ & no & $x_A + x_B \approx 666{,}7 < 750$ \\
$(500;\ 250;\ 250)$ & sì & tutti i vincoli rispettati \\
$(1000;\ 0;\ 0)$ & sì & stabile, anche se iniquo \\
\bottomrule
\end{tabular}
\end{center}

\begin{attenzione}[\texttt{max(v)} trova la grande coalizione, ma per un motivo sottile]
Il notebook ricava la grande coalizione con \texttt{max(characteristic\_function)}. Funziona,
ma non perché confronti le dimensioni: per i \texttt{frozenset} l'operatore \texttt{>}
significa ``sovrainsieme proprio''. \texttt{max} tiene il primo elemento e lo sostituisce
con ogni elemento che lo contiene strettamente; la grande coalizione contiene tutte le altre,
quindi appena la incontra non la lascia più. Se la grande coalizione mancasse dal dizionario,
\texttt{max} restituirebbe una coalizione qualsiasi senza segnalare nulla.
\end{attenzione}

\begin{osservazione}[perché non si calcola il core]
Il notebook si limita a \emph{verificare} se un esito sta nel core. Calcolarlo tutto
richiederebbe di trovare i vertici del poliedro definito dai vincoli, con una sequenza di
programmi lineari. Per sapere solo se il core è \emph{vuoto} basta però un programma
lineare: si cerca un $x$ qualsiasi che soddisfi i vincoli. Applicato al gioco con
$v(\{1,2\}) = 7$, $v(\{1,3\}) = 6$, $v(\{2,3\}) = 5$, $v(N) = 8$ il risolutore conferma che
non esiste.
\end{osservazione}

\subsection{Valore di Shapley}

\begin{lstlisting}
def shapley_value(player, v):
    N = max(v)
    n = len(N)
    phi = 0.0
    for C in v:
        if player in C:
            s = len(C)
            weight = factorial(n - s) * factorial(s - 1) / factorial(n)
            phi += weight * (v[C] - v.get(C - {player}, 0))
    return phi

def shapley(v):
    return {p: shapley_value(p, v) for p in max(v)}
\end{lstlisting}

Risultato sul gelato: $\{A: 583{,}33;\ B: 208{,}33;\ C: 208{,}33\}$, come atteso. Una seconda
implementazione con le permutazioni dà gli stessi valori.

\paragraph{Sottoinsiemi o permutazioni.}
Il notebook raccomanda la formula sui sottoinsiemi: $O(2^n)$ invece di $O(n!)$. Su giochi
casuali con $n$ giocatori, tempi misurati per calcolare il valore di \emph{tutti} i giocatori:
\begin{center}
\renewcommand{\arraystretch}{1.2}
\begin{tabular}{rrr}
\toprule
$n$ & Sottoinsiemi & Permutazioni \\
\midrule
6  & 0,1 ms & 3,8 ms \\
8  & 0,8 ms & 304 ms \\
9  & 1,9 ms & 3\,409 ms \\
10 & 5,5 ms & $3\,628\,800$ ordini, non eseguito \\
\bottomrule
\end{tabular}
\end{center}
I valori calcolati con i due metodi coincidono ogni volta. Il codice sopra scorre tutte le
coalizioni per ogni giocatore, quindi il costo totale è $n\cdot 2^n$: comunque esponenziale.
Per giochi con molti giocatori si usano approssimazioni (per esempio campionando permutazioni
a caso) o rappresentazioni compatte, come anticipa la fine del notebook.

\begin{attenzione}[``il valore di Shapley è individualmente razionale'']
Il notebook elenca tra le proprietà $\phi(i,v) \ge v(\{i\})$. Vale per i giochi
\emph{superadditivi}: lì ogni contributo marginale è
$v(C \cup \{i\}) - v(C) \ge v(\{i\})$, e quindi lo è anche la media. Senza superadditività
può fallire. Il notebook stesso ricorda subito dopo l'altro limite: il valore di Shapley
non è necessariamente nel core.
\end{attenzione}

\subsection{In più: struttura di coalizione e nucleolo}

\paragraph{Programmazione dinamica.}
Implementazione diretta dell'equazione della teoria, coalizione per coalizione in ordine di
dimensione:
\begin{lstlisting}
def coalition_structure_dp(v):
    N = max(v)
    best, struct = {}, {}
    for C in sorted(powerset(sorted(N)), key=len):
        best[C], struct[C] = v[C], [C]
        members = sorted(C)
        for k in range(1, len(C)):
            for part in combinations(members, k):
                C1 = frozenset(part); C2 = C - C1
                if min(C1) != min(C):          # ogni divisione una volta sola
                    continue
                if best[C1] + best[C2] > best[C]:
                    best[C] = best[C1] + best[C2]
                    struct[C] = struct[C1] + struct[C2]
    return best, struct
\end{lstlisting}
Sul gioco dell'esempio restituisce esattamente la tabella corretta: $v^*(\{B,C\}) = 400$ e
$v^*(N) = 1050$ con $CS^* = \{\{A,C\},\{B\}\}$.

\paragraph{Nucleolo.}
Implementato con l'algoritmo di Kopelowitz e \texttt{linprog}, fissando a ogni passo solo le
coalizioni il cui eccesso non si può abbassare. Risultati:
\begin{center}
\renewcommand{\arraystretch}{1.2}
\begin{tabular}{lll}
\toprule
Gioco & Nucleolo & Vettore degli eccessi \\
\midrule
Simmetrico ($v(ij) = 1$, $v(N) = 3$) & $(1, 1, 1)$ & $(-1,-1,-1,-1,-1,-1)$ \\
Conta solo $A$ ($v(N) = 750$) & $(750, 0, 0)$ & $(0,0,0,0,0,-750)$ \\
Gelato ($v(N) = 1000$) & $(750, 125, 125)$ & $(-125,-125,-125,-125,-250,-750)$ \\
\bottomrule
\end{tabular}
\end{center}
Nel gelato nucleolo e Shapley sono diversi: $(750; 125; 125)$ contro
$(583{,}3; 208{,}3; 208{,}3)$. Entrambi stanno nel core, ma il nucleolo ``bilancia'' le
insoddisfazioni (le coppie $\{A,B\}$ e $\{A,C\}$ e i singoli $B$ e $C$ hanno tutti eccesso
$-125$), mentre Shapley premia i contributi marginali.

% ---------------------------------------------------------------------
\begin{riepilogo}
\begin{itemize}
  \item Gioco in funzione caratteristica $(N,v)$: $v(C)$ è quanto guadagna la coalizione
        $C$; utilità trasferibile.
  \item Struttura di coalizione: partizione di $N$ che massimizza
        $\sum v(C)$. Le partizioni sono il numero di Bell; la programmazione dinamica la
        trova in $O(3^n)$.
  \item Superadditività: fondersi non fa mai perdere, quindi si forma la grande coalizione e
        un esito è un vettore $x$ con $\sum x_i = v(N)$.
  \item \textbf{Core}: esiti in cui ogni coalizione riceve almeno $v(C)$. Può essere vuoto e
        può contenere esiti iniqui.
  \item \textbf{Nucleolo}: minimizza lessicograficamente gli eccessi
        $e(S,x) = v(S) - x(S)$; esiste, è unico, sta nel core se questo non è vuoto; si
        calcola con una sequenza di programmi lineari (Kopelowitz). Può essere iniquo.
  \item \textbf{Shapley}: contributo marginale medio su tutti gli ordini di ingresso; unico
        a soddisfare efficienza, giocatore nullo, simmetria, additività. Equo ma non
        necessariamente stabile.
  \item \textbf{Lab}: \texttt{frozenset} come chiavi, superadditività, verifica del core,
        Shapley sui sottoinsiemi ($n\cdot2^n$ contro $n!$); in più programmazione dinamica e
        nucleolo.
\end{itemize}
\end{riepilogo}


\end{document}
